% concatenated from main.tex
\documentclass[12pt,fleqn,a4paper]{amsart}

% --- Encoding and layout ---
\usepackage[utf8]{inputenc}
\usepackage[T1]{fontenc}
\usepackage{a4wide}

% --- Maths packages ---
\usepackage{amsmath,amssymb,amsthm,latexsym,hyperref,enumitem}
\usepackage{mathrsfs}
\usepackage{url}

\renewcommand{\le}{\leqslant}
\renewcommand{\ge}{\geqslant}

\setlist[enumerate,1]{label=\textup{(\roman*)},leftmargin=*}
\setlist[itemize,1]{leftmargin=*}

\numberwithin{equation}{section}
\raggedbottom

% --- Theorems ---
\newtheorem{theorem}{Theorem}[section]
\newtheorem{mainthm}[theorem]{Main Theorem}
\newtheorem*{theoremA}{Theorem A}
\newtheorem*{theoremB}{Theorem B}
\newtheorem*{theoremC}{Theorem C}
\newtheorem*{theoremD}{Theorem D}
\newtheorem*{theoremE}{Theorem E}
\newtheorem{lemma}[theorem]{Lemma}
\newtheorem{corollary}[theorem]{Corollary}
\newtheorem{proposition}[theorem]{Proposition}
\theoremstyle{definition}
\newtheorem{definition}[theorem]{Definition}
\newtheorem{remark}[theorem]{Remark}
\newtheorem{example}[theorem]{Example}
\newtheorem{question}[theorem]{Question}
\newtheorem{convention}[theorem]{Convention}

% --- Shortcuts ---
\newcommand{\N}{\mathbb{N}}
\newcommand{\Q}{\mathbb{Q}}
\newcommand{\C}{\mathbb{C}}
\newcommand{\Z}{\mathbb{Z}}
\newcommand{\R}{\mathbb{R}}
\newcommand{\RR}{\mathbb{R}}

\newcommand{\SB}{\mathrm{SB}}
\newcommand{\Ideal}{\mathrm{Ideal}}
\newcommand{\dist}{\mathrm{dist}}
\newcommand{\id}{\mathrm{id}}
\newcommand{\Min}{\mathrm{Min}}
\newcommand{\Con}{\mathrm{Con}}
\newcommand{\NSub}{\mathrm{NSub}}
\newcommand{\Aut}{\mathrm{Aut}}
\newcommand{\Inn}{\mathrm{Inn}}
\newcommand{\Id}{\mathrm{Id}}
\newcommand{\Sub}{\mathrm{Sub}}
\newcommand{\WF}{\mathrm{WF}}
\newcommand{\CTbl}{\mathrm{CTbl}}

\title[Polish spaces for countable and separable structures]
{Polish spaces for countable and separable structures\\ through quotient encodings}

\author[T.~Kania]{Tomasz Kania}
\address[T.~Kania]{Mathematical Institute\\Czech Academy of Sciences\\\v Zitn\'a 25 \\115 67 Praha 1\\Czech Republic  and  Institute of Mathematics and Computer Science\\ Jagiellonian University\\ {\L}ojasiewicza 6, 30-348 Krak\'{o}w, Poland
}
\email{kania@math.cas.cz, tomasz.marcin.kania@gmail.com}
\thanks{RVO: 67985840.}
\date{\today}

% --- MSC and Keywords ---
\subjclass[2020]{Primary 03E15, 46L05; Secondary 03C15, 46H05, 46L80, 54B20}
\keywords{Descriptive set theory, Borel complexity, Polish space, Wijsman topology, $C^*$-algebra, Banach algebra, hyperspace, countable structure, quotient encoding, operator algebra $K$-theory, $\Pi^1_1$-complete}

\begin{document}

\begin{abstract}
We develop a quotient-based framework for locating natural properties of countable
algebraic structures and separable Banach-type structures in the Borel hierarchy.
The common idea is to present an object as a quotient of a fixed generator and to
read definability from the corresponding kernel or congruence.

For separable Banach-type structures, including Banach algebras, $C^*$-algebras and
TROs, admissible kernels form Polish spaces; in the Wijsman topology the quotient-norm
functional $K\mapsto \|x+K\|$ is continuous.  This gives a uniform definability
scheme with explicit Borel upper bounds.  For countable algebraic structures,
congruence spaces are compact zero-dimensional Polish spaces and atomic predicates are
clopen.

For Banach algebras we obtain, among other estimates, closedness of commutativity,
the abstract uniform-algebra norm identity and Dedekind finiteness, and $G_\delta$
bounds for topological stable-rank bounds.  In the unital $C^*$-algebra coding based
on $C^*_{\max}(F_\infty)$ we obtain closedness of stable finiteness and existence of
a tracial state, $G_\delta$ bounds for AF-ness, MF-ness, approximate divisibility and
real-rank bounds; a $\Pi^0_3$ bound for nuclear dimension; Borelness of nuclearity
and simplicity, and analyticity of $D$-absorption for fixed exact $D$.  The
$G_\delta$ bounds for AF-ness, real-rank bounds and topological stable-rank bounds
are shown to be sharp by continuous reductions from a canonical $\Pi^0_2$-complete
set.  We give an internal Borel coding of the
$K_0$-assignment: every coordinate section is $F_\sigma$, and the resulting map into
the standard subgroup coding of countable abelian groups is of Baire class~$2$.
Suspension and Bott periodicity, combined with the known standard coding computations,
yield Borel codings of $K_1$ and all higher $K$-groups.

The framework also gives a compact quotient treatment of countable groups, rings,
lattices, Boolean algebras and abelian groups.  We include continuous completeness
reductions for several low-rank algebraic properties and exhibit natural
$\Pi^1_1$-complete properties: separability of the dual in the commutative
$C^*$-quotient coding and superatomicity in the Boolean-algebra quotient coding.
\end{abstract}

\maketitle

\section{Introduction}

Descriptive set theory gives a language for measuring the complexity of mathematical
classification problems, but classification complexity and property complexity are not
the same question.  The isomorphism relation for separable $C^*$-algebras is already
maximally complicated among orbit equivalence relations; Sabok proved this even for
separable simple AI $C^*$-algebras~\cite{SabokCompleteness}.  Nevertheless, once a
single algebra is presented by a code, one can still ask where concrete properties such
as stable finiteness, AF-ness, nuclear dimension bounds, traciality or $K$-theory lie
inside the Borel hierarchy.  The purpose of this paper is to give a quotient-based
framework in which many such computations become uniform.

The guiding principle is simple.  Instead of embedding all objects into a universal
ambient object, we present them as quotients of a fixed generator.  This is forced in
some Banach-algebraic categories: no separable Banach algebra contains an isomorphic
copy of every separable Banach algebra~\cite{KaniaNoUniversalBanachAlg}.  By contrast,
quotient generators are abundant: for example $C^*_{\max}(F_\infty)$ generates the
separable unital $C^*$-algebras by quotients, and free objects play the same role for
countable algebraic structures.  Thus the natural parameter is not an embedded copy of
the object, but the kernel, ideal or congruence defining the quotient.

This viewpoint also fits the existing descriptive-set-theoretic codings.  Kechris
introduced a standard Borel parameterisation of separable $C^*$-algebras
\cite{KechrisAsianCstar}, and Farah--Toms--T\"ornquist developed several mutually
Borel-equivalent codings together with Borel computations of invariants
\cite{FarahTomsTornquist,FarahTomsTornquistCrelle}.  On the Banach-space side,
C\'uth, Dole\v{z}al, Doucha and Kurka studied Polish spaces of separable Banach
spaces via norm and pseudonorm codings~\cite{CuthDolezalDouchaKurka2022,CuthDolezalDouchaKurka2024}.
The present paper is complementary: it keeps the quotient presentation visible and uses
the topology of the kernel space to read off Borel ranks.

The main results are organised around the following general theorems.  They are stated
informally here; the precise versions, including the hypotheses on the chosen quotient
generators and topologies, appear in the body of the paper.  The concrete references
are collected after the statements so that the statements themselves remain readable.

\begin{theoremA}[Continuous quotient spaces and definability]
Let $P$ be a separable quotient generator in a Banach-type category with countably many
continuous operations.  The admissible kernels in $P$ form a Polish space for every
admissible hyperspace topology.  In the Wijsman topology the quotient-norm functional
\[
        K\longmapsto \|x+K\|_{P/K}
\]
is continuous for each $x\in P$.  Consequently, every property expressible by the
quotient-norm language and stable existential relations is Borel, with Wijsman-rank
bounded by the visible quantifier alternation.
\end{theoremA}

\begin{theoremB}[Countable quotient spaces]
For countable algebraic structures, congruences on a free countable object form compact
zero-dimensional Polish spaces.  Atomic quotient predicates are clopen, the canonical
coding of quotients on domain $\mathbb N$ is Borel, and first-order algebraic
complexity is reflected directly by the Borel hierarchy.
\end{theoremB}

\begin{theoremC}[$C^*$-algebraic regularity bounds]
In the unital quotient coding based on $A=C^*_{\max}(F_\infty)$, stable finiteness
and the existence of a tracial state are closed.  AF-ness, MF-ness, approximate
divisibility, real-rank bounds and topological stable-rank bounds are $G_\delta$.
The $G_\delta$ upper bounds for AF-ness, real-rank bounds and topological stable-rank
bounds are optimal: these classes are $\Pi^0_2$-complete.  Quasidiagonality and
nuclear dimension $\le n$ are $\Pi^0_3$; for nuclear dimension and property~(SP) we
also record commutative $\Pi^0_2$ lower bounds, but no $\Pi^0_3$-hardness is claimed.
Nuclearity and simplicity are Borel, and $D$-absorption is analytic for each fixed
separable unital exact $D$.
\end{theoremC}

\begin{theoremD}[Borel $K$-theory and tensor kernels]
There is an internal Borel assignment
\[
        \kappa_{K_0}:\Ideal(A)\longrightarrow \Sub(\mathbb Z^{(\mathbb N)})
\]
such that $\mathbb Z^{(\mathbb N)}/\kappa_{K_0}(K)\cong K_0(A/K)$; each fixed
coordinate section is $F_\sigma$, and the map is of Baire class~$2$.  Suspension,
Bott periodicity and the standard coding computations yield Borel assignments for
$K_1$ and all higher $K$-groups.  Tensor-product kernel maps are Borel in general and
continuous in the exact cases considered below.
\end{theoremD}

\begin{theoremE}[Coanalytic calibration]
Some natural quotient-coded properties are not Borel.  In the commutative
$C^*$-quotient coding, having separable dual is $\Pi^1_1$-complete.  In the
Boolean-algebra quotient coding, superatomicity is $\Pi^1_1$-complete.
\end{theoremE}

The precise versions behind these summaries are as follows.  Theorem~A is proved as
Theorems~\ref{thm:continuous-main} and~\ref{thm:definability-rank}; Theorem~B is
Theorem~\ref{thm:countable-main}.  The upper-bound part of Theorem~C is proved across
Section~\ref{sec:Cstar-properties}; its sharpness part is Theorem~\ref{thm:Cstar-Pi2-complete},
and the $D$-absorption assertion is Theorem~\ref{thm:Dstab}.  The $K$-theoretic part of
Theorem~D is Theorems~\ref{thm:K0} and~\ref{thm:K1}, together with
Corollary~\ref{cor:K0-map-complexity}; the tensor-kernel part is Section~\ref{sec:tensor-products},
especially Theorem~\ref{thm:tensor-continuity}.  Theorem~E is
Theorems~\ref{thm:dual-separable-Pi11} and~\ref{thm:superatomic-Pi11}.

The catalogue results below are deliberately conservative.  When a low finite rank is
claimed, the proof supplies an explicit countable scheme in the Wijsman topology.  When
only Borelness is asserted, no hidden rank optimisation is intended.  In particular, the
paper does not claim a $G_\delta$ or $\Pi^0_2$ bound for simplicity of separable
$C^*$-algebras or Banach algebras; the Borel statement for $C^*$-algebras is obtained by
comparison with the standard codings and known definability results.

\subsection{Organisation}

Section~\ref{sec:prelim} fixes notation, recalls admissible hyperspace topologies,
Wijsman topology and Lusin--Novikov uniformisation.  Section~\ref{sec:continuous}
proves Theorem~A.  Section~\ref{sec:countable} proves Theorem~B.  Section~\ref{sec:examples}
records quotient generators.  Sections~\ref{sec:banach-complexity} and
\ref{sec:Cstar-properties} establish the Banach-algebraic and $C^*$-algebraic
complexity estimates, with the unital $C^*$-algebra scope made explicit in
Remark~\ref{rem:Cstar-unital-scope}.  Section~\ref{sec:K0} proves the
$K$-theory assignment theorem, and Section~\ref{sec:tensor-products} treats
tensor-product kernels.  Section~\ref{sec:D-stability} proves analyticity of
$D$-absorption.  Section~\ref{sec:TRO} records the free TRO construction and the
TRO estimates that follow from stable finite-dimensional relations.  Section~\ref{sec:countable-complexity}
gives the countable-structure catalogue and the low-rank completeness reductions.
Section~\ref{sec:Pi11-example} proves the $\Pi^1_1$-complete calibration results.
\section{Preliminaries and notation}
\label{sec:prelim}
%%%%%%%%%%%%%%%%%%%%%%%%%%%%%%%%%%%%%%%%%%%%%%%%%%%%%%%%%%%%%%

Throughout, $X$ denotes a separable Banach space.  For a non-empty closed set
$F\subseteq X$ and $x\in X$ we write
\[
d(x,F)=\inf\{\|x-y\|:y\in F\}.
\]
If $F$ is a closed linear subspace then $d(x,F)=\|x+F\|_{X/F}$.
For each separable space $X$ under discussion we fix a countable dense set
$D^X=\{u_m:m\in\N\}$.

We shall use standard descriptive set theory; see~\cite{Kechris,BeckerKechris,GaoDST}.
We recall the key notions for the reader's convenience.

\subsection{The Borel and projective hierarchies}\label{subsec:borel-hierarchy}

Let $X$ be a Polish space (separable, completely metrisable).
The \emph{Borel hierarchy} stratifies the Borel $\sigma$-algebra of $X$ into
countable levels.  At the base,
$\Sigma^0_1$ denotes the open sets and $\Pi^0_1$ the closed sets.
Inductively, for $n\ge 1$:
$\Sigma^0_{n+1}$ consists of countable unions of $\Pi^0_n$ sets, and
$\Pi^0_{n+1}$ consists of countable intersections of $\Sigma^0_n$ sets.
Equivalently, $\Pi^0_n$ is the class of complements of $\Sigma^0_n$ sets.
In low levels:
$\Sigma^0_2 = F_\sigma$ (countable unions of closed sets),
$\Pi^0_2 = G_\delta$ (countable intersections of open sets),
$\Sigma^0_3 = G_{\delta\sigma}$, and $\Pi^0_3 = F_{\sigma\delta}$.

Beyond the Borel hierarchy lie the \emph{projective} (or \emph{analytic}) classes.
A set is $\Sigma^1_1$ (\emph{analytic}) if it is the continuous image of a Borel set
in another Polish space, or equivalently the projection of a closed subset of
$X\times\N^\N$.  A set is $\Pi^1_1$ (\emph{coanalytic}) if its complement is analytic.
Suslin's theorem asserts that a set is Borel iff it is both analytic and coanalytic.

\smallskip
\noindent\textbf{Hardness and completeness.}
Let $\Gamma$ be a pointclass ($\Sigma^0_n$, $\Pi^0_n$, $\Sigma^1_1$, $\Pi^1_1$, etc.).
A set $A\subseteq X$ is \emph{$\Gamma$-hard} if for every set $B\in\Gamma(Y)$ (where $Y$
is any Polish space) there exists a continuous map $f:Y\to X$ with $f^{-1}(A)=B$;
informally, $A$ is ``at least as complex as any $\Gamma$ set''.
The set $A$ is \emph{$\Gamma$-complete} if it belongs to $\Gamma$ and is $\Gamma$-hard.
Completeness for $\Gamma$ implies that $A$ does not belong to the dual class
$\check\Gamma$ (e.g.\ a $\Sigma^0_2$-complete set is not $\Pi^0_2$, and a
$\Pi^1_1$-complete set is not Borel).
Proving a property $\Gamma$-complete therefore establishes that its upper bound
within the hierarchy is \emph{optimal} and cannot be improved.

\smallskip
\noindent\textbf{Connection to definability.}
In this paper, upper bounds on Borel complexity arise from
expressing properties via countable quantifier alternations over
continuous predicates.  Each universal quantifier (over a countable domain) preserves
$\Pi^0_n$; each existential quantifier preserves $\Sigma^0_n$; alternation
increments the level by one.  Lower bounds (completeness results) are proved
by explicit continuous reductions from canonical complete sets, such as the set of
eventually-zero sequences in $2^\N$ ($\Sigma^0_2$-complete) or the set of
well-founded trees ($\Pi^1_1$-complete).

\subsection{Hyperspaces and admissible topologies}\label{sec:hyperspace}

Let $X$ be a separable metric space and let $\mathscr F(X)$ denote the hyperspace of
non-empty closed subsets of $X$.

\begin{definition}\label{def:wijsman}
The \emph{Wijsman topology} on $\mathscr F(X)$ is the coarsest topology making each map
$F\mapsto d(x,F)$ continuous for every fixed $x\in X$.
Equivalently, a net $(F_\alpha)$ converges to $F$ in the Wijsman topology iff
$d(x,F_\alpha)\to d(x,F)$ for every $x\in X$; in other words, the distance functions
converge pointwise.

If $X$ is separable, the Wijsman topology is metrisable
(indeed, fixing a countable dense set $\{u_m\}$ one can take
$d_W(F,G)=\sum_m 2^{-m}\min\bigl(1,|d(u_m,F)-d(u_m,G)|\bigr)$
as a compatible metric).
If $X$ is Polish, then $(\mathscr F(X),\mathrm{Wijsman})$ is Polish;
see Beer~\cite{Beer}.
\end{definition}

The Wijsman topology is the natural topology on our parameter spaces throughout
this paper: for Banach-type structures, it makes the quotient-norm functional
$K\mapsto\|x+K\|$ continuous by definition, and for countable structures
the Wijsman topology on closed subsets of a discrete space coincides with the
Vietoris topology.

\begin{lemma}\label{lem:wijsman-approx}
Let $F_n\to F$ in the Wijsman topology.
\begin{enumerate}
\item If $x\in F$, there exist $x_n\in F_n$ with $x_n\to x$.
\item If $F_n\subseteq G_n$ for all $n$ and $G_n\to G$ (Wijsman), then $F\subseteq G$.
\end{enumerate}
\end{lemma}

\begin{proof}
(i) Since $d(x,F)=0$ and $d(x,F_n)\to d(x,F)=0$, choose $x_n\in F_n$ with
$\|x-x_n\|\le d(x,F_n)+1/n$.  Then $\|x-x_n\|\to 0$.

(ii) Let $x\in F$.  Wijsman convergence $F_n\to F$ gives $d(x,F_n)\to d(x,F)=0$.
Since $F_n\subseteq G_n$ for all $n$, we have $d(x,G_n)\le d(x,F_n)$ for all $n$,
hence $d(x,G_n)\to 0$.  As $G_n\to G$ in the Wijsman topology,
$d(x,G_n)\to d(x,G)$, so $d(x,G)=0$ and therefore $x\in G$.
\end{proof}

The Wijsman topology is one instance of a broader class of
hyperspace topologies introduced by Godefroy and
Saint-Raymond~\cite{GodefroySaintRaymond2018}.

\begin{definition}
\label{def:admissible}
A Polish topology $\tau$ on $\mathscr F(X)$ is \emph{admissible} if:
\begin{enumerate}
\item for each open $U\subseteq X$,
$E^+(U)=\{F\in\mathscr F(X):F\cap U\ne\varnothing\}$ is $\tau$-open;
\item $\tau$ has a subbase consisting of countable unions of sets of the form
$E^+(U)\setminus E^+(V)$ with $U,V$ open in $X$;
\item the membership relation $\{(x,F):x\in F\}\subseteq X\times\mathscr F(X)$
is $\tau$-closed.
\end{enumerate}
\end{definition}

Condition~(i) says that the topology detects hits:
if a closed set meets an open region, nearby sets (in the $\tau$-sense) also
do.  Condition~(ii) is a regularity constraint ensuring $\tau$
is generated by set-theoretically definable data.
Condition~(iii) is the key geometric requirement: it says precisely that the
singleton embedding $x\mapsto \{x\}$ is a homeomorphism from $X$ onto a
closed subset of $(\mathscr F(X),\tau)$ --- in other words, points are
``robustly detected'' by the topology, and any sequence of closed sets converging
in~$\tau$ to a point mass $\{x\}$ must eventually carry elements near~$x$.

The Wijsman topology is admissible.
The first two conditions are standard (see~\cite{Beer}).  For the third:
if $(x_n,F_n)\to(x,F)$ in $X\times\mathscr F(X)$ and $x_n\in F_n$ for all~$n$,
then $d(x,F_n)\le d(x,x_n)+d(x_n,F_n)=d(x,x_n)\to 0$,
so $d(x,F)=\lim_n d(x,F_n)=0$, hence $x\in F$ and the membership relation is closed.

Let $\SB(X)$ denote the family of closed linear subspaces of~$X$.  For admissible~$\tau$,
$\SB(X)$ is a $G_\delta$ subset of $(\mathscr F(X),\tau)$ and hence Polish
\cite{GodefroySaintRaymond2018}.

\begin{proposition}
\label{prop:SBX-Wijsman-closed}
Let $F_n\in\SB(X)$ and $F_n\to F$ in the Wijsman topology on $\mathscr F(X)$.
Then $F\in\SB(X)$; hence $\SB(X)$ is closed in $(\mathscr F(X),\mathrm{Wijsman})$.
\end{proposition}

\begin{proof}
By Wijsman convergence, $d(0,F_n)=0$ for all $n$ implies $d(0,F)=\lim_n d(0,F_n)=0$,
so $0\in F$.

Let $x,y\in F$ and $a,b\in\Q$.  By the Wijsman approximation lemma, there exist
$x_n,y_n\in F_n$ with $x_n\to x$ and $y_n\to y$ in $X$.
Since each $F_n$ is a linear subspace, we have $z_n:=ax_n+by_n\in F_n$.
By continuity of addition and scalar multiplication, $z_n\to z:=ax+by$.
Therefore $d(z,F_n)\le \|z-z_n\|\to 0$,
and Wijsman convergence yields
$d(z,F)=\lim_n d(z,F_n)=0$, so $z\in F$.
Hence $F$ is closed under rational linear combinations.
Finally, for arbitrary $t\in\R$ and $x\in F$, choose rationals $t_k\to t$;
then $t_k x\in F$ for all~$k$, and since $F$ is closed, $tx=\lim_k t_k x\in F$.
Thus $F$ is a closed linear subspace.
\end{proof}

\paragraph{Substructure variant.}
The same argument does not use linearity per se.  If $X$ is a separable metric
space equipped with countably many continuous operations
$\omega_i:X^{n_i}\to X$, then the family of all non-empty closed subsets
$F\subseteq X$ that are closed under each $\omega_i$ is Wijsman-closed in
$\mathscr F(X)$.  Indeed, if $F_n\to F$ in the Wijsman topology and each $F_n$
is a closed substructure, then for any $x_1,\dots,x_{n_i}\in F$ one may choose
$x_j^{(n)}\in F_n$ with $x_j^{(n)}\to x_j$ by
Lemma~\ref{lem:wijsman-approx}.  Continuity gives
$\omega_i(x_1^{(n)},\dots,x_{n_i}^{(n)})\to \omega_i(x_1,\dots,x_{n_i})$, and since
$\omega_i(x_1^{(n)},\dots,x_{n_i}^{(n)})\in F_n$ for all $n$, Wijsman convergence
implies $\omega_i(x_1,\dots,x_{n_i})\in F$.  (For nullary operations, the same
argument applies to the distinguished constants.)  In particular, for a separable
$C^*$-algebra $A$, the family of closed $*$-subalgebras of $A$ is Wijsman-closed.

\begin{remark}\label{rem:SBX-Gdelta}
One should \emph{not} expect $\SB(X)$ to be $\tau$-closed for every admissible~$\tau$
in full generality.  The defining condition
``$\forall x,y\in F$ and $\forall a,b\in\Q$, we have $ax+by\in F$'' is universal
and thus naturally produces countable intersections of open conditions, but universal
formulae do not automatically yield closed sets for arbitrary admissible hyperspace
topologies.  For Wijsman, however, the above proposition gives closedness.
\end{remark}

The key tool for working with admissible topologies is the following selection theorem,
which provides continuous dense selections and is the workhorse that allows us to
pass from the abstract hyperspace to concrete dense-set data.

\begin{theorem}\label{thm:GSR}
Let $X$ be a separable Banach space and $\tau$ an admissible topology on $\SB(X)$.
There exist continuous maps $f_k:(\SB(X),\tau)\to X$ ($k\in\N$) such that
\[
f_k(F)\in F\ \ \text{for all }F\in\SB(X)\text{ and }k\in\N,
\quad
\overline{\{f_k(F):k\in\N\}}=F\quad (F\in\SB(X)).
\]
\end{theorem}

\begin{proof}
See \cite{GodefroySaintRaymond2018}.
\end{proof}

\subsection{Matrix-level quotient norms}\label{subsec:matrix-Phi}

In the $C^*$-algebra parts of this paper, several arguments require norms
in \emph{matrix amplifications} $M_n(A/K)$ rather than merely in~$A/K$ itself
(see \cite[II.6.6]{BlackadarOA} or \cite[Chapter~1]{BrownOzawa} for background on
matrix norms and operator-space structure).
We introduce a systematic notation.

\begin{definition}\label{def:Phi-matrix}
Let $A$ be a separable $C^*$-algebra and $K\in\Ideal(A)$.
Because $M_n$ is nuclear, the canonical identification $M_n(A)\cong A\otimes M_n$
carries $M_n(K)=K\otimes M_n$ isometrically.  For $X\in M_n(A)$ define
\[
\Phi^{(n)}(K,X):=\dist\bigl(X,\,M_n(K)\bigr)=\|X+M_n(K)\|_{M_n(A/K)}.
\]
\end{definition}

\begin{lemma}\label{lem:Phi-n-continuous}
For each fixed $n\in\N$ and $X\in M_n(A)$, the map
$K\mapsto \Phi^{(n)}(K,X)$ is continuous in the Wijsman topology on $\Ideal(A)$.
\end{lemma}

\begin{proof}
It is enough to check sequential continuity, since the Wijsman topology on
$\Ideal(A)$ is metrisable.  Suppose $K_m\to K$ Wijsman and put
$T=\N\cup\{\infty\}$, with $K_\infty=K$.  For $a\in A$ the function
\[
        t\longmapsto \|a+K_t\|=\dist(a,K_t)
\]
is continuous on~$T$.  By Dixmier's criterion for quotient continuous fields
\cite[Prop.~10.3.2]{Dixmier1977}, the quotients $A/K_t$ form a continuous field over
$T$ with the canonical sections $t\mapsto a+K_t$.

Finite matrix amplification preserves continuous fields: applying the same criterion
to the matrix sections
\[
        t\longmapsto [x_{ij}+K_t]\in M_n(A/K_t),
        \qquad X=[x_{ij}]\in M_n(A),
\]
gives continuity of
\[
        t\longmapsto \|[x_{ij}+K_t]\|_{M_n(A/K_t)}.
\]
Under the canonical quotient identification
$M_n(A)/M_n(K_t)\cong M_n(A/K_t)$ this norm is precisely
$\dist(X,M_n(K_t))$.  Hence
$\Phi^{(n)}(K_m,X)\to\Phi^{(n)}(K,X)$.
\end{proof}

\begin{remark}\label{rem:Phi-n-dense}
Fix a countable dense $*$-subalgebra $D^A\subseteq A$.  Then the set 
\[
    M_n(D^A)=\{(a_{ij})_{i,j\le n}:a_{ij}\in D^A\}
\]
is a countable dense
subset of~$M_n(A)$, and it is closed under multiplication, $*$, and rational
linear combinations.
All arguments below that invoke $\Phi^{(n)}$ may therefore be evaluated on
this countable dense set with no loss.
\end{remark}

\subsection{Projection of locally closed sets along compact spaces}

The following elementary observation is used several times (most explicitly in Proposition~\ref{prop:nucdim}).

\begin{lemma}\label{lem:proj-locally-closed}
Let $X$ be a compact metric space, let $Y$ be a Polish space, and let $C\subseteq X\times Y$ be closed and $U\subseteq X\times Y$ open.  Then $\pi_Y(C\cap U)$ is $F_\sigma$ in~$Y$.
\end{lemma}

\begin{proof}
Write
\[
U=\bigcup_{m\ge 1}F_m,
\qquad
F_m:=\bigl\{z\in X\times Y:d\bigl(z,(X\times Y)\setminus U\bigr)\ge 1/m\bigr\}.
\]
Each $F_m$ is closed in $X\times Y$ and satisfies $F_m\subseteq U$.
(We do \emph{not} claim that $F_m$ is compact, since only the $X$-factor is compact.)
Now $C\cap F_m$ is closed in $X\times Y$, and because $X$ is compact and $Y$ is
Hausdorff, the projection $\pi_Y(C\cap F_m)$ is closed in~$Y$.
Therefore
\[
\pi_Y(C\cap U)=\bigcup_{m\ge 1}\pi_Y(C\cap F_m)
\]
is $F_\sigma$.
\end{proof}


\subsection{Lusin--Novikov uniformisation}

We shall invoke Lusin--Novikov uniformisation for countable Borel fibres; we record it
for reference.

\begin{theorem}\label{thm:LN}
Let $X$ and $Y$ be standard Borel spaces and let $B\subseteq X\times Y$ be Borel.
If each section $B_x=\{y:(x,y)\in B\}$ is countable, then there exist Borel maps
$f_n:X\to Y$ ($n\in\N$) such that
\[
B_x=\{f_n(x):n\in\N\}
\quad\text{for all }x\in X
\]
(allowing repetitions).
\end{theorem}

\begin{proof}
See~\cite[Theorem~18.10]{Kechris}.
\end{proof}

\begin{convention}\label{conv:enumeration}
Every countable algebraic object (free group $F_\infty$, free ring $R_\infty$, free
algebra $F_L$, etc.)\ is identified with $\N$ via a fixed effective enumeration chosen
once and for all.
This makes expressions such as $\min\{[a]_S:a\in F_\infty\}$ well-defined throughout.
\end{convention}

%%%%%%%%%%%%%%%%%%%%%%%%%%%%%%%%%%%%%%%%%%%%%%%%%%%%%%%%%%%%%%
\section{Main theorem for separable continuous structures}\label{sec:continuous}
%%%%%%%%%%%%%%%%%%%%%%%%%%%%%%%%%%%%%%%%%%%%%%%%%%%%%%%%%%%%%%

This section develops the analytic half of our general framework:
we construct Polish parameter spaces for separable Banach-type structures
(such as Banach algebras, $C^*$-algebras, Banach lattices, and ternary rings of
operators) and establish a general definability theorem linking algebraic formulae
to Borel complexity.  The key idea is to view such structures not abstractly, but
concretely as \emph{quotients} of a single universal object endowed with
finitely or countably many continuous multilinear operations.

\subsection{Categories with continuous operations}

\begin{definition}\label{def:category-ops}
A \emph{concrete category $\mathcal{C}$ with continuous operations} consists of:
\begin{itemize}
\item \textbf{Objects}: Separable Banach spaces $X$ equipped with countably many
continuous operations
$\omega_i: X^{n_i} \to X$, $i\in I$ (countable), where $n_i\in\N\cup\{0\}$ is the arity
of~$\omega_i$.
\item \textbf{Morphisms}: Bounded linear maps $T: X \to Y$ preserving all operations:
\[
T(\omega_i^X(x_1, \ldots, x_{n_i}))
= \omega_i^Y(T(x_1), \ldots, T(x_{n_i})).
\]
\item \textbf{Quotients}: If $T: X \to Y$ is surjective with kernel $K = \ker(T)$,
then $Y \cong X/K$ as objects in $\mathcal{C}$, where operations descend to quotients.
\end{itemize}
\end{definition}

\begin{definition}\label{def:admissible-kernel}
Let $X\in\mathcal{C}$.  A closed linear subspace $K\subseteq X$ is an
\emph{admissible kernel} (or \emph{$\mathcal{C}$-ideal}) if for every $i\in I$ and all
$x_1,\ldots,x_{n_i},x_1',\ldots,x_{n_i}'\in X$ with $x_j-x_j'\in K$ for
$1\le j\le n_i$, one has
\[
\omega_i(x_1,\dots,x_{n_i})-\omega_i(x_1',\dots,x_{n_i}')\in K.
\]
Equivalently, $\omega_i$ descends to a well-defined operation on $X/K$;
in the Banach setting, this means that the equivalence relation
$E_K$ on $X$ given by $x\,E_K\,y$ iff $x-y\in K$ is a congruence for all the
operations $\omega_i$.  We write $\Ideal_{\mathcal C}(X)$ for the collection of
admissible kernels.
\end{definition}

\paragraph{Congruence viewpoint.}
The Banach-space formulation is chosen because later arguments exploit the quotient
norm.  Abstractly, the same quotient-coding setup can be stated for separable
Polish structures with continuous operations by replacing closed linear subspaces
$K\subseteq X$ with closed congruences $E\subseteq X^2$.  In the Banach setting,
$K$ and $E_K$ are interchangeable, and Section~\ref{sec:countable} is the
discrete counterpart of this congruence-based viewpoint.

\begin{remark}
\label{rem:banach-form-vs-congruence}
Up to this point, the definitions and closure arguments admit a parallel
formulation for arbitrary separable Polish structures with continuous operations,
using closed congruences $E\subseteq X^2$ rather than kernels $K\subseteq X$.
We do not develop that notation systematically in the sequel only because the
later applications repeatedly use quotient norms, distance-to-kernel functions,
and linear perturbation arguments.  Thus the paper is written in Banach form for
economy of notation, but whenever a later argument is purely formal it may be read
equally in the congruence language.
\end{remark}

\begin{example}\label{ex:standard-cats}
We list canonical examples fitting the framework of
Definition~\ref{def:category-ops}.

\begin{enumerate}
\item \emph{Banach algebras.}  $\omega_1(x,y)=xy$ (multiplication).
Admissible kernels are closed two-sided ideals.

\item \emph{$C^*$-algebras.}  $\omega_1(x,y)=xy$, $\omega_2(x)=x^*$.
Admissible kernels are closed $*$-ideals.

\item \emph{Banach lattices.}
$\omega_1(x,y)=x\vee y$, $\omega_2(x,y)=x\wedge y$, $\omega_3(x)=|x|$.
Admissible kernels are closed lattice ideals.

\item \emph{Ternary rings of operators (TROs).}
$\omega_1(x,y,z)=[x,y,z]=xy^*z$.
Admissible kernels are closed ternary ideals.

\item \emph{Operator systems (not treated here).}
Operator systems do not fit Definition~\ref{def:category-ops} as stated.
Their natural quotient theory is by completely order ideals (kernels of unital
completely positive maps), and an Archimedeanisation step is generally needed to
form the quotient operator system~\cite{KPTT}.  We therefore exclude operator
systems from the present framework.

\item \emph{Banach $A$-modules (encoding caveat).}
Fix a separable Banach algebra $A$ and a countable dense subset
$(a_m)$ of its unit ball.
A left Banach $A$-module may be encoded by the countable family of
unary operations $\omega_m(x)=a_mx$.  We do not use this variant below.

\item \emph{Jordan--Banach algebras (JB or JB$^*$).}
The fundamental operation is the Jordan product
$\omega_1(x,y)=\tfrac12(xy+yx)$; in the JB$^*$ case one adds
$\omega_2(x)=x^*$.
Admissible kernels are closed Jordan ideals.

\item \emph{Banach--Lie algebras.}
$\omega_1(x,y)=[x,y]$ (the Lie bracket).
Admissible kernels are closed Lie ideals.

\item \emph{Operator spaces (encoding caveat).}
Operator spaces are naturally described by a countable family of matrix-norm
predicates, rather than only by operations $X^{n_i}\to X$.  They therefore fall
outside the literal scope of Definition~\ref{def:category-ops}; a multi-sorted
or predicate-enriched version of the framework should cover them.
\end{enumerate}
\end{example}

\begin{remark}\label{rem:BL-separate}
While Banach lattices fit the general quotient-kernel framework at the level of
closed lattice ideals, the descriptive set-theoretic analysis of the natural Polish
parameter spaces of separable Banach lattices involves additional lattice-specific
topological subtleties.  These will be treated separately in forthcoming work of
M.~Niwi{\'n}ski~\cite{NiwinskiBL}.
\end{remark}

\subsection{Closedness of the admissible-kernel space}

\begin{proposition}\label{prop:ideal-closed}
Let $X\in\mathcal C$ be separable and let $\tau$ be an admissible topology on $\SB(X)$.
Then $\Ideal_{\mathcal C}(X)$ is $\tau$-closed in $\SB(X)$, hence Polish.
\end{proposition}

\begin{proof}
Fix a countable dense set $D^X=\{u_m:m\in\N\}$.  By Theorem~\ref{thm:GSR}, there
exist continuous selections $f_k:(\SB(X),\tau)\to X$ with $f_k(F)\in F$ and
$\overline{\{f_k(F):k\in\N\}}=F$ for all $F\in\SB(X)$.

For $i\in I$, $1\le j\le n_i$, and $\bar m=(m_1,\dots,m_{n_i})\in\N^{n_i}$ define
a continuous map $\Delta_{i,j,\bar m}:X\to X$ by
\[
\Delta_{i,j,\bar m}(x):=
\omega_i(u_{m_1},\dots,u_{m_{j-1}},u_{m_j}+x,u_{m_{j+1}},\dots,u_{m_{n_i}})
-\omega_i(u_{m_1},\dots,u_{m_{n_i}}).
\]
For $k\in\N$ set
\[
\mathcal K_{i,j,\bar m,k}
:=\bigl\{F\in\SB(X):\ \Delta_{i,j,\bar m}\bigl(f_k(F)\bigr)\in F\bigr\}.
\]
The map $F\mapsto \Delta_{i,j,\bar m}(f_k(F))$ is continuous since
$f_k$ and $\Delta_{i,j,\bar m}$ are continuous.
By admissibility, the membership relation $\{(x,F):x\in F\}\subseteq X\times\SB(X)$
is $\tau$-closed, hence each $\mathcal K_{i,j,\bar m,k}$ is $\tau$-closed.

Set $\mathcal K=\bigcap_{i,j,\bar m,k}\mathcal K_{i,j,\bar m,k}$,
a countable intersection of $\tau$-closed sets, hence $\tau$-closed.
We show $\mathcal K=\Ideal_{\mathcal C}(X)$.

Suppose that $F\in\Ideal_{\mathcal C}(X)$.  Then $f_k(F)\in F$ for all~$k$
(by Theorem~\ref{thm:GSR}), so each defining condition of $\mathcal K_{i,j,\bar m,k}$ holds.
Hence $F\in\mathcal K$.

Conversely, let $F\in\mathcal K$.  Fix $i$, $j$, elements
$x_1,\dots,x_{n_i}\in X$, and $h\in F$.
Choose $u_{m_\ell(t)}\to x_\ell$ with $u_{m_\ell(t)}\in D^X$,
and choose $f_{k(t)}(F)\to h$ (this is possible because $f_k(F)\in F$ for all~$k$ and
the selections $\{f_k(F):k\in\N\}$ are dense in~$F$ by Theorem~\ref{thm:GSR}).
Since $F\in\mathcal K$, for each~$t$,
\[
\Delta_{i,j,\bar m(t)}\bigl(f_{k(t)}(F)\bigr)\in F,
\]
where $\bar m(t)=(m_1(t),\dots,m_{n_i}(t))$.
By continuity of~$\omega_i$ and closedness of~$F$, the limit
$\omega_i(x_1,\dots,x_j+h,\dots,x_{n_i})-\omega_i(x_1,\dots,x_{n_i})$
lies in~$F$. This establishes invariance under perturbation of one coordinate by an element of~$F$.
To deduce the full admissible-kernel condition
($x_j-x_j'\in F$ for all~$j$ implies that indeed 
$\omega_i(x_1,\dots,x_{n_i})-\omega_i(x_1',\dots,x_{n_i}')\in F$),
telescope over coordinates: replace $x_1$ by $x_1'$ while fixing all other entries
(applying the single-coordinate invariance), then replace $x_2$ by $x_2'$, and so on.
After $n_i$ applications, $F$ absorbs the full difference.
Hence $F\in\Ideal_{\mathcal C}(X)$.

\end{proof}

\subsection{Quotient norms and definability}

\begin{definition}\label{def:Phi}
For $x\in X$ and $F\in\Ideal_{\mathcal C}(X)$, define
\[
\Phi(F,x)=d(x,F)=\inf_{z\in F}\|x-z\|.
\]
If $F$ is a subspace, $\Phi(F,x)=\|x+F\|_{X/F}$.
\end{definition}

\begin{lemma}\label{lem:Phi-continuous}
Under the Wijsman topology, for each fixed $x\in X$, the map
$F\mapsto \Phi(F,x)$ is continuous on $\Ideal_{\mathcal C}(X)$.
\end{lemma}

\begin{proof}
Immediate from Definition~\ref{def:wijsman}: Wijsman convergence
$F_n\to F$ means $d(x,F_n)\to d(x,F)$ for all~$x$.
\end{proof}

\begin{lemma}\label{lem:wijsman-subbasic}
Let $X$ be separable and let $\tau_W$ be the Wijsman topology on $\mathscr F(X)$.
For $x\in X$ and $r\in\Q$ the sets
\[
U(x,r):=\{F:\ d(x,F)<r\},\qquad
V(x,r):=\{F:\ d(x,F)\le r\}
\]
are $\tau_W$-open and $\tau_W$-closed respectively.
If $F\subseteq X$ is a closed subspace, then $d(x,F)=\Phi(F,x)$.
\end{lemma}

\begin{proof}
For each fixed $x$, the map $F\mapsto d(x,F)$ is continuous (Definition~\ref{def:wijsman}).
Thus $U(x,r)$ is the preimage of the open ray $(-\infty,r)$ and $V(x,r)$ is the
preimage of the closed ray $(-\infty,r]$.
\end{proof}

\begin{lemma}\label{lem:dense-op-closed}
Let $X\in\mathcal C$ be separable with countably many
operations $\omega_i$ of finite arity.  There exists a countable dense subset
$D\subseteq X$ closed under all operations $\omega_i$ and under rational
linear combinations.
\end{lemma}

\begin{proof}
Start with any countable dense $D_0$.  Define inductively $D_{n+1}$ to be the
set obtained from $D_n$ by closing under all operations~$\omega_i$ applied to tuples
from~$D_n$ and under rational linear combinations.  Put $D=\bigcup_n D_n$.
\end{proof}

\begin{definition}\label{def:projective}
An object $P\in\mathcal C$ is a \emph{separable quotient generator} if
$P$ is separable and every separable object $Y\in\mathcal C$ is isomorphic to
$P/K$ for some $K\in\Ideal_{\mathcal C}(P)$.
\end{definition}

\begin{remark}\label{rem:terminology}
In the categorical literature, a \emph{projective generator} additionally satisfies
a lifting property with respect to epimorphisms.  We do not require such a lifting
property; our arguments use only the quotient-universality encoded in
Definition~\ref{def:projective}.  We therefore use the more descriptive term
`separable quotient generator' throughout.
\end{remark}

\subsection{Main theorem and definability scheme}

\begin{mainthm}[Polish parameter space for continuous structures]\label{thm:continuous-main}
Let $\mathcal{C}$ be a~concrete category with continuous operations admitting a
separable quotient generator~$P$.  Fix:
\begin{itemize}
\item a countable dense subset $D^P=\{u_m:m\in\N\}\subseteq P$ closed under all
operations~$\omega_i$ and possible rational (or complex-rational, depending on the ground field) linear combinations\footnote{Formally, there may be no rational/complex-rational vector space structure in which case this condition becomes vacuous.};
\item an admissible topology $\tau$ on $\SB(P)$ (e.g.\ Wijsman).
\end{itemize}
Then:
\begin{enumerate}
\item \textbf{(Polishness.)}
$\Ideal_{\mathcal C}(P)$ is a Polish space with the topology~$\tau$.

\item \textbf{(Borel encoding.)}
Define a code $c(K)\in\R^\N$ for $K\in\Ideal_{\mathcal C}(P)$ by
concatenating all values
\[
\Phi\Bigl(K,\sum_{j=1}^\ell q_j u_{m_j}\Bigr)
\quad\text{and}\quad
\Phi\bigl(K,\omega_i(u_{m_1},\dots,u_{m_{n_i}})-u_r\bigr),
\]
enumerated over all rational data.  For any admissible topology~$\tau$,
each coordinate $K\mapsto\Phi(K,x)=d(x,K)$ is upper semicontinuous
(being the infimum of continuous functions via the
Godefroy--Saint-Raymond selections), hence Borel, and
$c:\Ideal_{\mathcal C}(P)\to\R^\N$ is Borel.
Under the \emph{Wijsman} topology, each coordinate is continuous.

\item \textbf{(Faithfulness.)}
If $c(K)=c(L)$, then $K=L$.  In particular, $P/K\cong P/L$ as objects in~$\mathcal C$
via the unique isometric isomorphism sending $u_m+K$ to $u_m+L$.
(\emph{Proof:} $c(K)=c(L)$ implies $\Phi(K,u_m)=\Phi(L,u_m)$ for all~$m$, hence
$d(\cdot,K)=d(\cdot,L)$ on a dense set, hence everywhere, so $K=L$.)

\item \textbf{(Surjectivity.)}
Every separable object $Y\in\mathcal C$ is isomorphic to $P/K$ for some
$K\in\Ideal_{\mathcal C}(P)$.  In particular, $c$ gives a Borel coding of
quotient presentations of separable objects in~$\mathcal C$.
We do \emph{not} claim injectivity on isomorphism classes: distinct kernels may
still yield isomorphic quotients.

\item \textbf{(Definability scheme.)}
Any property $\mathcal P$ of $P/K$ expressible by a formula built from:
\begin{itemize}
\item atomic predicates of the form
$F(\Phi(K,x_1),\ldots,\Phi(K,x_m))\bowtie r$,
where $x_j\in D^P$, $F:\R^m\to\R$ is continuous,
$\bowtie\in\{<,\le,=,\ge,>\}$, and $r\in\Q$;
\item Boolean connectives;
\item countable quantification over variables in $D^P$ or $\N$;
\end{itemize}
defines a Borel subset of $\Ideal_{\mathcal C}(P)$ for \emph{any} admissible
topology~$\tau$.  Under the \emph{Wijsman} topology specifically,
the Borel rank is bounded by the quantifier alternation depth:
$\forall$ alone gives~$\Pi^0_2$,
$\forall\exists$ gives~$\Pi^0_3$,
$\forall\exists\forall$ gives~$\Pi^0_4$, and so on.
\end{enumerate}
\end{mainthm}

\begin{proof}
Part~(1) is Proposition~\ref{prop:ideal-closed}.
For~(2), each coordinate map $K\mapsto d(x,K)$ is upper semicontinuous for any
admissible topology (as the infimum of continuous Godefroy--Saint-Raymond selections),
hence Borel; under the Wijsman topology it is continuous by
Lemma~\ref{lem:Phi-continuous}.
Part~(3): as argued in the statement, $c(K)=c(L)$ forces $K=L$.
Part~(4) holds because $P$ is a quotient generator.
Part~(5): for any admissible $\tau$, each atomic predicate is Borel
(since $K\mapsto\Phi(K,x)$ is u.s.c.\ by~(2)),
Boolean combinations and countable quantification preserve Borelness.
Under the Wijsman topology, each $\Phi$-coordinate is continuous,
so each atomic predicate is open (for~$<,>$)
or closed (for~$\le,\ge,=$), and countable
quantification over a countable set raises the Borel rank by at most one per alternation.
\end{proof}

\begin{definition}\label{def:stable-rel}
Fix $m\in\N$ and a finite family of $\mathcal C$-terms
$t_1(\bar x),\dots,t_r(\bar x)$ in variables $\bar x=(x_1,\dots,x_m)$.
For an object $Y\in\mathcal C$ and $\bar y\in Y^m$ set the \emph{defect}
\[
\mathrm{def}_{\bar t}^Y(\bar y):=\max_{\ell\le r}\|t_\ell^Y(\bar y)\|.
\]
Given $M\in\N$, we say that the relation
\[
\Sigma_{\bar t,M}(\bar x):\ \ \bigl(\|x_j\|\le M\ (1\le j\le m)\bigr)\ \wedge\
\bigl(t_1(\bar x)=\cdots=t_r(\bar x)=0\bigr)
\]
is \emph{stable on quotients of $P$} if for every $\varepsilon>0$ there exists
$\delta>0$ such that for every $K\in\Ideal_{\mathcal C}(P)$ and every
$\bar y\in (P/K)^m$ with $\|y_j\|\le M$,
\[
\mathrm{def}_{\bar t}^{P/K}(\bar y)<\delta\ \Longrightarrow\
\exists\bar z\in(P/K)^m:
\ \Sigma_{\bar t,M}(\bar z)\ \ \text{and}\ \
\max_j\|y_j-z_j\|<\varepsilon.
\]
\end{definition}

\begin{lemma}\label{lem:stable-exists-dense}
Assume $\Sigma_{\bar t,M}$ is stable on quotients of $P$.
Fix $\varepsilon=1$ and let $\delta>0$ be the corresponding stability constant.
Then for every $K\in\Ideal_{\mathcal C}(P)$ the following are equivalent:
\begin{enumerate}
\item there exists $\bar x\in (P/K)^m$ with $\Sigma_{\bar t,M}(\bar x)$;
\item there exists $\bar u\in (D^P)^m$ such that
$\max_j\Phi(K,u_j)< M$ and $\mathrm{def}_{\bar t}^{P/K}(\bar u+K)<\delta$.
\end{enumerate}
In particular, the set of $K$ satisfying $\exists\bar x\,\Sigma_{\bar t,M}(\bar x)$
is \emph{open} in the Wijsman topology on $\Ideal_{\mathcal C}(P)$.
\end{lemma}

\begin{proof}
(1)$\Rightarrow$(2).
Let $\bar x\in (P/K)^m$ be an exact solution with $\|x_j\|\le M$.
By continuity of the finitely many term maps $t_\ell^{P/K}$ at $\bar x$,
choose $\eta>0$ such that
\[
\max_j\|y_j-x_j\|<\eta \ \Longrightarrow\ 
\mathrm{def}_{\bar t}^{P/K}(\bar y)<\delta.
\]
Pick $\bar v\in(D^P)^m$ with $\max_j\|v_j+K-x_j\|<\eta/4$.
Then for each $j$,
\[
\Phi(K,v_j)=\|v_j+K\|\le \|x_j\|+\eta/4\le M+\eta/4.
\]

For each $j$, if $\Phi(K,v_j)<M$ set $u_j=v_j$.
Otherwise $M\le \Phi(K,v_j)\le M+\eta/4$.
Consider the open interval
\[
\Bigl(1-\frac{\eta/2}{\Phi(K,v_j)},\,\frac{M}{\Phi(K,v_j)}\Bigr).
\]
Its length equals
\[
\frac{M-\Phi(K,v_j)+\eta/2}{\Phi(K,v_j)}
\ \ge\ \frac{M-(M+\eta/4)+\eta/2}{\Phi(K,v_j)}
\ =\ \frac{\eta/4}{\Phi(K,v_j)}\ >\ 0,
\]
so the interval is non-empty. Choose a rational $\lambda_j$ in this interval and set
$u_j=\lambda_j v_j\in D^P$ (using closure of $D^P$ under rational scalars).
Then $\Phi(K,u_j)=\lambda_j\Phi(K,v_j)<M$ and
\[
\|u_j+K-(v_j+K)\|=|1-\lambda_j|\Phi(K,v_j)<\eta/2.
\]
Hence
\[
\|u_j+K-x_j\|\le \|u_j+K-(v_j+K)\|+\|v_j+K-x_j\|
<\eta/2+\eta/4<\eta,
\]
and therefore $\mathrm{def}_{\bar t}^{P/K}(\bar u+K)<\delta$.

(2)$\Rightarrow$(1).
Let $\bar u\in(D^P)^m$ satisfy the inequalities in~(2), and set $\bar y=\bar u+K$.
Then $\|y_j\|=\Phi(K,u_j)<M$ and $\mathrm{def}_{\bar t}^{P/K}(\bar y)<\delta$.
By stability there exists $\bar z$ with $\Sigma_{\bar t,M}(\bar z)$.

Finally, for fixed $\bar u$ the conditions
$\max_j\Phi(K,u_j)<M$ and $\mathrm{def}_{\bar t}^{P/K}(\bar u+K)<\delta$
are strict inequalities of continuous functions of $K$, hence open.
Taking the union over countably many $\bar u\in(D^P)^m$ yields openness.
\end{proof}

\begin{theorem}\label{thm:definability-rank}
Work in the setting of Theorem~\ref{thm:continuous-main} with the Wijsman topology.
Let $\mathcal L$ be the language generated by:

\begin{enumerate}[label=\textup{(\alph*)}]
\item \emph{quotient-norm atoms:}
predicates of the form
\[
F\bigl(\Phi(K,t_1),\dots,\Phi(K,t_m)\bigr)\ \bowtie\ r,
\]
where $t_j$ are $\mathcal C$-terms over $D^P$, $F:\R^m\to\R$ is continuous,
$r\in\Q$, and $\bowtie\in\{<,\le,=,\ge,>\}$;

\item \emph{stable-existence atoms:}
for each stable relation $\Sigma_{\bar t,M}$ (Definition~\ref{def:stable-rel}),
a predicate asserting the existence of a solution in $P/K$:
\[
\exists \bar x\ \Sigma_{\bar t,M}(\bar x).
\]
Equivalently (Lemma~\ref{lem:stable-exists-dense}), one may use a fixed rational
tolerance $\delta>0$ and quantify only over tuples from $D^P$.
\end{enumerate}

Build formulae from these atoms using Boolean connectives and countable quantifiers
over $D^P$ and~$\N$.

Then for every such formula $\varphi$, the satisfaction set
\[
[\varphi]=\{K\in\Ideal_{\mathcal C}(P):\ P/K\models\varphi\}
\]
is Borel in $(\Ideal_{\mathcal C}(P),\mathrm{Wijsman})$.
Moreover, if $\varphi$ is in prenex normal form with $s$ alternations, then the
following uniform upper bound holds:
$[\varphi]\in\Sigma^0_{2+s}$ or $[\varphi]\in\Pi^0_{2+s}$ according to the leading
quantifier block (with the convention that stable-existence atoms count as
\emph{open} atoms in the rank bookkeeping).  The estimate is deliberately an upper
bound; closed atomic predicates, such as equalities, often lower the actual rank.
\end{theorem}

\begin{proof}
By Lemma~\ref{lem:Phi-continuous}, each map $K\mapsto\Phi(K,t)$ is continuous.
Quotient-norm atoms of type~(a) are therefore either open (for~$<,>$) or closed
(for~$\le,\ge,=$) by Lemma~\ref{lem:wijsman-subbasic}.
Stable-existence atoms of type~(b) are open by
Lemma~\ref{lem:stable-exists-dense}.
Boolean combinations preserve Borelness.
Countable $\exists$ corresponds to a countable union, countable $\forall$ to a
countable intersection.  Each alternation raises the Borel rank by at most one.
\end{proof}

\begin{remark}\label{rem:dense-method-scope}
Theorem~\ref{thm:definability-rank} applies to properties that can be expressed using:
\begin{itemize}
\item norms of \emph{named terms} in the quotient (via $\Phi$), and
\item existential assertions whose defining relations are \emph{stable} on bounded balls,
so that exact witnesses can be replaced by approximate witnesses from the fixed
countable dense set $D^P$ (Lemma~\ref{lem:stable-exists-dense}).
\end{itemize}
This includes, for instance, the existence of projections, partial isometries,
matrix units, finite-dimensional $*$-subalgebras, and semiprojective relations, hence
the standard approximation/regularity properties treated later (AF, QD/MF, nuclear
dimension bounds, $D$-absorption, etc.); see Lemma~\ref{lem:stable-fd} and
Remark~\ref{rem:stable-relations}.

By contrast, properties whose definitions require \emph{unbounded} existential
quantification without a uniform stability modulus (\emph{e.g.}\ ring-theoretic phenomena
with unbounded witnesses in reduced products) or that intrinsically quantify over
\emph{uncountable} parameter sets may fall outside this scheme.
\end{remark}

\begin{remark}\label{rem:upper-bounds}
The ranks given in Theorem~\ref{thm:definability-rank} are \emph{uniform upper bounds}
that accommodate both open and closed atomic predicates.  In many concrete instances
the actual Borel class is lower, because particular atoms may be open (strict
inequalities) rather than merely closed (equalities or non-strict inequalities).
Throughout this paper, each stated complexity estimate should be read as a safe upper
bound unless explicitly accompanied by a matching hardness or completeness result.

When we write ``closed (so in particular $\Pi^0_2$)'', we mean the set is $\Pi^0_1$ (closed) and
therefore \emph{a fortiori} $\Pi^0_2$; we sometimes report $\Pi^0_2$ in
such cases because our bookkeeping is geared to quantifier alternation
starting at the $\Pi^0_2$ level, and we do not always optimise the base level.
\end{remark}

\begin{remark}\label{rem:encoding}
The code $c(K)$ captures the \emph{norm structure} of $P/K$ via quotient norms of
all rational combinations, and the \emph{operations} via
$\Phi(K,\omega_i(u_{m_1},\ldots)-u_r)$ (measuring when operations take specific values).
Since $D^P$ is dense and operation-closed, this information determines $P/K$ up to
isometric isomorphism.
\end{remark}

%%%%%%%%%%%%%%%%%%%%%%%%%%%%%%%%%%%%%%%%%%%%%%%%%%%%%%%%%%%%%%
\section{Main theorem for countable algebraic structures}\label{sec:countable}
%%%%%%%%%%%%%%%%%%%%%%%%%%%%%%%%%%%%%%%%%%%%%%%%%%%%%%%%%%%%%%

\subsection{Congruences on free algebras}

Let $L=\{f_i:i\in I\}$ be a countable finitary signature and let $F_\infty$ be the
free $L$-algebra on generators $\{x_n:n\in\N\}$.

\begin{theorem}
\label{thm:free-objects-closed}
Let $\mathcal V$ be a finitary algebraic variety over a countable signature
$L=\{f_i:i\in I\}$.  Let $F_\infty$ be the free $L$-algebra on the countable set of
generators $\{x_n:n\in\N\}$.  Consider the following Polish spaces, all with the
product (Cantor) topology:

\begin{enumerate}
\item \textbf{Congruences (general case).}
$\Con_L(F_\infty)\subseteq 2^{F_\infty\times F_\infty}$, the family of
$L$-congruences on~$F_\infty$.

\item \textbf{Groups.}
$\NSub(F_\infty)\subseteq 2^{F_\infty}$, the family of normal subgroups of the
free group on countably many generators, $F_\infty=\langle x_n:n\in\N\rangle$.

\item \textbf{Rings.}
$\Id(R_\infty)\subseteq 2^{R_\infty}$, the family of two-sided ideals in the
free unital ring generated by countably many free variables, $R_\infty=\Z\langle X_n:n\in\N\rangle$.

\item \textbf{Abelian groups.}
$\Sub(\Z^{(\N)})\subseteq 2^{\Z^{(\N)}}$, the family of subgroups of
$\Z^{(\N)}=\bigoplus_{n\in\N}\Z e_n$.

\item \textbf{Lattices.}
$\Con_{\wedge,\vee}(L_\infty)\subseteq 2^{L_\infty\times L_\infty}$, the family of
lattice congruences on the free lattice~$L_\infty$.
\end{enumerate}

Then each of the above parameter spaces is a \emph{closed} subset of the ambient
Cantor cube, hence compact, zero-dimensional, and Polish.
\end{theorem}

\begin{proof}
We treat~(1) in full generality; the remaining cases are specialisations.

A congruence $R\subseteq F_\infty\times F_\infty$ must be an equivalence relation
compatible with each operation.  Reflexivity, symmetry, and transitivity are
universal Horn conditions: for instance, transitivity states that for all $a,b,c$,
if $(a,b)\in R$ and $(b,c)\in R$ then $(a,c)\in R$; the violation for a fixed
triple $(a,b,c)$ forms the clopen cylinder
$\{R:(a,b)\in R,\ (b,c)\in R,\ (a,c)\notin R\}$.
Taking complements and intersecting over all triples yields a closed set.

For $L$-compatibility, fix an operation $f_i$ of arity~$n_i$ and tuples
$\bar a,\bar b\in F_\infty^{n_i}$.  The violating set
$\{R:\bigwedge_j (a_j,b_j)\in R\text{ but }(f_i(\bar a),f_i(\bar b))\notin R\}$
depends on finitely many coordinates and is clopen.  Taking complements and
intersecting gives a closed set.

Thus $\Con_L(F_\infty)$ is a countable intersection of closed sets, hence closed
in $2^{F_\infty\times F_\infty}$.  By compactness of the Cantor cube,
$\Con_L(F_\infty)$ is compact Polish.
\end{proof}

\begin{remark}\label{rem:abelian-closed}
In the group case the subclass of \emph{abelian groups} is the closed subspace
$\{N\in\NSub(F_\infty):[F_\infty,F_\infty]\subseteq N\}$
$=\bigcap_{g,h\in F_\infty}\{N:[g,h]\in N\}$,
and in the ring case \emph{commutative rings} form the closed subspace
$\{I\in\Id(R_\infty):\forall a,b,\ ab-ba\in I\}$.
\end{remark}

\subsection{Canonical encoding of quotients on a fixed domain}

\begin{definition}\label{def:canonical-encoding}
For a subobject~$S$ of the appropriate kind (congruence, normal subgroup, ideal, or
subgroup) such that $F_\infty/S$ is \emph{countably infinite},
we encode the quotient
canonically on the fixed domain~$\N$ as follows.  For each $a\in F_\infty$, let
$[a]_S$ be its $S$-coset or $S$-class.  Define
$\Min_S=\{\min([a]_S):a\in F_\infty\}\subseteq\N$,
the set of minimal representatives.
Since $F_\infty/S$ is infinite, $\Min_S$ is infinite;
let $\rho_S:\N\to\Min_S$ enumerate $\Min_S$ increasingly and let $q_S:F_\infty\to\N$
map each element to the index of its class representative.  For each basic operation
$f_i$ of arity $n_i$, define
\[
f_i^S(k_1,\dots,k_{n_i})
=q_S\bigl(f_i^{F_\infty}(\rho_S(k_1),\dots,\rho_S(k_{n_i}))\bigr).
\]
Then $\Theta(S)=(f_i^S)_{i\in I}$ is an $L$-structure on $\N$ isomorphic
to $F_\infty/S$.

\smallskip\noindent
\emph{Finite quotients.}\
If $F_\infty/S$ is finite (with $n$~elements, say), $\Min_S$ has exactly
$n$~elements and $\rho_S$ is not defined as a map $\N\to\Min_S$.
The canonical encoding $\Theta$ is therefore restricted to
infinite quotients; finite quotients are coded directly
by~$S$ (which determines the quotient structure up to isomorphism).
\end{definition}

\begin{lemma}\label{lem:rho-q-Borel}
On the $G_\delta$ subspace
$\mathcal S_\infty:=\{S\in\mathcal S:F_\infty/S\text{ is infinite}\}$,
the maps $S\mapsto\rho_S$ and $S\mapsto q_S$ are Borel, and hence $\Theta$ is Borel.
\end{lemma}

\begin{proof}
Fix $k\in\N$.  The value $\rho_S(k)$ is the least $a\in\N$ such that the family of
$S$-classes met by $\{0,1,\ldots,a\}$ has size $k{+}1$.

For $b,b'\in F_\infty$, the predicate ``$b$ and $b'$ represent the same element of $F_\infty/S$''
is clopen in the Cantor topology on~$\mathcal S$, since it is one of:
\begin{itemize}
\item $(b,b')\in S$ if $S$ is a congruence on $F_\infty$;
\item $bb'^{-1}\in S$ if $S$ is a normal subgroup of $F_\infty$;
\item $b-b'\in S$ if $S$ is an additive subgroup/ideal (abelian groups or rings).
\end{itemize}
Therefore the relation ``among $\{0,\ldots,a\}$ there are at least $k{+}1$ distinct $S$-classes''
is Borel in $S$ for each fixed $(a,k)$, and so $S\mapsto\rho_S(k)$ is Borel.

Similarly, for fixed $a\in F_\infty$, the value $q_S(a)$ is determined by finitely many
tests of whether $a$ is equivalent (mod~$S$) to $\rho_S(j)$ for $j\le a$, so $S\mapsto q_S(a)$ is Borel.
\end{proof}

\begin{mainthm}\label{thm:countable-main}
Let $\mathcal V$ and $F_\infty$ be as above, and let $\mathcal S$ denote any of the
closed coding spaces from Theorem~\ref{thm:free-objects-closed}.
Set
\[
\mathcal S_\infty:=\{S\in\mathcal S:\ F_\infty/S\text{ is infinite}\}.
\]
Then:
\begin{enumerate}
\item $\mathcal S$ is a compact zero-dimensional Polish space, and $\mathcal S_\infty$
is a $G_\delta$ subspace of $\mathcal S$ (hence Polish).

\item The map
\[
\Theta:\mathcal S_\infty\to \mathrm{Str}_L(\N)=\prod_{i\in I}\N^{\N^{n_i}},
\qquad
S\mapsto (f_i^S)_{i\in I},
\]
is Borel.

\item The induced map from $\mathcal S_\infty$ to the set of isomorphism classes
of countably infinite $L$-algebras in~$\mathcal V$ is surjective.  Equivalently,
for every countably infinite $\mathbf A\in\mathcal V$ there exists
$S\in\mathcal S_\infty$ with $\Theta(S)\cong \mathbf A$.

\item For each first-order sentence $\varphi$ in the language~$L$, the set
\[
\{S\in\mathcal S:\ F_\infty/S\models\varphi\}
\]
is Borel in~$\mathcal S$.
Quantifier-free sentences yield clopen sets; one leading existential quantifier
yields open; one leading universal quantifier yields closed; alternations give the
expected Borel ranks.
\end{enumerate}
\end{mainthm}

\begin{proof}
Part~(1): $\mathcal S$ is compact Polish by Theorem~\ref{thm:free-objects-closed}.
Moreover,
\[
\mathcal S_\infty
=\bigcap_{n\in\N}\{S:\ F_\infty/S\text{ has at least }n\text{ classes}\}
\]
is a countable intersection of open sets, hence $G_\delta$.

Part~(2) is Lemma~\ref{lem:rho-q-Borel}.
Part~(3) holds by the universal property of the free object.
Part~(4): atomic formulae are decided by finitely many coordinates of~$S$,
hence clopen.  Existential quantification adds a countable union; universal
quantification adds a countable intersection; further alternations raise the Borel rank accordingly.
\end{proof}

\begin{remark}
\label{rem:standard-countable-coding}
For a countable finitary signature $L=\{f_i:i\in I\}$, a classical parameter space for
countable $L$-structures is the Polish space
\[
\mathrm{Str}_L(\N):=\prod_{i\in I}\N^{\N^{n_i}},
\]
equipped with the product of the discrete topologies, where a point
$\mathbf A=(f_i^{\mathbf A})_{i\in I}$ represents the $L$-structure on domain $\N$ with
basic operations $(f_i^{\mathbf A})$.
The descriptive set theory of countable structures is typically developed in this
coding (together with the natural logic action of $S_\infty$ and its orbit equivalence
relation of isomorphism); see, for instance,
Kechris~\cite{Kechris}, Becker--Kechris~\cite{BeckerKechris},
Friedman--Stanley~\cite{FriedmanStanley}, Hjorth~\cite{Hjorth}, and Gao~\cite{GaoDST}.

Our quotient coding via congruences on the free algebra $F_\infty$ is compatible with
the usual $\mathrm{Str}_L(\N)$-coding for \emph{infinite} quotients.
(Recall that $\mathrm{Str}_L(\N)$ codes structures on the countably infinite domain~$\N$;
finite quotients are handled directly by the congruence~$S$ and do not enter the
$\mathrm{Str}_L(\N)$ translation.)
Concretely, the map $\Theta$ of Definition~\ref{def:canonical-encoding} produces from
a congruence $S$ a canonical quotient structure on $\N$.
Conversely, given $\mathbf A\in\mathrm{Str}_L(\N)$, the evaluation homomorphism
$F_\infty\to \mathbf A$ sending $x_n\mapsto n$ has kernel a congruence, and this
assignment is Borel.
Thus one may freely translate Borel-complexity statements between the compact
congruence/normal-subgroup spaces used here and the more customary
$\mathrm{Str}_L(\N)$ parameter spaces.
\end{remark}

\begin{remark}\label{rem:novikov-coding}
In algorithmic group theory, Novikov's proof of the existence of finitely presented
groups with unsolvable word problem proceeds by manipulating presentations and relators
encoded as finite words in a free group; see~\cite{Novikov1955}.
From a modern descriptive-set-theoretic perspective, a (finite or countable) group
presentation $\langle X\mid R\rangle$ is determined by the subset $R$ of the countable
set of reduced words in the free group $F(X)$, hence by a point of the Cantor space
$2^{F(X)}$.
Passing to the normal closure $\langle\!\langle R\rangle\!\rangle\triangleleft F(X)$
yields a normal subgroup and therefore a quotient group
$F(X)/\langle\!\langle R\rangle\!\rangle$.
In this sense, Novikov's presentation-based arguments may be viewed (in hindsight) as
working effectively inside what we would now regard as a natural standard Borel
parameter space of countable groups \emph({e.g.}\ a subspace of $\NSub(F_\infty)$ as in
Theorem~\ref{thm:free-objects-closed}).

At the same time, Novikov's paper is not phrased in the language of Polish spaces or
Borel sets: the topology/standard-Borel viewpoint and its systematic deployment for
definability and classification problems belongs to the later descriptive set theory
literature (\emph{cf}.\ \cite{Kechris,BeckerKechris,Hjorth,GaoDST}).
\end{remark}

\begin{remark}\label{rem:disc-vs-cont}
For rings, commutativity is \emph{closed} (hence~$\Pi^0_1$):
\[
        \forall a,b\quad (ab,ba)\in\theta .
\]
For Banach algebras the corresponding condition
$\forall a,b\in D^A,\ \Phi(K,ab-ba)=0$
is~$\Pi^0_2$, since $K\mapsto\Phi(K,ab-ba)$ is continuous and universal
quantification over the countable dense set raises the complexity by one level.
\end{remark}

%%%%%%%%%%%%%%%%%%%%%%%%%%%%%%%%%%%%%%%%%%%%%%%%%%%%%%%%%%%%%%
\section{Examples of quotient generators}\label{sec:examples}
%%%%%%%%%%%%%%%%%%%%%%%%%%%%%%%%%%%%%%%%%%%%%%%%%%%%%%%%%%%%%%

\begin{remark}\label{rem:proj-gens}
We list standard choices of~$P$; all are separable and surject onto every separable
object in the category via a morphism whose kernel is the appropriate admissible ideal.

\begin{enumerate}[label=(\alph*), itemsep=.4em]
\item \emph{Banach algebras.}
Every separable Banach algebra is generated by a sequence of elements of norm at
most~$1$.  By linearising and completing the corresponding free object, one obtains a
projectively universal quotient generator.  Specifically, for the category of
\emph{non-unital} Banach algebras, let $S$ be the free semigroup on countably many
generators and take $P=\ell_1(S)$.  For the category of \emph{unital} Banach algebras,
one must instead take $S$ to be the free monoid on countably many generators (so the
empty word acts as the multiplicative identity), yielding the unital quotient generator
$P=\ell_1(S)$.  The fact that these free objects act as universally free quotient
generators is a direct consequence of general categorical principles surrounding UFOs;
see the author and Balcerzak~\cite{BalcerzakKania2023}.

\item \emph{$C^*$-algebras.}
Let $F_\infty$ be the free group on countably many generators and take
$P=C^*_{\max}(F_\infty)$ \cite[{\S}V.2.1]{Blackadar}.
For non-unital algebras one may take $C^*_{\max}(F_\infty)\otimes\mathcal K$.

\item \emph{Banach lattices.}
Take $P=\mathrm{FBL}(\ell_1)$; see \cite{FreeBLl1}.

\item \emph{TROs.}
A free separable TRO exists (Section~\ref{sec:TRO}) and is a quotient generator.

\item \emph{Operator systems.}
The free operator system on countably many generators exists~\cite{KPTT}.

\item \emph{Operator spaces.}
Pisier's non-commutative $\ell^1$,
$P=(\bigoplus_{n\in\N} M_n)_{\ell_1}$, is a quotient generator
\cite{EffrosRuanBook,PisierOS}.

\item \emph{Banach modules over a fixed separable Banach algebra~$A$.}
One may take $P=\ell_1(\N;A)$ with the obvious module structure.
\end{enumerate}
\end{remark}

%%%%%%%%%%%%%%%%%%%%%%%%%%%%%%%%%%%%%%%%%%%%%%%%%%%%%%%%%%%%%%
\section{Banach algebras: complexity bounds}\label{sec:banach-complexity}
%%%%%%%%%%%%%%%%%%%%%%%%%%%%%%%%%%%%%%%%%%%%%%%%%%%%%%%%%%%%%%

Let $A$ be a fixed separable Banach-algebra quotient generator and work on
$\Ideal(A)$ with the Wijsman topology.
Fix a countable dense multiplicatively closed $D^A$.

\begin{remark}\label{rem:rank-convention}
All admissible topologies on a given ideal space generate the same standard Borel
$\sigma$-algebra~\cite{GodefroySaintRaymond2018}.  In particular, the statement
``property $\mathcal P$ is Borel'' is \emph{topology-independent}: it holds for
one admissible topology iff it holds for all.
However, the \emph{Borel rank} --- the precise level $\Sigma^0_\alpha$ or
$\Pi^0_\alpha$ at which $\mathcal P$ sits --- depends on the choice of topology,
because it involves the interplay between open and closed sets, which can differ
from one admissible topology to another.

Unless explicitly stated otherwise, all references to specific Borel classes
$\Sigma^0_\alpha$, $\Pi^0_\alpha$ throughout
Sections~\ref{sec:banach-complexity}--\ref{sec:Cstar-properties} are with respect to
the Wijsman topology on the relevant ideal space $\Ideal(P)$.  The Wijsman topology
is the natural choice for rank computations because it makes the quotient-norm
functional $\Phi$ continuous (not merely Borel), so that strict inequalities in
$\Phi$ define open sets and non-strict inequalities define closed sets.
\end{remark}

\begin{convention}\label{conv:proper-quotients}
Whenever a property below explicitly involves the unit
(topological stable rank, direct finiteness, stable finiteness, simplicity, tracial
states, and so on), we tacitly restrict to proper ideals $K\subsetneq A$.
Equivalently, we ignore the singleton corresponding to the zero quotient.
Since that singleton is closed, none of the stated upper bounds is affected.
\end{convention}

\subsection{Commutativity and uniformity}

\begin{proposition}\label{prop:comm}
The set of $K$ such that $A/K$ is commutative is closed (so in particular $\Pi^0_2$) in the Wijsman topology.
\end{proposition}

\begin{proof}
$A/K$ is commutative if and only if $\Phi(K,uv-vu)=0$ for all $u,v\in D^A$.
Each condition is closed and we take a countable intersection.
\end{proof}

\begin{proposition}\label{prop:unif}
The set of $K$ such that $A/K$ is an abstract uniform algebra is closed in the Wijsman
topology.
\end{proposition}

\begin{proof}
Here ``abstract uniform algebra'' means a Banach algebra whose norm satisfies
\[
        \|z^2\|=\|z\|^2 \qquad(z\in B).
\]
Equivalently, by the standard abstract characterisation of uniform algebras, such a
Banach algebra is isometrically isomorphic to a closed subalgebra of a commutative
$C(K)$ with the supremum norm.  Thus the descriptive-set-theoretic condition to be
checked is exactly the square-norm identity.

Let
\[
        E_K=\{z\in A/K:\|z^2\|=\|z\|^2\}.
\]
Since multiplication and the norm are continuous in $A/K$, the set $E_K$ is closed in
$A/K$.  Hence the identity holds for every element of $A/K$ if and only if it holds on
the dense set $D^A+K$.  For each fixed $u\in D^A$, the condition
\[
        \Phi(K,u^2)=\Phi(K,u)^2
\]
is closed, since both sides are continuous functions of~$K$ in the Wijsman topology.
The desired class is the countable intersection of these closed sets.
\end{proof}

\subsection{Right topological stable rank}

The \emph{topological stable rank} of a unital Banach algebra $B$, denoted $\mathrm{tsr}(B)$,
is the smallest integer $n\ge 1$ such that the set of right-unimodular $n$-tuples
$U_n(B)$ is dense in $B^n$ (see Rieffel~\cite{Rieffel1983}).  For $C^*$-algebras, the
topological stable rank coincides with the Bass stable
rank~\cite{HermanVaserstein1984}.

\begin{proposition}\label{prop:tsr}
Fix $n\in\N$.  The set of $K\in\Ideal(A)$ such that $A/K$ is unital and
$\mathrm{tsr}(A/K)\le n$ is $G_\delta$ in the Wijsman topology.
\end{proposition}

\begin{proof}
Write $B=A/K$.  The condition $\mathrm{tsr}(B)\le n$ says that the right-unimodular
$n$-tuples are dense in~$B^n$.  It is enough to approximate tuples from $(D^A)^n$.
For fixed $\bar a=(a_1,\ldots,a_n)\in(D^A)^n$ and $m\in\N$, let $U_{\bar a,m}$ be
the set of all $K$ for which there exist $\bar b,\bar v\in(D^A)^n$ such that
\[
        \max_j \Phi(K,a_j-b_j)<\frac1m,
        \qquad
        \Phi\Bigl(K,\sum_{j=1}^n b_jv_j-1\Bigr)<1 .
\]
For fixed witnesses this is open, and therefore $U_{\bar a,m}$ is open.  If the second
inequality holds in $B$, then $s=\sum_j(b_j+K)(v_j+K)$ is invertible and
\[
        \sum_j (b_j+K)\bigl((v_j+K)s^{-1}\bigr)=1,
\]
so $\bar b+K$ is right-unimodular.  Conversely, if a right-unimodular tuple
approximating $\bar a+K$ exists, approximate both that tuple and a right inverse by
members of $D^A$ to obtain the displayed strict inequalities.  Hence
\[
        \{K:\mathrm{tsr}(A/K)\le n\}
        =\bigcap_{\bar a\in(D^A)^n}\bigcap_{m\in\N} U_{\bar a,m},
\]
a countable intersection of open sets.
\end{proof}

\subsection{Uniformly open multiplication}

Uniformly open multiplication in Banach algebras has been studied
in connection with topological stable rank and differential
embeddings; see~\cite{DragaKania,KaniaMaslany}.

\begin{definition}\label{def:moduli}
Let
\[
\mathcal D:=\{\delta:(0,1]\cap\Q\to(0,1]\cap\Q:\delta \text{ is non-decreasing}\}.
\]
We equip the range $(0,1]\cap\Q$ with the discrete topology and endow
$\mathcal D$ with the subspace topology inherited from the product space
\[
\prod_{r\in(0,1]\cap\Q}\bigl((0,1]\cap\Q\bigr)_{\mathrm{disc}}.
\]
\end{definition}

\begin{remark}\label{rem:moduli-polish}
The space $\mathcal D$ is Polish.  Indeed, since $(0,1]\cap\Q$ is countably infinite,
the ambient product
\[
\prod_{r\in(0,1]\cap\Q}\bigl((0,1]\cap\Q\bigr)_{\mathrm{disc}}
\]
is homeomorphic to the Baire space $\N^\N$.
Moreover,
\[
\mathcal D=
\bigcap_{\substack{r,s\in(0,1]\cap\Q\\ r\le s}}
\{\delta:\delta(r)\le \delta(s)\},
\]
and for each fixed pair $(r,s)$ the set on the right is closed in the product topology
(since the coordinate spaces are discrete).  Hence $\mathcal D$ is a closed subspace of
a Polish space, and therefore Polish.
\end{remark}

\begin{definition}\label{def:UAO-UO}
Let $(X,d_X)$ and $(Y,d_Y)$ be metric spaces and let $f:X\to Y$ be continuous.
Fix $\delta\in\mathcal D$.

\begin{enumerate}
\item We say that $f$ is \emph{$\delta$-uniformly almost open} if for every
$r\in(0,1]\cap\Q$ and every $x\in X$,
\[
B_Y\bigl(f(x),r\bigr)\subseteq \overline{f\bigl(B_X(x,\delta(r))\bigr)}.
\]

\item We say that $f$ is \emph{$\delta$-uniformly open} if for every
$r\in(0,1]\cap\Q$ and every $x\in X$,
\[
B_Y\bigl(f(x),r\bigr)\subseteq f\bigl(B_X(x,\delta(r))\bigr).
\]

\item We say that $f$ is \emph{uniformly almost open} (resp.\ \emph{uniformly open})
if it is $\delta$-uniformly almost open (resp.\ $\delta$-uniformly open) for some
$\delta\in\mathcal D$.
\end{enumerate}
\end{definition}

\begin{lemma}\label{lem:schauder}
Let $X$ and $Y$ be metric spaces and let $f:X\to Y$ be continuous.
If $X$ is complete and $f$ is uniformly almost open, then $f$ is uniformly open.
\end{lemma}

\begin{proof}
This is Schauder's lemma; see the author and Draga~\cite[Lemma~2.1]{DragaKania}.
\end{proof}

\begin{lemma}\label{lem:mult-balls-cont}
Let $B$ be a Banach algebra and let $m:B\times B\to B$ be multiplication.
Equip $B\times B$ with the max metric
\[
d_\infty\bigl((a,b),(a',b')\bigr):=\max\{\|a-a'\|,\|b-b'\|\}.
\]
Fix $\varepsilon>0$ and define a set-valued map
\[
F_\varepsilon:B\times B\to\mathscr F(B),\qquad
F_\varepsilon(a,b):=\overline{m\bigl(B_\varepsilon(a)\times B_\varepsilon(b)\bigr)}.
\]
Then $F_\varepsilon$ is Wijsman-continuous: for each fixed $y\in B$ the map
$(a,b)\mapsto \dist\bigl(y,\,F_\varepsilon(a,b)\bigr)$
is continuous on $B\times B$.
\end{lemma}

\begin{proof}
Let $(a_n,b_n)\to(a,b)$ in $d_\infty$ and set $u_n:=a_n-a$, $v_n:=b_n-b$.
Fix $x\in B_\varepsilon(a)$ and $y\in B_\varepsilon(b)$. Then
$x+u_n\in B_\varepsilon(a_n)$, $y+v_n\in B_\varepsilon(b_n)$, and
\[
(x+u_n)(y+v_n)-xy = x v_n + u_n y + u_n v_n.
\]
Since $x$ and $y$ range over the bounded sets $B_\varepsilon(a)$ and $B_\varepsilon(b)$,
there is a constant $C=C(a,b,\varepsilon)$ such that
\[
\sup_{\substack{x\in B_\varepsilon(a)\\y\in B_\varepsilon(b)}}
\|(x+u_n)(y+v_n)-xy\|
\le C\bigl(\|u_n\|+\|v_n\|\bigr)+\|u_n\|\,\|v_n\|\to 0.
\]
The same estimate holds in the reverse direction (shifting from $(a_n,b_n)$ to $(a,b)$).
Hence the Hausdorff distance between the bounded sets
$m(B_\varepsilon(a_n)\times B_\varepsilon(b_n))$ and
$m(B_\varepsilon(a)\times B_\varepsilon(b))$
tends to $0$, and therefore $F_\varepsilon(a_n,b_n)\to F_\varepsilon(a,b)$
in the Wijsman topology.
\end{proof}

\begin{proposition}\label{prop:UO}
Let $A$ be a separable Banach-algebra quotient generator and work on $\Ideal(A)$
with the Wijsman topology. Fix a countable dense multiplicatively closed $D^A$ and fix
$\delta\in\mathcal D$.

\begin{enumerate}
\item The set of $K\in\Ideal(A)$ such that multiplication in $B=A/K$ is
$\delta$-uniformly almost open is $\Pi^0_3$.

\item The set of $K\in\Ideal(A)$ such that multiplication in $A/K$ is
uniformly open (\emph{i.e.}\ admits some modulus) is analytic.
\end{enumerate}
\end{proposition}

\begin{proof}
Write $B=A/K$ and let $m_B:B\times B\to B$ be multiplication, equipped with the max metric
on $B\times B$.

We begin by constructing a countable $\Phi$-scheme that captures
$\delta$-uniform almost openness.
Fix $r\in(0,1]\cap\Q$, $a,b,y\in D^A$, and $n\in\N$.
Consider the condition on $K$:
\begin{multline}\label{eq:UAO-scheme}
\Phi(K,y-ab)\ge r
\quad\lor\\
\exists a',b'\in D^A:\ 
\Phi(K,a'-a)<\delta(r)\ \wedge\ \Phi(K,b'-b)<\delta(r)
\ \wedge\ \Phi(K,a'b'-y)<\tfrac1n .
\end{multline}
For fixed $a',b'$, the conjunction on the right is a finite intersection of strict
inequalities in Wijsman-continuous functions of $K$, hence is open; therefore the right
disjunct is open (countable union over $a',b'\in D^A$).
The left disjunct is closed. Hence \eqref{eq:UAO-scheme} defines a $\Sigma^0_2$ subset of
$\Ideal(A)$.

Let $\mathcal U_\delta$ be the intersection of \eqref{eq:UAO-scheme} over all
$(r,a,b,y,n)\in ((0,1]\cap\Q)\times(D^A)^3\times\N$.
Then $\mathcal U_\delta$ is a countable intersection of $\Sigma^0_2$ sets, hence
$\mathcal U_\delta\in\Pi^0_3$.

We now verify that $\mathcal U_\delta$ coincides with
$\delta$-uniform almost openness of $m_B$.
Fix $r\in(0,1]\cap\Q$ and set $\varepsilon=\delta(r)$.
For $a,b\in B$, define the closed set
\[
F_\varepsilon(a,b):=\overline{m_B\bigl(B_\varepsilon(a)\times B_\varepsilon(b)\bigr)}
\subseteq B.
\]
By Lemma~\ref{lem:mult-balls-cont}, for fixed $\varepsilon$ the map
$(a,b)\mapsto F_\varepsilon(a,b)$
is Wijsman-continuous, and hence the function
$(a,b,z)\mapsto \dist\bigl(z,F_\varepsilon(a,b)\bigr)$
is continuous on $B\times B\times B$.

Assume $K\in\mathcal U_\delta$.
Take $a,b\in D^A$ (viewed in $B$) and $z\in B$ with $\|z-ab\|<r$.
Since $D^A+K$ is dense in $B$ and $F_\varepsilon(a,b)$ is closed, it suffices to show
$z\in F_\varepsilon(a,b)$ for $z$ ranging over $D^A+K$.
Fix $y\in D^A$ with $\Phi(K,y-ab)<r$.
Then for each $n$ the left disjunct in \eqref{eq:UAO-scheme} fails, and we obtain
$a'_n,b'_n\in D^A$ with
\[
\Phi(K,a'_n-a)<\varepsilon,\qquad \Phi(K,b'_n-b)<\varepsilon,\qquad
\Phi(K,a'_nb'_n-y)<\tfrac1n.
\]
In $B$ this means $(a'_n+K)\in B_\varepsilon(a+K)$, $(b'_n+K)\in B_\varepsilon(b+K)$ and
$\|(a'_n+K)(b'_n+K)-(y+K)\|<1/n$, hence $y+K\in F_\varepsilon(a+K,b+K)$.
Density of $D^A+K$ now gives $B_r(ab)\subseteq F_\varepsilon(a,b)$ for all
$a,b\in D^A+K$.

If $\delta$-uniform almost openness failed somewhere in $B\times B$, there would exist
$a,b\in B$ and $z\in B$ with $\|z-ab\|<r$ but $\dist(z,F_\varepsilon(a,b))>0$.
By continuity of $(a,b,z)\mapsto (\|z-ab\|,\dist(z,F_\varepsilon(a,b)))$ and density of
$D^A+K$, such a failure would already occur for some $a,b,z\in D^A+K$, a contradiction.
Therefore multiplication in $B$ is $\delta$-uniformly almost open.

Conversely, if multiplication in $B$ is $\delta$-uniformly almost open, then for any
$r,a,b,y,n$ with $\Phi(K,y-ab)<r$ we have $y+K\in F_{\delta(r)}(a+K,b+K)$, so there exist
$a',b'\in B$ within $\delta(r)$ of $a,b$ with $(a'b')+K$ within $1/n$ of $y+K$.
Approximating $a',b'$ by elements of $D^A$ inside the open balls and using continuity of
multiplication yields the right-hand disjunct in \eqref{eq:UAO-scheme}. Hence
$K\in\mathcal U_\delta$.
This proves~(1).

For genuine uniform openness, let
\[
\mathcal U:=\bigl\{(\delta,K)\in\mathcal D\times\Ideal(A):
K\in\mathcal U_\delta\bigr\}.
\]
Unwinding \eqref{eq:UAO-scheme} shows $\mathcal U$ is Borel in $\mathcal D\times\Ideal(A)$.
Hence its projection $\pi(\mathcal U)\subseteq\Ideal(A)$ is analytic.
If multiplication in $A/K$ is uniformly open, it is in particular uniformly almost open, so
$K\in\pi(\mathcal U)$.
Conversely, if $K\in\pi(\mathcal U)$ then multiplication in $A/K$ is uniformly almost open;
since $(A/K)\times(A/K)$ is complete, Lemma~\ref{lem:schauder} yields that multiplication is
uniformly open.
This proves~(2).
\end{proof}

\subsection{Pure infiniteness and norm control}\label{subsec:PI}

Let $B$ be a unital Banach algebra.  Following Daws--Horv\'ath~\cite{DawsHorvath},
define for $a\in B$
\[
C_{\mathrm{pi}}^B(a):=\inf\{\|b\|\,\|c\|:\ b,c\in B,\ bac=1\}\in(0,\infty].
\]
In particular, $C_{\mathrm{pi}}^B(a)<\infty$ iff $1\in BaB$.

\begin{definition}\label{def:PI-unif}
A unital Banach algebra $B$ is \emph{purely infinite} (in the sense of
\cite{DawsHorvath}) if $B\not\cong\C$ and $C_{\mathrm{pi}}^B(a)<\infty$ for every
non-zero $a\in B$.

We say that $B$ is \emph{uniformly purely infinite} if $B\not\cong\C$ and there exists
$M\in\N$ such that
\[
\sup\{C_{\mathrm{pi}}^B(a):\ \|a\|=1\}\le M.
\]
By \cite[Prop.~2.2]{DawsHorvath}, uniform pure infiniteness is equivalent to pure
infiniteness of every ultrapower $B^{\mathcal U}$ for countably incomplete ultrafilters.
\end{definition}

\begin{proposition}\label{prop:UPI}
Let $A$ be a unital Banach-algebra quotient generator and work on $\Ideal(A)$ with
the Wijsman topology.  The set of $K$ such that $A/K$ is uniformly purely infinite
is $\Sigma^0_3$.
\end{proposition}

\begin{proof}
Write $B=A/K$.  For $N\in\N$ set
\[
r_N:=\frac{1}{8N^2}\in(0,1)
\]
and let $\mathcal U_N\subseteq\Ideal(A)$ be the set of all $K$ such that
\begin{multline}\label{eq:UPI-net}
\forall a\in D^A:\ 
\bigl|\Phi(K,a)-1\bigr|>r_N
\quad\lor\\
\exists b,c\in D^A:\ 
\Phi(K,b)<N,\ \Phi(K,c)<N,\ \Phi(K,bac-1)<\tfrac18.
\end{multline}
For fixed $a\in D^A$, both disjuncts in \eqref{eq:UPI-net} define open subsets of
$\Ideal(A)$, so $\mathcal U_N$ is $G_\delta$, \emph{i.e.}\ $\Pi^0_2$.

Let $\mathrm{NonScal}$ denote the set of $K$ such that $A/K\not\cong\C$.
This set is $F_\sigma$: indeed,
\[
K\in\mathrm{NonScal}
\iff
\exists a\in D^A\ \exists m\in\N\ \forall \lambda\in\Q(i):\ \Phi(K,a-\lambda 1)\ge \tfrac1m.
\]
For fixed $(a,m)$ this is the superlevel condition
\[
\inf_{\lambda\in\Q(i)} \Phi(K,a-\lambda 1)\ge \tfrac1m.
\]
Since each map $K\mapsto \Phi(K,a-\lambda 1)$ is continuous in the Wijsman topology,
the infimum over the countable set $\Q(i)$ is upper semicontinuous.  Hence its
superlevel set is closed.  Taking the countable union over $(a,m)\in D^A\times\N$
shows that $\mathrm{NonScal}$ is $F_\sigma$.

We \emph{claim} that
\[
\mathrm{UPI}=\mathrm{NonScal}\cap\bigcup_{N\in\N}\mathcal U_N.
\]

\smallskip\noindent
($\Rightarrow$)
Assume that $B$ is uniformly purely infinite.
Choose $M\in\N$ with
$\sup\{C_{\mathrm{pi}}^B(x):\ \|x\|=1\}\le M$,
and then choose $N\in\N$ with $N^2>M$.
Fix $a\in D^A$ with $|\|a+K\|-1|\le r_N$ and put
$u:=(a+K)/\|a+K\|$.
Then $\|u\|=1$ and $\|u-(a+K)\|\le r_N$.
Since $C_{\mathrm{pi}}^B(u)\le M$, there exist $b_0,c_0\in B$ with $b_0uc_0=1$ and
$\|b_0\|\,\|c_0\|\le M$.
Rescale by
\[
\lambda:=\Bigl(\frac{\|c_0\|}{\|b_0\|}\Bigr)^{1/2},
\qquad
b:=\lambda b_0,
\qquad
c:=\lambda^{-1}c_0.
\]
Then $buc=1$ and
$\|b\|=\|c\|=\sqrt{\|b_0\|\,\|c_0\|}\le \sqrt M < N$.
Hence
\[
\|b(a+K)c-1\|
=\|b(a+K-u)c\|
\le \|b\|\,\|a+K-u\|\,\|c\|
< N^2 r_N
=\tfrac18.
\]
Approximating $b,c$ by elements of $D^A$ preserves the strict bounds
$\Phi(K,b)<N$, $\Phi(K,c)<N$, and $\Phi(K,bac-1)<1/8$.
Thus $K\in\mathcal U_N$.
Since $B\not\cong\C$ by definition, $K\in\mathrm{NonScal}$ as well.

\smallskip\noindent
($\Leftarrow$)
Assume $K\in\mathrm{NonScal}\cap\mathcal U_N$.
Let $x\in B$ with $\|x\|=1$ and choose $a\in D^A$ with
$\|a+K-x\|<r_N/2$.
Then $|\|a+K\|-1|\le r_N/2$, so the left disjunct in \eqref{eq:UPI-net} fails.
Hence there exist $b,c\in D^A$ with
$\|b+K\|<N$, $\|c+K\|<N$, and $\|b(a+K)c-1\|<1/8$.
Therefore
\[
\|bxc-1\|
\le \|b(x-a-K)c\|+\|b(a+K)c-1\|
< N^2\cdot\frac{r_N}{2}+\frac18
=\frac{3}{16}
<\frac12.
\]
So $bxc$ is invertible and $\|(bxc)^{-1}\|<2$.
Putting $c':=c(bxc)^{-1}$ gives $bxc'=1$ and
$\|b\|\,\|c'\|\le \|b\|\,\|c\|\,\|(bxc)^{-1}\|< 2N^2$.
Thus $C_{\mathrm{pi}}^B(x)\le 2N^2$ for every $\|x\|=1$.
Since $K\in\mathrm{NonScal}$, we also have $B\not\cong\C$.
Therefore $B$ is uniformly purely infinite.

\medskip\noindent
Finally, $\mathrm{NonScal}$ is $\Sigma^0_2$ and $\bigcup_N\mathcal U_N$ is
$\Sigma^0_3$, so
$\mathrm{UPI}=\mathrm{NonScal}\cap\bigcup_{N\in\N}\mathcal U_N$
is $\Sigma^0_3$.
\end{proof}

\begin{remark}\label{rem:PI-nonuniform}
The bare condition ``$C_{\mathrm{pi}}^{A/K}(a)<\infty$ for all non-zero $a$'' involves
existential witnesses with \emph{no uniform norm control}.  As in the directly finite
case, such unbounded-witness phenomena can
obstruct reductions to a single bounded fragment.  We therefore isolate the uniform
variant (Definition~\ref{def:PI-unif}), which is the version naturally compatible with
the dense-set quantification method and with ultrapower permanence~\cite{DawsHorvath}.
At present we do not know whether the bare non-uniform class of purely infinite quotients
is Borel, analytic, or genuinely more complicated in this coding.
\end{remark}

\subsection{Dedekind (in)finiteness and unbounded witnesses}\label{subsec:df-banach}

A unital ring (or Banach algebra) $B$ is \emph{Dedekind finite} (or \emph{directly
finite}) if $ab=1$ implies $ba=1$ for all $a,b\in B$; otherwise $B$ is
\emph{Dedekind infinite}.

In Banach algebras, a central subtlety is that direct finiteness is not controlled by
uniformly bounded fragments: Daws--Horv\'ath show that ring-theoretic (in)finiteness
behaves delicately under reduced products of Banach algebras, and in particular that
the norms of witnesses to Dedekind infiniteness cannot in general be bounded uniformly
along families~\cite{DawsHorvathCJM}.

For complexity purposes it is useful to isolate bounded
\textit{strict} witness fragments.  For $M\in\N$ let $\mathrm{Bad}_M$ be the set of
$K\in\Ideal(A)$ such that, in $B=A/K$, there exist $a,b\in B$ with
\[
        \|a\|<M,\qquad \|b\|<M,\qquad ab=1,\qquad ba\ne 1 .
\]
Then $B$ is Dedekind infinite iff $K\in\bigcup_M\mathrm{Bad}_M$, because any pair of
bounded witnesses has norm strictly smaller than some integer~$M$.

\begin{proposition}\label{prop:df-banach}
The class of Dedekind finite quotients $A/K$ is closed in the Wijsman topology.
Equivalently, the class of Dedekind infinite quotients is open.
\end{proposition}

\begin{proof}
Write $B=A/K$.  For $M\in\N$ say that $B$ has an \emph{open $M$-bounded
Dedekind-infinite witness} if there are $a,b\in B$ with
\[
        \|a\|<M,
        \qquad \|b\|<M,
        \qquad ab=1,
        \qquad ba\ne 1 .
\]
A unital Banach algebra is Dedekind infinite iff it has such a witness for some
integer~$M$.

Fix $M\in\N$.  For $k,n\in\N$ with $M-1/k>0$, set
\[
\delta_{M,k,n}:=\min\Bigl\{\frac12,\frac{1}{2kM},\frac{1}{4nM^2+1}\Bigr\}.
\]
Consider the following open condition on~$K$:
\[
\begin{aligned}
&\exists u,v\in D^A:\quad
\Phi(K,u)<M-\frac1k,\qquad
\Phi(K,v)<M-\frac1k,\\
&\hspace{35mm}
\Phi(K,uv-1)<\delta_{M,k,n},\qquad
\Phi(K,vu-1)>\frac1n .
\end{aligned}
\]
All inequalities are strict and involve Wijsman-continuous functions, so the condition
is open.

We claim that $B$ has an open $M$-bounded Dedekind-infinite witness iff the displayed
condition holds for some $k,n$.  If $a,b$ are exact witnesses with norms $<M$, choose
$k$ such that both norms are $<M-1/k$ and choose $n$ with $\|ba-1\|>3/n$; density of
$D^A+K$ then gives $u,v\in D^A$ satisfying the displayed strict inequalities.
Conversely, suppose the displayed condition holds and put $u_0=u+K$, $v_0=v+K$.
Since $\|u_0v_0-1\|<\delta_{M,k,n}<1$, the element $u_0v_0$ is invertible.  Set
\[
        \widetilde a=u_0,
        \qquad
        \widetilde b=v_0(u_0v_0)^{-1}.
\]
Then $\widetilde a\widetilde b=1$ and, by the choice of $\delta_{M,k,n}$,
\[
        \|\widetilde b\|
        \le (M-1/k)(1-\delta_{M,k,n})^{-1}<M .
\]
Moreover
\[
\|\widetilde b\widetilde a-1\|
\ge \|v_0u_0-1\|-
     \|v_0((u_0v_0)^{-1}-1)u_0\|
> \frac1n- M^2\frac{\delta_{M,k,n}}{1-\delta_{M,k,n}}>0.
\]
Thus $\widetilde a,\widetilde b$ are genuine open $M$-bounded witnesses.

The Dedekind-infinite class is the countable union of the open sets just described,
over $M,k,n$ and $u,v\in D^A$.  Hence it is open, and its complement is closed.
\end{proof}

\begin{remark}\label{rem:df-banach-nontrivial}
The Daws--Horv\'ath analysis~\cite{DawsHorvathCJM} shows that norms of witnesses to
Dedekind infiniteness cannot in general be bounded uniformly along arbitrary families.
The preceding proof avoids any uniform global bound: it writes Dedekind infiniteness as
a countable union of open bounded fragments, using open balls and a Neumann-series
correction inside each fragment.
\end{remark}

%%%%%%%%%%%%%%%%%%%%%%%%%%%%%%%%%%%%%%%%%%%%%%%%%%%%%%%%%%%%%%
\section{\texorpdfstring{$C^*$}{C*}-algebras: decomposition, approximation, and regularity}
\label{sec:Cstar-properties}
%%%%%%%%%%%%%%%%%%%%%%%%%%%%%%%%%%%%%%%%%%%%%%%%%%%%%%%%%%%%%%

Throughout, let $A=C^*_{\max}(F_\infty)$ and fix a countable dense $*$-subalgebra
$D^A$ closed under multiplication and~$*$.  Work on $\Ideal(A)$ with the Wijsman
topology and write $B=A/K$.

\begin{remark}\label{rem:Cstar-unital-scope}
The parameter space in Sections~\ref{sec:Cstar-properties}--\ref{sec:D-stability}
is based on the \emph{unital} quotient generator $A=C^*_{\max}(F_\infty)$.
Accordingly, the statements proved there are literally statements about quotients of
this unital generator, hence about unital separable $C^*$-algebras (aside from the
zero quotient).  The standard non-unital variant is obtained by replacing $A$ with
$A\otimes\mathcal K$; predicates involving $1$ should then be rewritten in the
unitisation, or equivalently using a fixed strictly positive contraction.  Since this
modification is routine but notationally cumbersome, we keep the body of the paper in
the unital coding.
\end{remark}

\begin{remark}
\label{rem:other-codings}
Fix a separable infinite-dimensional Hilbert space $H$.
The first standard Borel parametrisation of separable $C^*$-algebras was introduced
by Kechris~\cite{KechrisAsianCstar}.  Farah--Toms--T\"ornquist later introduced and
used several mutually Borel-equivalent standard Borel/Polish parameter spaces for
separable $C^*$-algebras; we briefly recall two of them.

\smallskip
\noindent{\bf (1) The concrete coding $\Gamma(H)$.}
Let $\Gamma(H):=B(H)^{\N}$ equipped with the product of the strong operator topology.
Each $\gamma=(\gamma_n)_{n\in\N}\in\Gamma(H)$ codes the separable $C^*$-algebra
$C^*(\gamma)\subseteq B(H)$ generated by $\{\gamma_n:n\in\N\}$.

\smallskip
\noindent{\bf (2) The seminorm coding $\Xi$.}
Fix an enumeration $(p_n)_{n\in\N}$ of all formal $*$-polynomials over $\Q(i)$
in countably many non-commuting variables $X_0,X_1,\dots$.
The space $\Xi$ consists of all functions
$\xi:\N\to\R_{\ge 0}$ that arise as
$\xi(n)=\|p_n(a_0,a_1,\dots)\|$
for some separable $C^*$-algebra $B$ and some sequence $(a_k)$ in $B$.
Equivalently, $\xi$ encodes a $C^*$-seminorm on the free $*$-algebra;
the completion of the corresponding quotient is the coded algebra~$B_\xi$.

\smallskip
Farah--Toms--T\"ornquist show that $\Gamma(H)$ and $\Xi$ (as well as several related
codings, including abstract countable-structure codings $\widehat\Gamma$ and
$\widehat\Xi$) are equivalent parameterisations: there are Borel maps between the
coding spaces sending a code to a code for an isomorphic algebra; see
\cite[Section~2]{FarahTomsTornquistCrelle} and~\cite{FarahTomsTornquist}.
These are also the codings used in
\cite{ElliottFarahPaulsenRosendalTomsTornquist}.

\smallskip\noindent
\textbf{How we transfer Borelness.}
The existence of a Borel isomorphism between two bare standard Borel spaces is not by
itself enough to transfer a class of algebras: the map must send a code to a code for
an isomorphic algebra.  We therefore use only explicit Borel maps with this
isomorphism-preserving property.  Borel Schr\"oder--Bernstein may be used to identify
the underlying standard Borel spaces once injections in both directions are known, but
none of the rank or Borelness arguments below relies on an uncontrolled
Schr\"oder--Bernstein bijection.
\end{remark}

\begin{lemma}
\label{lem:pullback-ideals}
Let $\sigma:C\twoheadrightarrow D$ be a surjective $*$-homomorphism between separable
$C^*$-algebras.  Then the pullback map
\[
\sigma^*:\Ideal(D)\to\Ideal(C),\qquad J\mapsto\sigma^{-1}(J),
\]
is a Wijsman-continuous injection.
\end{lemma}

\begin{proof}
Injectivity: if $J\neq L$ then $\sigma^{-1}(J)\neq\sigma^{-1}(L)$ because $\sigma$ is
surjective.  For Wijsman continuity, fix $c\in C$.  Because $\sigma$ is a surjective
$*$-homomorphism, the quotient $C/\sigma^{-1}(J)\cong D/J$ isometrically (via the
canonical isomorphism), and therefore
\[
\dist\bigl(c,\,\sigma^{-1}(J)\bigr)
=\|\sigma(c)+J\|_{D/J}
=\dist\bigl(\sigma(c),\,J\bigr).
\]
Since $J\mapsto\dist(\sigma(c),J)$ is continuous in the Wijsman topology on $\Ideal(D)$,
it follows that $J\mapsto\dist(c,\sigma^{-1}(J))$ is continuous on $\Ideal(C)$.
As $c$ was arbitrary, $\sigma^*$ is Wijsman-continuous.
\end{proof}

\begin{remark}
\label{rem:Xi-ideals}
The seminorm coding $\Xi$ is naturally identified with an ideal space.  Let
\[
U:=C^*\bigl(x_n\ (n\in\N)\ :\ \|x_n\|\le 1\bigr)
\]
be the universal unital $C^*$-algebra generated by countably many contractions.  Then $U$
is separable, and there is a canonical Borel bijection $\Xi\cong\Ideal(U)$: a seminorm
$\xi\in\Xi$ corresponds to the kernel $I_\xi$ of the canonical surjection $U\to B_\xi$,
and conversely $I\in\Ideal(U)$ yields the seminorm $\xi_I(p)=\|p+I\|_{U/I}$, which is
continuous coordinatewise in~$I$.

\smallskip
\noindent{\bf (3) A continuous injection $\Ideal(A)\hookrightarrow\Xi$.}
Let $(g_m)_{m\in\N}$ denote the canonical free unitary generators of
$A=C^*_{\max}(F_\infty)$.  Since each $g_m$ is a contraction, the universal property of

\[
U:=C^*\bigl(x_m\ (m\in\N): \|x_m\|\le 1\bigr)
\]

gives a surjective $*$-homomorphism
$\sigma:U\twoheadrightarrow A$ with $\sigma(x_m)=g_m$.  By
Lemma~\ref{lem:pullback-ideals},

\[
        K\longmapsto \sigma^{-1}(K)
\]

is a Wijsman-continuous injection from $\Ideal(A)$ into $\Ideal(U)\cong\Xi$, and the
$\Xi$-code obtained in this way represents the quotient $A/K$.

\smallskip
\noindent{\bf (4) A continuous injection $\Xi\hookrightarrow\Ideal(A)$.}
Because $A=C^*_{\max}(F_\infty)$ is a quotient generator, $U$ (being separable) is a
quotient of~$A$.  Fix once and for all a surjective $*$-homomorphism
$\pi:A\twoheadrightarrow U$.
By Lemma~\ref{lem:pullback-ideals}, the pullback $I\mapsto\pi^{-1}(I)$ is a
Wijsman-continuous injection $\Ideal(U)\hookrightarrow\Ideal(A)$.
Under the identification $\Xi\cong\Ideal(U)$, this gives a continuous injection
$\Xi\hookrightarrow\Ideal(A)$.  The Borel Schr\"oder--Bernstein theorem therefore gives
an abstract Borel isomorphism between the two underlying standard Borel spaces.  We shall
not use that abstract bijection for transferring algebraic properties; for that purpose
one must use the explicit quotient-preserving maps just described.

\smallskip
\noindent{\bf (5) Borel compatibility with $\Gamma(H)$.}
Composing the continuous injections of~(3) and~(4) with the Borel equivalences
between $\Xi$ and $\Gamma(H)$ from
\cite{FarahTomsTornquistCrelle}, we obtain Borel maps $\Gamma(H)\to\Ideal(A)$ and
$\Ideal(A)\to\Gamma(H)$ which send a code to a code for an isomorphic $C^*$-algebra.
In particular, any Borel-complexity statement for a $C^*$-algebraic property in our
ideal-quotient coding transfers to the equivalent statement in the FTT codings, and
vice versa.
\end{remark}

\begin{lemma}\label{lem:fixed-ideal-to-Xi}
Let $C$ be a fixed separable $C^*$-algebra and let $(c_j)_{j\in\N}$ be a dense
sequence in the unit ball generating~$C$ as a $C^*$-algebra.  There is a
Wijsman-continuous map from $\Ideal(C)$ to the appropriate standard seminorm coding
$\Xi$ which sends $I\in\Ideal(C)$ to a code for $C/I$.
\end{lemma}

\begin{proof}
If $C$ is non-unital, use the standard non-unital seminorm coding, obtained by
enumerating rational $*$-polynomials without a constant term; this coding is Borel
equivalent to the concrete coding $\Gamma(H)$ used by Farah--Toms--T\"ornquist.  If
$C$ is unital, use the unital version described above.  In either case, for the fixed
enumeration $(p_n)$ of the relevant rational $*$-polynomials, define
\[
        \xi_I(n):=\|p_n(c_0,c_1,\ldots)+I\|_{C/I}.
\]
For each $n$, the element $p_n(c_0,c_1,\ldots)\in C$ is fixed, and the coordinate
$I\mapsto\xi_I(n)$ is continuous in the Wijsman topology.  The resulting seminorm is
precisely the quotient seminorm of the dense $*$-subalgebra generated by the images of
the $c_j$, so the completion coded by $\xi_I$ is $C/I$.
\end{proof}

\paragraph{Nuclearity.}
Because nuclearity is invariant under $*$-isomorphism and the explicit Borel maps in
Remark~\ref{rem:Xi-ideals} send quotient codes to standard codes for isomorphic
algebras, the class
\[
\{K\in\Ideal(A):A/K\text{ is nuclear}\}
\]
is Borel in the Wijsman topology.  Indeed, one route is via the CPAP
reformulation of nuclearity: Effros proved Borelness of the nuclear class in the
standard parameterisations; see Section~5 of \cite{FarahTomsTornquistCrelle}.
This gives a quick external route to Borelness in the present coding.

It seems plausible that one can sharpen this to a low finite Borel rank by combining the
CPAP with an internal coding lemma for completely positive contractive maps from
finite-dimensional $C^*$-algebras into quotients $A/K$.  Since that lemma is not
developed in the present paper, we record only the Borel upper bound here.

\begin{remark}
\label{rem:ideal-coding-history}
The ideal-space coding of separable $C^*$-algebras as quotients of a fixed universal
separable $C^*$-algebra should not be presented as historically new.  Nate Brown
suggested this quotient/ideal coding around 2010, as communicated to the author by
Ilijas Farah.  Related independent forthcoming work of Austin Shiner, using quotient
encodings to study functorial complexity for $C^*$-algebras, is noted at the end of
the Introduction.
\end{remark}


\paragraph{Continuous logic.}
It is also natural to compare the present quotient/Wijsman formalism with the
continuous-logic treatment of axiomatizable and definable classes of $C^*$-algebras;
see \cite{FarahHartLupiniRobertTikuisisVignatiWinter}.  After transferring between
our ideal coding and the standard FTT codings, this gives an alternative abstract
route to Borelness for many axiomatizable classes.  We nevertheless retain the
direct arguments here because they keep track of concrete Borel ranks inside the
present parameter space.

\subsection{Coding projections, unitaries, and matrix units}

\begin{lemma}\label{lem:basic-coding}
For $x\in D^A$ and $K\in\Ideal(A)$:
\begin{enumerate}
\item $x+K$ is a projection in $B$ iff $\Phi(K,x^2-x)=0$ and $\Phi(K,x^*-x)=0$.
\item $u+K$ is unitary in $B$ iff $\Phi(K,u^*u-1)=0$ and $\Phi(K,uu^*-1)=0$.
\end{enumerate}
For fixed finite families, the predicate ``forms a system of matrix units'' is closed.
\end{lemma}

\begin{proof}
All assertions follow directly from identities in $B$ expressed by $\Phi$ and
the continuity of $\Phi$ in~$K$.
\end{proof}

\begin{lemma}\label{lem:positivity-coding}
Let $B=A/K$.
\begin{enumerate}
\item An element $a+K\in B$ is positive iff for every $m\in\N$ there exists
$y\in D^A$ with $\Phi(K,a-y^*y)<1/m$.  In particular, the set of $K$ for which a
fixed $a\in D^A$ has positive image is~$\Pi^0_2$.
\item Every positive element of $B$ is approximated by elements of the form $b^*b+K$
with $b\in D^A$.  Hence in quantifications
``$\forall$ positive $a\in B_+$'', it suffices to range $a$ over $\{b^*b:b\in D^A\}$,
which is countable.
\item For $a\in D^A$ and a projection $p\in B$ (coded as in
Lemma~\ref{lem:basic-coding}), set $x=p(a^*a+K)p\in pBp$.
Then
\[
x\ge \tfrac{1}{n}p
\quad\Longleftrightarrow\quad
x \text{ is invertible in }pBp \text{ and }\|x^{-1}\|\le n.
\]
Moreover, if there exists $y\in pBp$ with $\|y\|\le n$ and
\[
\|xy-p\|<\tfrac12,\qquad \|yx-p\|<\tfrac12,
\]
then $x$ is invertible in $pBp$ and $\|x^{-1}\|\le 2n$; hence
$x\ge \tfrac{1}{2n}p$.
Thus, for the descriptive-set-theoretic arguments below, the relation
$x\ge \tfrac1n p$ may be replaced by this bounded stable witness scheme,
at the expense of renumbering $n$ by a fixed multiplicative constant.
\end{enumerate}
\end{lemma}

\begin{proof}
(i) If $a+K\ge 0$ in $B$, let $c$ be a positive square root, and approximate $c$ by
$y_m\in D^A$ to get $\|a+K-y_m^*y_m+K\|<1/m$.  Conversely, if $a+K=\lim(y_m^*y_m+K)$,
then $a+K\ge 0$ since the positive cone is closed.

(ii) follows from (i) and the fact that $D^A$ is closed under~$*$ and multiplication.

(iii) In the unital $C^*$-algebra $pBp$ (with unit $p$), a positive element $x\ge 0$
satisfies $x\ge \tfrac{1}{n}p$ iff $x$ is invertible with $\|x^{-1}\|\le n$
(this is a standard spectral fact: $x\ge\tfrac{1}{n}p$ iff
$\mathrm{sp}(x)\subseteq [1/n,\infty)$, iff $x^{-1}$ exists with
$\|x^{-1}\|=1/\min(\mathrm{sp}(x))\le n$).
If there exists $y\in pBp$ with $\|y\|\le n$ and $\|xy-p\|<1/2$, then the Neumann
series gives $\|x^{-1}\|\le \|y\|\,\|(xy)^{-1}\|\le n\cdot 2=2n$, so $x\ge\frac{1}{2n}p$.
The constant $2$ is harmless for our purposes: replacing $n$ by $2n$ does not change
the quantifier structure.
Approximating $p$ and $y$ by elements of $D^A$ and using stability of
projections/invertibles (Lemma~\ref{lem:stable-fd} and
Remark~\ref{rem:stable-relations}), this is detected by strict $\Phi$-inequalities
on bounded balls.
\end{proof}

\begin{lemma}\label{lem:stable-fd}
For each $d\in\N$ and $\varepsilon>0$ there exists $\delta>0$ such that the following
holds in every $C^*$-algebra~$B$.

If $(x_{ij})_{1\le i,j\le d}\subseteq B$ satisfies
\[
\begin{aligned}
\max_{i,j}\|x_{ij}^*-x_{ji}\| &<\delta,\\
\max_{i,j,k,\ell}\|x_{ij}x_{k\ell}-\delta_{jk}x_{i\ell}\| &<\delta,\\
\Bigl\|\sum_{i=1}^d x_{ii}-1\Bigr\| &<\delta,
\end{aligned}
\]
then there exists a system of matrix units $(e_{ij})_{i,j\le d}\subseteq B$ such that
$\max_{i,j}\|e_{ij}-x_{ij}\|<\varepsilon$.
In particular, the $C^*$-subalgebra generated by $(e_{ij})$ is (unitally) isomorphic
to~$M_d$.
\end{lemma}

\begin{proof}
This is a standard perturbation/semiprojectivity fact for finite-dimensional
$C^*$-algebras; see~\cite{Loring}.
\end{proof}

\begin{remark}\label{rem:stable-relations}
Several arguments below (for AF, QD, MF, and related approximation properties, as well
as for the $K_0$ coding in Section~\ref{sec:K0}) involve quantification over the
countable dense $*$-subalgebra~$D^A$ using exact equalities such as $\Phi(K,\cdot)=0$
to assert that certain quotient elements are projections, matrix units, or unitaries.
Such exact witnesses need not belong to~$D^A$.

The resolution is standard: one replaces exact relations by \emph{approximate} ones
(expressed by strict inequalities $\Phi(K,\cdot)<\delta$ on elements of $D^A$) and
appeals to perturbation stability (Lemma~\ref{lem:stable-fd} and its generalisations)
to promote approximate witnesses to genuine ones in the quotient.
Since strict inequalities are open predicates, the quantifier-complexity analysis is
unchanged.
\end{remark}

\begin{lemma}
\label{lem:ucp-fd}
Fix $m\in\N$.
\begin{enumerate}
\item \textbf{(UCP maps with prescribed values.)}
Let $\bar x=(x_1,\dots,x_r)\subseteq D^A$ and let $Y_1, \dots, Y_r \in M_m(\Q(i))$.
The condition that there exists a unital completely positive (ucp) map
$\varphi:A/K\to M_m$ satisfying $\varphi(x_j+K)=Y_j$ for all $1\le j\le r$
is a \emph{closed} (and in particular Borel) predicate in $K \in \Ideal(A)$.

\item \textbf{(Approximate embeddings of f.d.\ $C^*$-algebras.)}
A family $(e_{ij})_{i,j\le d}\subseteq D^A$ whose images in $B$
\emph{approximately} satisfy matrix-unit relations (within tolerance~$\delta$) yields,
by Lemma~\ref{lem:stable-fd}, a genuine unital copy of $M_d$ in~$B$.
More precisely, the predicate ``$\exists$ a unital $*$-homo\-morphism
$\pi:F_0\to B$ with $F_0\cong\bigoplus_{j=1}^\ell M_{d_j}$ and
$\max_s\dist(x_s+K,\pi(F_0))<\varepsilon$'' is Borel in~$K$:
quantify over all choices of approximate matrix-unit families for each direct summand
$M_{d_j}$ with entries in $D^A$ (a countable search), require approximate matrix-unit
relations to hold within tolerance~$\delta$ via strict $\Phi$-inequalities, and require
$\varepsilon$-closeness of the target elements.  This contributes one existential
quantifier block.
\end{enumerate}
\end{lemma}

\begin{proof}
Part~(i): The existence of a ucp map $\varphi:A/K\to M_m$ with
$\varphi(x_j+K)=Y_j$ is equivalent to the existence of a ucp map
$\psi:A\to M_m$ with $\psi(x_j)=Y_j$ for $1\le j\le r$ and $\psi$ vanishing on~$K$.
Since any ucp map is contractive, vanishing on $K$ is equivalent to
$\|\psi(a)\| \le \|a+K\|_{A/K} = \Phi(K,a)$ for all $a\in A$; by density
and continuity it suffices to check this inequality for $a\in D^A$.

Let $\mathrm{UCP}(A,M_m)$ denote the space of all ucp maps $A\to M_m$, equipped
with the point-norm topology.
Since $M_m$ is finite-dimensional, the point-norm topology on $\mathrm{UCP}(A,M_m)$
coincides with pointwise convergence of the $m^2$ matrix-entry functionals in the
weak$^*$ topology of $A^*$; hence $\mathrm{UCP}(A,M_m)$ is compact metrisable by
Banach--Alaoglu and separability of $A$.
Consider the constraint on pairs
$(\psi, K) \in \mathrm{UCP}(A,M_m) \times \Ideal(A)$:
\[
\forall j\le r\ \bigl(\psi(x_j)=Y_j\bigr) \ \wedge \ 
\forall a\in D^A\ \bigl(\|\psi(a)\| \le \Phi(K,a)\bigr).
\]
For each fixed $a\in D^A$, the evaluation $\psi \mapsto \|\psi(a)\|$ is continuous,
and $K \mapsto \Phi(K,a)$ is continuous (Lemma~\ref{lem:Phi-continuous}), so
$\|\psi(a)\|\le\Phi(K,a)$ defines a closed subset of the product.
Similarly, each condition $\psi(x_j)=Y_j$ is closed.
A countable intersection of closed sets is closed, so the full constraint set
is closed in $\mathrm{UCP}(A,M_m) \times \Ideal(A)$.
Because $\mathrm{UCP}(A,M_m)$ is compact, the projection of this closed set onto
$\Ideal(A)$ is closed.

Part~(ii): approximate
matrix-unit families on a fixed finite set are quantified over
$(x_{ij})\subseteq D^A$, with the matrix-unit relations required to hold up to a
tolerance~$\delta$ in $B$ (written using strict inequalities in~$\Phi$).
Lemma~\ref{lem:stable-fd} then yields genuine matrix units in~$B$
nearby, hence an actual unital copy of the prescribed finite-dimensional algebra.
Closeness to a finite set is measured by~$\Phi$.
\end{proof}

\begin{lemma}\label{lem:cpc-fd}
Let $F_0$ be a finite-dimensional $C^*$-algebra, let
$\bar x=(x_1,\dots,x_r)\subseteq D^A$, and let $\bar y=(y_1,\dots,y_r)\subseteq F_0$.
Then the set of $K\in\Ideal(A)$ for which there exists a completely positive
contractive map
$\psi:A/K\to F_0$
satisfying $\psi(x_j+K)=y_j$ for all $1\le j\le r$ is closed.
\end{lemma}

\begin{proof}
Let $\mathrm{CPC}(A,F_0)$ denote the space of all completely positive contractive
maps $A\to F_0$, equipped with the point-norm topology.
Since $F_0$ is finite-dimensional, $\mathrm{CPC}(A,F_0)$ is compact metrisable:
using the dense sequence $D^A=\{u_m:m\in\N\}$, each
$\psi\in\mathrm{CPC}(A,F_0)$ is determined by the tuple
$(\psi(u_m))_{m\in\N}\in \prod_m B_{\|u_m\|}^{F_0}$, a compact metrisable product,
and the conditions of linearity, $*$-preservation, complete positivity, and
contractivity define a closed subset.

A c.p.c.\ map $\psi:A\to F_0$ factors through $A/K$ iff
\[
\|\psi(a)\| \le \|a+K\|_{A/K}=\Phi(K,a)\qquad(a\in A),
\]
and by density it suffices to check this on $a\in D^A$.
Hence the desired relation on pairs $(\psi,K)\in\mathrm{CPC}(A,F_0)\times\Ideal(A)$ is
\[
\forall j\le r\ \bigl(\psi(x_j)=y_j\bigr)
\quad\wedge\quad
\forall a\in D^A\ \bigl(\|\psi(a)\|\le \Phi(K,a)\bigr).
\]
Each condition $\psi(x_j)=y_j$ is closed, and for fixed $a\in D^A$ the inequality
$\|\psi(a)\|\le \Phi(K,a)$ is closed because $(\psi,K)\mapsto \|\psi(a)\|$ and
$(\psi,K)\mapsto\Phi(K,a)$ are continuous.
Thus the full relation is closed in
$\mathrm{CPC}(A,F_0)\times\Ideal(A)$.
Projecting this closed set along the compact factor $\mathrm{CPC}(A,F_0)$ yields a
closed subset of $\Ideal(A)$.
\end{proof}

\subsection{AF and related approximation classes}

A $C^*$-algebra $B$ is \emph{approximately finite-dimensional} (\emph{AF}) if it is the
closure of an increasing union of finite-dimensional $C^*$-subalgebras, or
equivalently, if for every finite $F\subseteq B$ and every $\varepsilon>0$ there exists
a finite-dimensional $C^*$-subalgebra $F_0\subseteq B$ approximating $F$
within~$\varepsilon$ (see \cite[II.8.2.2]{BlackadarOA}).

\begin{proposition}\label{prop:AF}
The class $\mathrm{AF}=\{K\in\Ideal(A):A/K\text{ is AF}\}$ is $\Pi^0_2$ in the Wijsman
topology.
\end{proposition}

\begin{proof}
Recall that $B$ is AF iff for every finite set $F\subseteq B$ and every
$\varepsilon>0$ there exists a finite-dimensional $C^*$-subalgebra approximating $F$
within~$\varepsilon$.  It suffices to test finite tuples $\bar x\in(D^A)^{<\N}$ and
tolerances $1/n$.

Fix such a pair $(\bar x,n)$.  A witness consists of a finite-dimensional algebra type
$F_0\cong\bigoplus_{j=1}^\ell M_{d_j}$, approximate matrix units from $D^A$ for this
algebra, and rational linear combinations of those approximate matrix units
approximating the entries of $\bar x$.  Requiring the matrix-unit relations to hold
within the stability tolerance from Lemma~\ref{lem:stable-fd}, and requiring the
approximations to be within $1/n$, is a finite system of strict inequalities in
Wijsman-continuous functions of~$K$.  For fixed finite data this is open, and the search
over all finite-dimensional types, approximate matrix-unit tuples, and rational
coefficients is countable.  Thus the set of $K$ admitting an AF witness for the fixed
pair $(\bar x,n)$ is open.

Intersecting these open sets over all finite tuples $\bar x$ and all $n\in\N$ gives a
$G_\delta$ set, hence a $\Pi^0_2$ set.
\end{proof}

Recall that a $C^*$-algebra is called \emph{AI} (approximately interval) if it is the
closure of an increasing union of subalgebras each isomorphic to a direct sum of
algebras of the form $C([0,1],M_n)$; it is \emph{AH} (approximately homogeneous)
if the building blocks are homogeneous $C^*$-algebras $p\,C(X,M_n)\,p$ with $X$ a
compact metrizable space; it is \emph{ASH} (approximately sub\-homogeneous) if the
building blocks are subhomogeneous; and it is \emph{A$\mathbb T$} (approximately circle)
if the building blocks are direct sums of $C(\mathbb T,M_n)$ and finite-dimensional
algebras.  For details see \cite[IV.3]{BlackadarOA}.

\begin{proposition}
\label{prop:AI-AH}
Let $\mathcal C=(C_j)_{j\in\N}$ be a countable family of separable $C^*$-algebras,
and assume that each $C_j$ is given by a stable finite presentation.  Let
$\mathrm{LocIm}(\mathcal C)$ be the class of $K\in\Ideal(A)$ such that $B=A/K$ has the
following property: for every finite $F\subseteq B$ and every $\varepsilon>0$ there are
$j\in\N$ and a $*$-homomorphism $\pi:C_j\to B$ such that
$\dist(x,\pi(C_j))<\varepsilon$ for all $x\in F$.

Then $\mathrm{LocIm}(\mathcal C)$ is $\Pi^0_2$ in the Wijsman topology.
\end{proposition}

\begin{proof}
Fix a finite tuple $\bar x\in(D^A)^{<\N}$ and $n\in\N$.  For each $j$, fix a stable
finite presentation of $C_j$ on generators $z^{(j)}_1,\dots,z^{(j)}_{r_j}$, and choose a
finite $1/(2n)$-net in a sufficiently large ball of the rational $*$-algebra generated
by these symbols.  Stability supplies a tolerance $\delta_j>0$ such that any tuple in a
quotient satisfying the defining relations within $\delta_j$ is within $1/(2n)$ of the
range of an actual representation of $C_j$.

For fixed $j$ and a tuple $\bar y\in(D^A)^{r_j}$, require the defining relations of
$C_j$ to hold for $\bar y+K$ within $\delta_j$, and require every element of $\bar x$ to
be within $1/n$ of the image, under the formal evaluation at $\bar y$, of one of the
chosen net elements.  These are finitely many strict inequalities involving
Wijsman-continuous quotient-norm functions, hence they define an open subset of
$\Ideal(A)$.

Taking the countable union over $j$ and over $\bar y$ gives, for the fixed pair
$(\bar x,n)$, an open set of codes admitting a local witness.  Intersecting over all
finite tuples $\bar x$ and all $n$ gives a $G_\delta$ set, that is, a $\Pi^0_2$ set.
\end{proof}

\begin{remark}
The proposition is deliberately formulated for homomorphic images.  Stable finite
presentations let one perturb approximate relations to genuine representations, but
they do not by themselves force the representation to be injective.  Therefore the
argument above does not prove that the usual AI, $A\mathbb T$, AH, or ASH classes are
$\Pi^0_2$ or $\Pi^0_3$ in this coding.  To recover those embedding versions one would
need additional Borel conditions controlling the spectrum or injectivity of the
represented block.
\end{remark}

A separable $C^*$-algebra $B$ is \emph{quasidiagonal} (QD) if every faithful
representation of $B$ on a separable Hilbert space admits an increasing approximate
unit of finite-rank projections that asymptotically commute with every operator in
the image, or equivalently if $B$ admits Voiculescu-type approximate embeddings into
matrix algebras (see \cite[V.4.2.1]{BlackadarOA}).
It is \emph{MF} (matricially finite) if it embeds into
$\prod_n M_{k_n}/(\bigoplus_n M_{k_n})_{c_0}$
for some sequence $(k_n)$
(see \cite[V.4.3.1]{BlackadarOA}).

\begin{proposition}\label{prop:QD-MF}
The class of quasidiagonal quotients is $\Pi^0_3$ in the Wijsman topology, and the
class of MF quotients is $\Pi^0_2$.
\end{proposition}

\begin{proof}
Write $B=A/K$.

\smallskip
\noindent\textbf{(1) Quasidiagonality.}
We use the standard finite-dimensional approximation formulation of
quasidiagonality for separable unital $C^*$-algebras:
$B$ is quasidiagonal iff for every finite set $F\subseteq B$ and every
$\varepsilon>0$ there exist $m\in\N$ and a ucp map $\varphi:B\to M_m$ such that
\begin{equation}\label{eq:QD-approx}
\|\varphi(xy)-\varphi(x)\varphi(y)\|<\varepsilon\quad (x,y\in F),
\qquad
\bigl|\|\varphi(x)\|-\|x\|\bigr|<\varepsilon\quad (x\in F).
\end{equation}
(See e.g.\ \cite[V.4.2.1]{BlackadarOA} for such Voiculescu-type formulations.)

As usual, it suffices to test finite sets drawn from $D^A+K$ and
$\varepsilon=1/n$.
Fix a finite tuple $\bar x=(x_1,\dots,x_r)\subseteq D^A$ and $n\in\N$.
Let $F(\bar x)$ be the finite $*$-closed set consisting of
\[
1,\ x_i,\ x_i^*,\ x_ix_j,\ x_i^*x_j,\ x_ix_j^*,\ x_i^*x_j^*
\qquad (1\le i,j\le r).
\]

For $m\in\N$, let $\mathrm{UCP}(A,M_m)$ be the compact metrisable space of ucp
maps $A\to M_m$.
Consider the set of pairs $(\varphi,K)\in \mathrm{UCP}(A,M_m)\times\Ideal(A)$ such that

\begin{enumerate}[label=\textup{(\alph*)}]
\item $\varphi$ descends to the quotient $A/K$, \emph{i.e.}
$\|\varphi(a)\|\le \Phi(K,a)$ for all $a\in D^A$;
\item the induced map on $B=A/K$ is $1/n$-approximately multiplicative on
$F(\bar x)$:
$\|\varphi(st)-\varphi(s)\varphi(t)\|<1/n$ for all $s,t,st\in F(\bar x)$;
\item the induced map is $1/n$-approximately isometric on $F(\bar x)$:
$\bigl|\|\varphi(s)\|-\Phi(K,s)\bigr|<1/n$ for all $s\in F(\bar x)$.
\end{enumerate}

Condition~(a) is closed in $(\varphi,K)$, by the same compactness argument as in
Lemma~\ref{lem:ucp-fd}(i).
Conditions~(b) and~(c) are open, since they are strict inequalities involving
continuous functions of $(\varphi,K)$.
Hence the total relation is locally closed in the compact-by-Polish space
$\mathrm{UCP}(A,M_m)\times\Ideal(A)$.
By Lemma~\ref{lem:proj-locally-closed}, its projection onto $\Ideal(A)$ is
$F_\sigma$.

Taking the union over $m\in\N$, we see that for fixed $(\bar x,n)$ the set of
$K$ for which there exists a quasidiagonal witness is $\Sigma^0_2$.
Intersecting over all finite tuples $\bar x$ and all $n\in\N$ yields that the
quasidiagonal class is $\Pi^0_3$.

\smallskip
\noindent\textbf{(2) MF.}
We use the standard microstate/norm-ultraproduct characterisation:
a separable unital $C^*$-algebra $B$ is MF iff for every finite set
$F\subseteq B$ and every $\varepsilon>0$ there exist $m\in\N$ and a map
$\theta:F\to M_m$ such that $\theta$ is approximately $*$-multiplicative on $F$
and approximately isometric on $F$.

As in the QD case, it suffices to test finite sets drawn from $D^A+K$ and
$\varepsilon=1/n$.
Fix $\bar x=(x_1,\dots,x_r)\subseteq D^A$ and let $F(\bar x)$ be the same finite
$*$-closed set as above.
A candidate MF witness is a matrix size $m$ and a family of rational matrices
$(Y_t)_{t\in F(\bar x)}\subseteq M_m(\Q(i))$
satisfying
\begin{align*}
\|Y_{st}-Y_sY_t\|&<\tfrac{1}{n}\qquad (s,t,st\in F(\bar x)),\\
\|Y_{s^*}-Y_s^*\|&<\tfrac{1}{n}\qquad (s,s^*\in F(\bar x)),\\
\bigl|\|Y_s\|-\Phi(K,s)\bigr|&<\tfrac{1}{n}\qquad (s\in F(\bar x)).
\end{align*}
For fixed $(m,(Y_t))$ these are strict inequalities in continuous functions of $K$,
hence define an open set of $K$.
Taking the union over the countable parameter set of such matrix families shows
that, for each fixed $(\bar x,n)$, the set of $K$ admitting an MF witness is open.
Intersecting over all $\bar x$ and $n$ yields that the MF class is a $G_\delta$ set,
that is, $\Pi^0_2$.
\end{proof}

\begin{proposition}\label{prop:approx-div}
A unital $C^*$-algebra $B$ is \emph{approximately divisible} if for every finite
$F\subseteq B$ and $\varepsilon>0$ there exist $p,q\ge 2$ and a unital embedding
$M_p\oplus M_q\hookrightarrow B$ whose image $\varepsilon$-commutes with every
element of~$F$ (see \cite[V.2.3.1]{BlackadarOA}).  The class of $K\in\Ideal(A)$ such
that $A/K$ is approximately divisible is $G_\delta$ in the Wijsman topology.
\end{proposition}

\begin{proof}
For fixed $\bar x\in(D^A)^{<\N}$ and $m\in\N$, the existence of an approximately
central unital copy of some $M_p\oplus M_q$ is witnessed by approximate matrix units
from $D^A$ for the two summands, with $p,q\ge2$, satisfying the finite-dimensional
relations and the centrality inequalities $\Phi(K,[e_{ij},x_s])<1/m$ with strict
margins.  Stability of finite-dimensional relations promotes such approximate data to
a genuine unital embedding after decreasing the tolerance.  For fixed witnesses the
conditions are open, and the search over $p,q$ and dense witnesses is countable; hence
the witness set for fixed $(\bar x,m)$ is open.  Intersecting over $\bar x$ and $m$
gives a $G_\delta$ set.
\end{proof}

\subsection{Stable finiteness}

A unital $C^*$-algebra $B$ is \emph{finite} if every isometry in $B$ is unitary, and
\emph{stably finite} if $M_n(B)$ is finite for every $n\in\N$
(see \cite[V.2.1.1]{BlackadarOA}).

\begin{proposition}\label{prop:stably-finite}
The class of stably finite quotients is closed (and in particular $\Pi^0_2$)
in the Wijsman topology.
\end{proposition}

\begin{proof}
Write $B=A/K$.
By definition, $B$ is stably finite iff $M_n(B)$ is finite for every $n\in\N$, \emph{i.e.}\
iff there is no $n$ for which $M_n(B)$ contains a non-unitary isometry.

We use the matrix-level quotient norms $\Phi^{(n)}$ of
Definition~\ref{def:Phi-matrix} and the countable dense set $M_n(D^A)$
(Remark~\ref{rem:Phi-n-dense}).
Fix the concrete constants
\[
\alpha:=\tfrac14,\qquad \beta:=\tfrac12.
\]

\smallskip\noindent
\emph{Claim.} $B$ is \emph{not} stably finite iff there exist $n\in\N$ and
$v\in M_n(D^A)$ such that
\begin{equation}\label{eq:sf-witness}
\Phi^{(n)}(K,v^*v-1)<\alpha
\qquad\text{and}\qquad
\Phi^{(n)}(K,vv^*-1)>\beta.
\end{equation}

\smallskip\noindent
($\Rightarrow$) If $B$ is not stably finite, then for some $n$ there exists an
isometry $u\in M_n(B)$ which is not unitary. Then $p:=uu^*$ is a proper projection,
so $\|p-1\|=1$, \emph{i.e.}\ $\|uu^*-1\|=1$.
Choose $v\in M_n(D^A)$ with $\|v+M_n(K)-u\|<\varepsilon$, where $\varepsilon>0$ is so
small that
\begin{align*}
\|(v+M_n(K))^*(v+M_n(K)) - u^*u\|&<\alpha,\\
\|(v+M_n(K))(v+M_n(K))^* - uu^*\|&<1-\beta.
\end{align*}
Since $u^*u=1$ and $\|uu^*-1\|=1$, this yields \eqref{eq:sf-witness}.

\smallskip\noindent
($\Leftarrow$) Suppose \eqref{eq:sf-witness} holds for some $n$ and $v\in M_n(D^A)$.
Let $x:=v+M_n(K)\in M_n(B)$ and set $s:=x^*x$.
Then $\|s-1\|<\alpha<1$, so $s$ is invertible with
$\|s^{-1}-1\|\le \alpha/(1-\alpha)$.
Define the polar correction $w:=x\,s^{-1/2}\in M_n(B)$.
Then $w^*w=s^{-1/2}s\,s^{-1/2}=1$, so $w$ is an isometry.

We show $w$ is not unitary. Since $ww^*=xs^{-1}x^*$, we have
$ww^*-xx^*=x(s^{-1}-1)x^*$, so
\[
\|ww^*-xx^*\|\le \|x\|^2\,\|s^{-1}-1\|.
\]
Moreover $\|x\|^2=\|s\|\le 1+\alpha$, so with $\alpha=\tfrac14$,
\[
\|ww^*-xx^*\|\le (1+\tfrac14)\cdot\frac{1/4}{1-1/4}
= \tfrac54\cdot\tfrac13=\tfrac{5}{12}<\tfrac12.
\]
Since $\|xx^*-1\|=\Phi^{(n)}(K,vv^*-1)>\beta=\tfrac12$, we conclude
\[
\|ww^*-1\|
\ge \|xx^*-1\|-\|ww^*-xx^*\|
> \tfrac12-\tfrac{5}{12}=\tfrac{1}{12}>0.
\]
Since $ww^*$ is a projection (as $w$ is an isometry), $\|ww^*-1\|>0$ forces
$ww^*\neq 1$, so $w$ is a non-unitary isometry and $B$ is not stably finite.

\smallskip\noindent
\emph{Borel complexity.}
For fixed $n$ and $v$, each inequality in \eqref{eq:sf-witness} is a strict
inequality in a Wijsman-continuous function of $K$ (Lemma~\ref{lem:Phi-n-continuous}),
hence defines an open set in $\Ideal(A)$. Taking the union over the countable set of
pairs $(n,v)$ shows the non-stably-finite codes form an open set. Therefore the class
of stably finite quotients is closed (and in particular $\Pi^0_2$).
\end{proof}

\begin{remark}\label{rem:stably-finite-sharp}
Proposition~\ref{prop:stably-finite} is stronger than the generic rank bound supplied by
Theorem~\ref{thm:definability-rank}: although stable finiteness is \emph{a priori} a universal
property, the complement is witnessed by the \emph{strict} inequalities
\eqref{eq:sf-witness}, so the non-stably-finite class is open and the stably finite class is
actually closed.  In particular, this is a genuine $\Pi^0_1$ phenomenon, not merely a
$\Pi^0_2$ one.
\end{remark}

\subsection{Real rank}

The \emph{real rank} of a unital $C^*$-algebra $B$, denoted $\mathrm{RR}(B)$, is the
smallest non-negative integer $n$ such that for every $(n\!+\!1)$-tuple of self-adjoint
elements $(x_0,\dots,x_n)$ in $B$ and every $\varepsilon>0$, there exists a
self-adjoint tuple $(y_0,\dots,y_n)$ with $\max_j\|x_j-y_j\|<\varepsilon$ such that
$\sum_{j=0}^n y_j^2$ is invertible (see \cite[V.3.2.3]{BlackadarOA}).

\begin{proposition}\label{prop:real-rank}
For fixed $n\in\N$, the class
$\{K\in\Ideal(A):\mathrm{RR}(A/K)\le n\}$ is $G_\delta$ in the Wijsman topology.
\end{proposition}

\begin{proof}
The self-adjoint part of $D^A$ is countable and dense in the self-adjoint part of each
quotient.  For a fixed self-adjoint tuple
$\bar x=(x_0,\ldots,x_n)\in(D^A\cap A_{\mathrm{sa}})^{n+1}$ and $m\in\N$, let
$U_{\bar x,m}$ be the set of $K$ for which there are self-adjoint
$\bar y=(y_0,\ldots,y_n)\in(D^A\cap A_{\mathrm{sa}})^{n+1}$ and $z\in D^A$ such that
\[
\max_j\Phi(K,x_j-y_j)<\frac1m,
\qquad
\Phi\Bigl(K,z\sum_{j=0}^n y_j^2-1\Bigr)<1,
\qquad
\Phi\Bigl(K,\Bigl(\sum_{j=0}^n y_j^2\Bigr)z-1\Bigr)<1 .
\]
For fixed witnesses this is open, and hence $U_{\bar x,m}$ is open.  The two inverse
inequalities imply that $\sum_j(y_j+K)^2$ is invertible after a Neumann-series
correction; conversely, an invertible positive element has an inverse which may be
approximated by an element of $D^A$.  Therefore
\[
        \{K:\mathrm{RR}(A/K)\le n\}
        =\bigcap_{\bar x}\bigcap_m U_{\bar x,m},
\]
which is $G_\delta$.
\end{proof}

\subsection{Simplicity}

A $C^*$-algebra $B$ is \emph{simple} if it has no non-trivial closed two-sided ideals
(see \cite[II.1.1]{BlackadarOA}).

\begin{proposition}\label{prop:simple}
The class
\[
\mathrm{Simple}=\{K\in\Ideal(A):A/K\text{ is simple}\}
\]
is Borel in the Wijsman topology.
\end{proposition}

\begin{proof}
Farah--Magidor proved that the class of simple separable $C^*$-algebras is Borel;
more generally, their Corollary~8.2 applies to classes uniformly definable by omitting
types~\cite[Corollary~8.2]{FarahMagidor2018}.  By
Remarks~\ref{rem:other-codings} and~\ref{rem:Xi-ideals}, there is an explicit Borel
map from $\Ideal(A)$ to a standard code for an isomorphic quotient algebra.  Pulling
the Farah--Magidor Borel set back along this isomorphism-preserving map gives the
claim.
\end{proof}



\begin{remark}
No finite Wijsman-rank bound for simplicity is asserted here.  A tempting dense-set
argument would try to test whether each non-zero dense element generates the whole
algebra as an ideal, but this does not pass from an arbitrary non-full element to a
nearby dense one: the set of full elements is open, and its complement is closed, but a
sequence converging to a point of a closed set need not eventually remain in that set.
This is why Proposition~\ref{prop:simple} records only Borelness, obtained through the
standard codings.
\end{remark}

\subsection{Nuclear dimension}

The \emph{nuclear dimension} of a $C^*$-algebra $B$, denoted
$\dim_{\mathrm{nuc}}(B)$, is the smallest non-negative integer $n$ such that for
every finite $F\subseteq B$ and every $\varepsilon>0$ there exist a
finite-dimensional $C^*$-algebra $F_0$, a c.p.c.\ map $\psi:B\to F_0$, and c.p.c.\
maps $\varphi_0,\dots,\varphi_n:F_0\to B$ of order zero such that
$\|\sum_{i=0}^n\varphi_i\psi(x)-x\|<\varepsilon$ for all $x\in F$
(see \cite{WinterZacharias}).

\begin{lemma}
\label{lem:order-zero-coding}
Let
\[
F_0=\bigoplus_{j=1}^{\ell} M_{d_j}
\]
be a finite-dimensional $C^*$-algebra, and for each $j$ fix matrix units
$\{e_{ab}^{(j)}:1\le a,b\le d_j\}$ for the summand $M_{d_j}$.

Then there exist finitely many noncommutative $*$-polynomials
\[
r_1(\bar z),\dots,r_M(\bar z)
\]
in variables
\[
\bar z=\bigl(z_{ab}^{(j)}:1\le j\le \ell,\ 1\le a,b\le d_j\bigr),
\]
and, for every $\varepsilon>0$, a tolerance $\delta_{F_0}(\varepsilon)>0$ such that for
every $C^*$-algebra $B$ the following hold.

\begin{enumerate}
\item If $\varphi:F_0\to B$ is completely positive, contractive, and order zero, then the
tuple
\[
z_{ab}^{(j)}:=\varphi(e_{ab}^{(j)})
\]
satisfies
\[
r_s(\bar z)=0\qquad (1\le s\le M).
\]

\item If a tuple
\[
\bar y=\bigl(y_{ab}^{(j)}\bigr)\subseteq B
\]
satisfies
\[
\|r_s(\bar y)\|<\delta_{F_0}(\varepsilon)\qquad (1\le s\le M),
\]
then there exists a completely positive contractive order-zero map
\[
\varphi:F_0\to B
\]
such that
\[
\|\varphi(e_{ab}^{(j)})-y_{ab}^{(j)}\|<\varepsilon
\]
for all $j,a,b$.
\end{enumerate}
\end{lemma}

\begin{proof}
This is the standard weak stability of the finite-dimensional order-zero relations.
By \cite[Theorem~2.3]{WinterZacharias}, c.p.c.\ order-zero maps
$F_0\to B$ are in bijection with $*$-homomorphisms
\[
C_0((0,1])\otimes F_0 \longrightarrow B.
\]
For finite-dimensional $F_0$, the cone $C_0((0,1])\otimes F_0$ admits a stable finite
presentation; equivalently, the relations defining the canonical generators
$t\otimes e_{ab}^{(j)}$ are weakly stable, see
\cite[Proposition~1.4]{WinterZacharias}.  Choosing such a finite stable presentation
gives the required polynomials $r_s$ and tolerances $\delta_{F_0}(\varepsilon)$.
\end{proof}

\begin{proposition}\label{prop:nucdim}
For fixed $n\in\N$, the class
$\{K\in\Ideal(A):\dim_{\mathrm{nuc}}(A/K)\le n\}$ is $\Pi^0_3$
in the Wijsman topology.
\end{proposition}

\begin{proof}
Fix a finite tuple
\[
\bar x=(x_1,\dots,x_r)\subseteq D^A
\]
and $m\in\N$.  Let $F_{\bar x,m}$ be the set of all $K$ for which there exist a
finite-dimensional $C^*$-algebra $F_0$, a c.p.c. map
$\psi:A/K\to F_0$, and c.p.c. order-zero maps
$\varphi_0,\dots,\varphi_n:F_0\to A/K$ such that
\[
\Bigl\|\sum_{i=0}^{n}\varphi_i\psi(x_t+K)-(x_t+K)\Bigr\|<\frac1m
\qquad (1\le t\le r).
\]
We shall construct a set $E_{\bar x,m}\in\Sigma^0_2$ satisfying
\[
        F_{\bar x,4m}\subseteq E_{\bar x,m}\subseteq F_{\bar x,m}.        \tag{*}
\]
Then
\[
\{K:\dim_{\mathrm{nuc}}(A/K)\le n\}
=
\bigcap_{\bar x\in(D^A)^{<\N}}\bigcap_{m\in\N} E_{\bar x,m},
\]
because the left-hand side is equivalently
$\bigcap_{\bar x,m}F_{\bar x,m}=\bigcap_{\bar x,m}F_{\bar x,4m}$.  Since each
$E_{\bar x,m}$ is $\Sigma^0_2$, this will give a $\Pi^0_3$ set.

Fix a finite-dimensional algebra type
\[
F_0=\bigoplus_{j=1}^{\ell}M_{d_j},
\qquad
N(F_0):=\sum_{j=1}^{\ell}d_j^2,
\]
and matrix units $e_{ab}^{(j)}$ for each block.  Put
\[
M_{\bar x}:=1+\max_{1\le t\le r}\|x_t\|,
\qquad
\varepsilon_{F_0,\bar x,m}:=
\frac{1}{4m(n+1)N(F_0)M_{\bar x}},
\]
and let
\[
\delta_{F_0,\bar x,m}:=\delta_{F_0}(\varepsilon_{F_0,\bar x,m})
\]
be the stability tolerance from Lemma~\ref{lem:order-zero-coding}.

For each $i=0,\dots,n$, a candidate order-zero map $F_0\to A/K$ is coded by a tuple
\[
\bar d^{(i)}=(d_{ab}^{(i,j)})_{j,a,b}\in(D^A)^{N(F_0)},
\]
and we write
\[
\bar d=(\bar d^{(0)},\dots,\bar d^{(n)})\in(D^A)^{(n+1)N(F_0)}.
\]
Let $r_1(\bar z),\dots,r_M(\bar z)$ be the finite family of order-zero relations from
Lemma~\ref{lem:order-zero-coding}.  For fixed $F_0$ and $\bar d$, define
$\Lambda_{F_0,\bar d}$ to be the set of all $K$ such that, for every $i\le n$ and
$1\le s\le M$,
\[
        \Phi(K,r_s(\bar d^{(i)}))<\delta_{F_0,\bar x,m}.
\]
This is open, since $K\mapsto\Phi(K,a)$ is Wijsman-continuous for fixed $a\in A$.

Let $\mathrm{CPC}(A,F_0)$ be the compact metrisable space of c.p.c. maps
$A\to F_0$.  For fixed $F_0$ and $\bar d$, let
\[
C_{F_0,\bar d}\subseteq \mathrm{CPC}(A,F_0)\times\Ideal(A)
\]
be the set of pairs $(\psi,K)$ such that
\[
        \|\psi(a)\|\le \Phi(K,a)\qquad(a\in D^A),                    \tag{1}
\]
so that $\psi$ descends to $A/K$, and for every $1\le t\le r$,
\[
\Phi\Bigl(
K,
\sum_{i=0}^{n}\sum_{j=1}^{\ell}\sum_{a,b\le d_j}
\psi(x_t)^{(j)}_{ab}d_{ab}^{(i,j)}-x_t
\Bigr)
<\frac{1}{2m}.                                                        \tag{2}
\]
Condition~(1) is closed, and condition~(2) is open in
$\mathrm{CPC}(A,F_0)\times\Ideal(A)$.  Hence $C_{F_0,\bar d}$ is locally closed.
By Lemma~\ref{lem:proj-locally-closed}, its projection
$P_{F_0,\bar d}$ to $\Ideal(A)$ is $F_\sigma$.  Set
\[
        E_{\bar x,m}:=
        \bigcup_{F_0}\bigcup_{\bar d\in(D^A)^{(n+1)N(F_0)}}
        (\Lambda_{F_0,\bar d}\cap P_{F_0,\bar d}),
\]
where the first union runs over all finite-dimensional algebra types.  This is a
countable union of $F_\sigma$ sets, hence belongs to $\Sigma^0_2$.

It remains to prove (*).  Suppose first that $K\in F_{\bar x,4m}$, witnessed by
$F_0$, a c.p.c. map $\psi_0:A/K\to F_0$, and c.p.c. order-zero maps
$\varphi_i:F_0\to A/K$.  Put $\psi=\psi_0\circ q_K$.  For each $i,j,a,b$ choose
$d_{ab}^{(i,j)}\in D^A$ sufficiently close to a lift of
$\varphi_i(e_{ab}^{(j)})$.  Since the exact order-zero relations vanish modulo $K$,
and since only finitely many coefficients $\psi(x_t)^{(j)}_{ab}$ occur, the choices may
be made so that $K\in\Lambda_{F_0,\bar d}$ and (2) holds with error $<1/(2m)$.
Thus $K\in E_{\bar x,m}$.

Conversely, suppose $K\in E_{\bar x,m}$.  Then for some $F_0$ and $\bar d$ we have
$K\in\Lambda_{F_0,\bar d}\cap P_{F_0,\bar d}$.  The first condition and
Lemma~\ref{lem:order-zero-coding} give c.p.c. order-zero maps
$\varphi_i:F_0\to A/K$ satisfying
\[
\|\varphi_i(e_{ab}^{(j)})-(d_{ab}^{(i,j)}+K)\|
<\varepsilon_{F_0,\bar x,m}.
\]
The second condition gives a descending c.p.c. map $\psi:A/K\to F_0$ satisfying (2).
For $\bar x_t=x_t+K$, contractivity gives
$|\psi(\bar x_t)^{(j)}_{ab}|\le M_{\bar x}$, and therefore
\[
\begin{aligned}
\Bigl\|\sum_{i=0}^{n}\varphi_i\psi(\bar x_t)-\bar x_t\Bigr\|
&\le
\Bigl\|
\sum_{i,j,a,b}\psi(\bar x_t)^{(j)}_{ab}(d_{ab}^{(i,j)}+K)-\bar x_t
\Bigr\| \\
&\quad+
\sum_{i,j,a,b}|\psi(\bar x_t)^{(j)}_{ab}|
\|\varphi_i(e_{ab}^{(j)})-(d_{ab}^{(i,j)}+K)\| \\
&<\frac1{2m}+(n+1)N(F_0)M_{\bar x}\varepsilon_{F_0,\bar x,m}
<\frac1m.
\end{aligned}
\]
Hence $K\in F_{\bar x,m}$, proving the second inclusion in (*).
\end{proof}

\subsection{Property (SP)}

\begin{proposition}\label{prop:SP}
A $C^*$-algebra $B$ has \emph{property (SP)} if every non-zero hereditary subalgebra
of $B$ contains a non-zero projection (see \cite[V.2.2.14]{BlackadarOA}).
The class of $K\in\Ideal(A)$ such that $A/K$ has property~(SP)
is $\Pi^0_3$ in the Wijsman topology.
\end{proposition}

\begin{proof}
Write $B=A/K$.
$B$ has property~(SP) iff for every non-zero positive $a\in B_+$ there exists a
\emph{non-zero} projection $p\in\overline{aBa}$.

By Lemma~\ref{lem:positivity-coding}(ii), it suffices to range over $a=b^*b+K$
with $b\in D^A$, and to express non-zeroness by $\Phi(K,b^*b)>0$.
Indeed, the set of non-zero positive $a\in B_+$ such that $\overline{aBa}$
contains a non-zero projection is open in $B_+\setminus\{0\}$: if
$q\neq 0$ is a projection in $\overline{aBa}$, choose
$x\in B$ with
$\|q-axa\|<1/4$.
(This is possible because $aBa$ is dense in the hereditary $C^*$-subalgebra $\overline{aBa}$.)
Then the same inequality persists for all $a'$ sufficiently close to $a$,
so $\overline{a'Ba'}$ also contains a non-zero projection.

Fix a tolerance $\eta>0$ small enough that if $y\in B$ satisfies
$\|y^2-y\|<\eta$ and $\|y^*-y\|<\eta$, then functional calculus produces a projection
$p(y)$ with $\|p(y)-y\|<1/4$ (any standard $\eta$ from \cite{Loring} suffices).
In particular, if additionally $\|y\|>1/2$, then $p(y)\neq 0$.

For $b\in D^A$ and $n\in\N$ consider the condition
\begin{multline*}
\Phi(K,b^*b)=0
\ \lor\ 
\exists p,x\in D^A:\\
\Phi(K,p^2-p)<\eta\ \wedge\ \Phi(K,p^*-p)<\eta\ \wedge\ \Phi(K,p)>\tfrac12\\
\wedge\ \Phi\bigl(K,p-b^*b\,x\,b^*b\bigr)<\tfrac1n.
\end{multline*}
If $\Phi(K,b^*b)>0$, the second disjunct asserts the existence of a non-zero
approximate projection $p+K$ lying within $1/n$ of the hereditary subalgebra
\[
    H:=\overline{(b^*b+K)B(b^*b+K)}.
\]
Choose $h\in H$ with $\|h-(p+K)\|<1/n$; for large $n$, the element $h$ is also
an approximate (self-adjoint) idempotent.
Since $H$ is a closed $C^*$-subalgebra, functional calculus \emph{inside $H$}
produces a genuine projection $q\in H$ close to $h$, and $\|q\|\ge\|p+K\|-1/4>0$
ensures $q\neq 0$.
Thus $H$ contains a non-zero projection.

For fixed parameters, the first disjunct is closed and the second is a countable union
of open sets (strict $\Phi$-inequalities), hence the whole condition is $\Sigma^0_2$.
Intersecting over $b\in D^A$ and $n\in\N$ gives a $\Pi^0_3$ set.
\end{proof}

\subsection{Tracial states}

\begin{proposition}\label{prop:trace}
The class of $K\in\Ideal(A)$ such that $A/K$ admits a tracial state is closed
(in particular $\Pi^0_1$) in the Wijsman topology.
\end{proposition}

\begin{proof}
Let $T(A)$ be the compact (metrisable) simplex of tracial states on $A$
with its weak$^*$ topology.
A tracial state on $B=A/K$ is the same as a tracial state $\tau\in T(A)$
vanishing on $K$ (compose with the quotient map $A\to A/K$, and conversely factor
through $A/K$).

Fix the dense $*$-subalgebra $D^A\subseteq A$.
For $\tau\in T(A)$ and $K\in\Ideal(A)$, the condition $K\subseteq\ker(\tau)$ is
equivalent to
\[
\forall a\in D^A:\ |\tau(a)|\le \|a+K\|_{A/K}=\Phi(K,a),
\]
since $\tau$ is contractive and $\Phi(K,k)=0$ for $k\in K$.

For each fixed $a\in D^A$, the map $(\tau,K)\mapsto |\tau(a)|-\Phi(K,a)$ is continuous
on $T(A)\times\Ideal(A)$ (weak$^*$ continuity in $\tau$ and Wijsman continuity in $K$),
hence the inequality defines a closed subset.
Therefore
\[
C:=\{(\tau,K)\in T(A)\times\Ideal(A):\ K\subseteq\ker(\tau)\}
\]
is closed.
Since $T(A)$ is compact, the projection of $C$ onto $\Ideal(A)$ is closed, and this
projection is exactly the set of ideals $K$ for which $A/K$ carries a tracial state.
\end{proof}

\subsection{Crossed products and Rokhlin property}

\begin{remark}\label{rem:crossed-Rokhlin}
A full treatment of crossed-product codes and action codes would require an
additional parameter-space construction for compact metrisable spaces, group
actions, and quotient actions in the present framework.
Such a treatment should lead to analytic/Borel upper bounds for the following classes:
for a fixed countable discrete group~$G$, the class of $K$ such that $A/K$ is
isomorphic to a crossed product $C(X)\rtimes_\alpha G$ should be analytic;
and for coded actions $\alpha:G\curvearrowright A/K$, the Rokhlin property should
be $\Pi^0_3$ in the pair $(K,\alpha)$.
These facts are not used elsewhere in the paper, so we omit the details.
\end{remark}

\begin{remark}\label{rem:AI-auto}
Let $B=A/K$.

\smallskip\noindent
\textbf{Coding automorphisms.}
In our quotient coding, it is convenient to regard an automorphism $\alpha\in\Aut(B)$
as being specified by its values on the named dense $*$-subalgebra $D^A+K\subseteq B$.
Concretely, a \emph{code} for $\alpha$ may be taken to be a sequence
$\bar y=(y_m)_{m\in\N}$ in $D^A$ intended to represent $\alpha(u_m+K)=y_m+K$,
together with the requirement that the assignment $u_m+K\mapsto y_m+K$ extends to a
$*$-automorphism of $B$.
This extension condition can be written as a countable family of constraints
expressing $*$-homomorphism identities on $D^A$ (e.g.\ preservation of products,
$*$, and rational linear combinations) and the existence of a coded inverse; these
are Borel conditions in the pair $(K,\bar y)$ by the general definability scheme
(Theorem~\ref{thm:definability-rank}).

\smallskip\noindent
\textbf{Approximate innerness.}
An automorphism $\alpha\in\Aut(B)$ is \emph{approximately inner} if it belongs to the
point-norm closure of the inner automorphism group $\Inn(B)$, equivalently:
for every finite $F\subseteq B$ and every $\varepsilon>0$ there exists a unitary
$u\in B$ such that
\[
\max_{x\in F}\|\alpha(x)-uxu^*\|<\varepsilon.
\]
Since $D^A+K$ is dense, it suffices to test finite $F\subseteq D^A+K$ and
$\varepsilon=1/n$.
Moreover, the unitary quantifier can be reduced to the fixed countable dense set
$D^A$ by stability of the unitary relation:
if $v\in B$ is sufficiently close to a unitary (\emph{i.e.}\ $\|v^*v-1\|$ and $\|vv^*-1\|$
are small), then functional calculus produces a genuine unitary $u(v)$ close to~$v$.
Thus, for a fixed tolerance $\eta>0$ (small enough for unitary stability),
approximate innerness of $\alpha$ is equivalent to the following countable scheme:
\begin{multline*}
\forall m\in\N\ \forall n\in\N\ \exists v\in D^A:\\ 
\Phi(K,v^*v-1)<\eta\ \wedge\ \Phi(K,vv^*-1)<\eta\\
\wedge\ 
\max_{j\le m}\Phi\!\bigl(K,\alpha(u_j+K)-v(u_j+K)v^*\bigr)<\tfrac{1}{n}.
\end{multline*}
Here $\alpha(u_j+K)$ is read from the chosen automorphism code (e.g.\ via the
representatives $y_j+K$).
The displayed predicate is a finite conjunction of strict $\Phi$-inequalities, hence
open; the quantifier pattern is $\forall\forall\exists(\text{open})$.
Therefore, the set of coded pairs $(K,\alpha)$ with $\alpha$ approximately inner is
a $\Pi^0_3$ subset of the ambient coding space.

\smallskip\noindent
\textbf{Context.}
For each fixed $B$ (separable), $\Aut(B)$ equipped with the point-norm topology is a
Polish group and $\Aut_{\mathrm{ai}}(B)$ is a closed normal subgroup (being the
closure of $\Inn(B)$).  Globally, in the quotient coding, the previous paragraph
identifies approximate innerness as a uniform Borel condition on \emph{pairs}
$(K,\alpha)$.
\end{remark}

\begin{remark}\label{rem:summary}
In the Wijsman quotient coding, the rank estimates proved in this section are as
follows.  Nuclearity and simplicity are Borel.  Stable finiteness and the existence of a
tracial state are closed.  AF, MF, approximate divisibility, real-rank bounds and
topological stable-rank bounds are $G_\delta$; equivalently, they are $\Pi^0_2$.
The bounds for AF, real-rank bounds and topological stable-rank bounds are optimal by
Theorem~\ref{thm:Cstar-Pi2-complete}.  Quasidiagonality, property~(SP), and nuclear
dimension~$\le n$ are $\Pi^0_3$; for property~(SP) and nuclear dimension the same
theorem gives $\Pi^0_2$ lower bounds, but not $\Pi^0_3$-hardness.
Local approximation by homomorphic images of a countable family of stably finitely
presented blocks is $G_\delta$.  For a fixed separable unital exact algebra $D$,
$D$-absorption is analytic.
\end{remark}

\subsection{Sharpness of some \texorpdfstring{$G_\delta$}{G-delta} bounds}

The previous estimates give upper bounds.  We now record lower bounds for three of
the $G_\delta$ classes.  The reductions are commutative and therefore live inside a
very small part of the quotient space.  This still proves hardness in the full quotient
coding, because the reduction constructed below is a continuous map into the
ambient Wijsman ideal space.

Let
\[
        P_2=\{\alpha\in 2^{\N\times\N}:\forall r\;\exists s\; \alpha(r,s)=1\}.
\]
This is a standard $\Pi^0_2$-complete subset of the Cantor space $2^{\N\times\N}$.

\begin{lemma}\label{lem:compact-Pi2-test}
There is a continuous map
\[
        \alpha\longmapsto E_\alpha
\]
from $2^{\N\times\N}$ into the hyperspace of non-empty compact subsets of $[0,1]$,
with the Hausdorff topology, such that:
\begin{enumerate}
\item if $\alpha\in P_2$, then $E_\alpha$ is countable, hence zero-dimensional;
\item if $\alpha\notin P_2$, then $E_\alpha$ contains a non-degenerate interval.
\end{enumerate}
\end{lemma}

\begin{proof}
Choose pairwise disjoint closed intervals
\[
        I_r=[2^{-r-2},2^{-r-1}]\qquad(r\in\N),
\]
which accumulate only at~$0$.  For each pair $(r,s)$ choose a finite set
$F_{r,s}\subseteq I_r$, containing the endpoints of $I_r$, such that
\[
        d_H(F_{r,s},I_r)<2^{-s},
\]
where $d_H$ denotes Hausdorff distance.  For $\alpha\in2^{\N\times\N}$ put
\[
        m_r(\alpha)=\min\{s:\alpha(r,s)=1\},
\]
with $m_r(\alpha)=\infty$ if the set is empty, and define
\[
        E_\alpha=
        \{0\}
        \cup\bigcup_{m_r(\alpha)=\infty} I_r
        \cup\bigcup_{m_r(\alpha)<\infty} F_{r,m_r(\alpha)} .
\]
This is compact because the intervals $I_r$ accumulate only at~$0$.
If every row of $\alpha$ contains a~$1$, then $E_\alpha$ is a countable compact set.
If some row contains no~$1$, then $E_\alpha$ contains the corresponding interval~$I_r$.

It remains only to check continuity.  Let $\alpha_j\to\alpha$.  Given $\varepsilon>0$,
choose $R$ so large that
\[
        \{0\}\cup\bigcup_{r>R} I_r \subseteq [0,\varepsilon].
\]
For each $r\le R$ with $m_r(\alpha)<\infty$, the first $m_r(\alpha)+1$ entries in row
$r$ are eventually constant, so the corresponding piece of $E_{\alpha_j}$ is eventually
exactly the same as the piece of $E_\alpha$.  For each $r\le R$ with
$m_r(\alpha)=\infty$, choose $S$ so large that $2^{-S}<\varepsilon$; eventually the
first $S$ entries in row~$r$ of $\alpha_j$ are all zero, and therefore the corresponding
piece of $E_{\alpha_j}$ is either $I_r$ or some $F_{r,s}$ with $s\ge S$, in both cases
within Hausdorff distance $<\varepsilon$ from $I_r$.  Combining the finitely many
large intervals with the small tail gives $d_H(E_{\alpha_j},E_\alpha)\to0$.
\end{proof}

For $d\ge0$ set
\[
        Y_{\alpha,d}:=E_\alpha\times [0,1]^d\subseteq [0,1]^{d+1},
\]
where $[0,1]^0$ is a point.  Fix once and for all a surjection
\[
        \rho_d:C^*_{\max}(F_\infty)\longrightarrow C([0,1]^{d+1}).
\]
For a compact set $Y\subseteq [0,1]^{d+1}$ write
\[
        I_Y=\{f\in C([0,1]^{d+1}):f|_Y=0\},
        \qquad
        K_Y=\rho_d^{-1}(I_Y).
\]
The map $Y\mapsto K_Y$ is Wijsman-continuous: for $a\in C^*_{\max}(F_\infty)$,
\[
        \dist(a,K_Y)=\|\rho_d(a)+I_Y\|=\|\rho_d(a)|_Y\|_\infty
        =\max_{y\in Y}|\rho_d(a)(y)|,
\]
and the last expression is continuous in $Y$ for the Hausdorff topology.
Since $Y\mapsto I_Y$ is injective, compactness of the hyperspace and metrisability of the
Wijsman ideal space show that $Y\mapsto K_Y$ is a topological embedding onto its image.
Thus a continuous hardness reduction landing in this commutative subfamily is already a
hardness reduction for the full quotient coding.

\begin{theorem}\label{thm:Cstar-Pi2-complete}
In the Wijsman quotient coding of separable unital $C^*$-algebras by ideals of
$C^*_{\max}(F_\infty)$, the following classes are $\Pi^0_2$-complete:
\begin{enumerate}
\item AF algebras;
\item algebras of real rank at most $d$, for each fixed $d\ge0$;
\item algebras of topological stable rank at most $n$, for each fixed $n\ge1$.
\end{enumerate}
Moreover, for each fixed $d\ge0$, the class of algebras with nuclear dimension at most
$d$ is $\Pi^0_2$-hard, and property~\textup{(SP)} is $\Pi^0_2$-hard.
\end{theorem}

\begin{proof}
Membership in $\Pi^0_2$ for AF, real-rank bounds and topological stable-rank bounds was
proved above.  We prove hardness by reducing $P_2$.

For AF-ness use $Y_{\alpha,0}=E_\alpha$.  If $\alpha\in P_2$, then $E_\alpha$ is
countable and compact, hence zero-dimensional, and $C(E_\alpha)$ is AF.  If
$\alpha\notin P_2$, then $E_\alpha$ contains a non-degenerate interval, so
$C(E_\alpha)$ is not AF.  Thus
\[
        \alpha\in P_2
        \quad\Longleftrightarrow\quad
        C^*_{\max}(F_\infty)/K_{E_\alpha}\cong C(E_\alpha)\text{ is AF}.
\]
Since $\alpha\mapsto K_{E_\alpha}$ is continuous, AF-ness is $\Pi^0_2$-hard.

For real rank, use $Y_{\alpha,d}=E_\alpha\times[0,1]^d$.  If $\alpha\in P_2$, then
$E_\alpha$ is zero-dimensional, so standard dimension theory gives
$\dim(Y_{\alpha,d})=d$.  If $\alpha\notin P_2$, then $Y_{\alpha,d}$ contains a copy of
$[0,1]^{d+1}$, so $\dim(Y_{\alpha,d})\ge d+1$.  Since real rank of a commutative
unital $C^*$-algebra is covering dimension~\cite{BrownPedersen1991}, $\mathrm{rr}(C(Y_{\alpha,d}))\le d$ holds
exactly for $\alpha\in P_2$.  This proves $\Pi^0_2$-hardness of the real-rank bound.

For topological stable rank~$\le n$, put $d=2n-1$.  The same dimension computation
gives dimension $2n-1$ when $\alpha\in P_2$ and dimension $2n$ otherwise; the
choice $2n-1$ is exactly the threshold between stable ranks $n$ and $n+1$.
Rieffel's formula
$\mathrm{tsr}(C(X))=\lfloor \dim(X)/2\rfloor+1$ for compact spaces of finite covering
dimension~\cite{Rieffel1983} therefore yields
\[
        \mathrm{tsr}(C(Y_{\alpha,2n-1}))\le n
        \quad\Longleftrightarrow\quad
        \alpha\in P_2.
\]
This proves $\Pi^0_2$-hardness of the stable-rank bound.

The final assertions are proved by the same commutative family.  Winter--Zacharias'
formula $\dim_{\mathrm{nuc}}(C(X))=\dim(X)$ for compact metrisable $X$
\cite{WinterZacharias} gives $\Pi^0_2$-hardness of nuclear dimension~$\le d$.
For property~(SP), $C(E_\alpha)$ has property~(SP) when $E_\alpha$ is zero-dimensional:
every non-empty open subset contains a non-empty clopen subset, whose characteristic
function is a non-zero projection.  If $E_\alpha$ contains an interval, a hereditary
subalgebra supported in a smaller open interval contains no non-zero projection.  Hence
property~(SP) also has $P_2$ as a continuous preimage.
\end{proof}

%%%%%%%%%%%%%%%%%%%%%%%%%%%%%%%%%%%%%%%%%%%%%%%%%%%%%%%%%%%%%%
\section{Borel assignment of \texorpdfstring{$K$}{K}-theory}\label{sec:K0}
%%%%%%%%%%%%%%%%%%%%%%%%%%%%%%%%%%%%%%%%%%%%%%%%%%%%%%%%%%%%%%

We give a fully internal Borel coding of the $K_0$-assignment for the unital
quotients considered in Sections~\ref{sec:Cstar-properties}--\ref{sec:D-stability}.
For $K_1$ and higher $K$-groups we use suspension and Bott periodicity together
with the already-known Borel computability of $K$-theory in the standard
$\Gamma/\Xi$ codings.  A completely internal treatment of the non-unital
suspension case would require the relative/unitised version of the $K_0$ coding.
These results can also be extended to ordered $K$-theory, but we do not pursue that
direction here.

Recall that for a separable $C^*$-algebra $B$, the group $K_0(B)$ may be defined as
the Grothendieck group of Murray--von Neumann equivalence classes of projections in
$(B\otimes\mathcal K)^+$ under direct sum.  When $B$ is unital, this agrees with the
Grothendieck group of the commutative monoid $V(B)$ of projections in
$\bigsqcup_d M_d(B)$ under direct sum.
The group $K_1(B)$ is defined as $\pi_0(\mathrm{GL}_\infty(\widetilde B))$, where
$\widetilde B$ is the unitisation of $B$, or equivalently as $K_0(SB)$ where
$SB=C_0((0,1))\otimes B$ is the suspension; the latter identification follows from
Bott periodicity.
For the precise definitions and basic properties of $K$-theory for operator algebras we
refer to R{\o}rdam--Larsen--Laustsen~\cite{RordamLarsenLaustsen},
Blackadar~\cite{Blackadar}, and Wegge-Olsen~\cite{WeggeOlsen}; we use these
definitions without further comment.


\begin{remark}
\label{rem:K-FTT-compare}
Farah--Toms--T\"ornquist proved that ordered $K_0$ and $K_1$ are Borel computable in
the standard $\Gamma/\Xi$ parameterisations; see Section~3 of
\cite{FarahTomsTornquist}.
By Remarks~\ref{rem:other-codings} and~\ref{rem:Xi-ideals}, those codings are Borel
equivalent to the present quotient-ideal coding.
Thus the existence of a Borel $K$-theory assignment in our framework is known
\emph{a priori}.  The point of the present section is to construct such assignments
internally in the quotient/Wijsman setting and to record some explicit complexity
information there.
\end{remark}

\subsection{The coding space of countable abelian groups}

Let $\Z^{(\N)}=\bigoplus_{n\in\N}\Z e_n$.  The space
\[
\Sub(\Z^{(\N)})=\{H\subseteq\Z^{(\N)}:H\text{ is a subgroup}\}\subseteq 2^{\Z^{(\N)}}
\]
is a closed subset of a Cantor cube and hence compact Polish.
Each $H$ codes the quotient group $\Z^{(\N)}/H$.

\subsection{An internal Borel presentation of \texorpdfstring{$K_0$}{K0}}

Fix a unital separable $C^*$-algebra quotient generator $A$ and a countable dense
$*$-subalgebra $D^A\subseteq A$.
For each $d\in\N$, let $M_d(D^A)_+$ denote the positive cone of $M_d(D^A)$.
Choose once and for all an enumeration
\[
(a_n)_{n\in\N}
\]
of the countable set
\[
\mathcal P^A:=\bigsqcup_{d\in\N} M_d(D^A)_+
\]
with the following properties:
\begin{enumerate}
\item $a_0=0\in M_1(D^A)$;
\item if $a_m\in M_{d(m)}(D^A)_+$ and $a_n\in M_{d(n)}(D^A)_+$, then there is a
computable index $m\oplus n$ such that
\[
a_{m\oplus n}=a_m\oplus a_n\in M_{d(m)+d(n)}(D^A)_+.
\]
\end{enumerate}
Since every positive element is of the form $b^*b$ and $D^A$ is dense,
$\mathcal P^A$ is dense in the positive cone of $A\otimes\mathcal K$.

\begin{lemma}
\label{lem:K0-correct-proj}
For every $\varepsilon\in(0,1/4)$ there exist a constant
$\eta(\varepsilon)>0$ and a continuous function
$\chi:[0,\infty)\to[0,1]$ with
\[
\chi|_{[0,1/4]}=0,
\qquad
\chi|_{[3/4,\infty)}=1,
\]
such that for every $C^*$-algebra $B$ and every positive element $x\in B_+$ satisfying
\[
\|x^2-x\|<\eta(\varepsilon),
\]
the element $\widehat x:=\chi(x)$ is a projection and
\[
\|\widehat x-x\|<\varepsilon.
\]
\end{lemma}

\begin{proof}
Choose $\eta(\varepsilon)>0$ small enough that
$|t^2-t|<\eta(\varepsilon)$ for $t\in[0,\infty)$ forces
$t\in[0,\varepsilon)\cup(1-\varepsilon,1+\varepsilon)$.
Functional calculus then gives the asserted spectral projection and the norm estimate.
\end{proof}


\begin{lemma}
\label{lem:K0-stable-MvN}
There exists $\delta_{\mathrm{MvN}}>0$ such that the following holds.

Let $B$ be a $C^*$-algebra, let $p,q,v\in B$, and assume
\[
\|p^2-p\|,\ \|p^*-p\|,\ \|q^2-q\|,\ \|q^*-q\|<\delta_{\mathrm{MvN}},
\]
and
\[
\|v^*v-p\|<\delta_{\mathrm{MvN}},
\qquad
\|vv^*-q\|<\delta_{\mathrm{MvN}}.
\]
Then the corrected projections $\widehat p$ and $\widehat q$
(from Lemma~\ref{lem:K0-correct-proj}) are Murray--von Neumann equivalent.

Conversely, if $e,f\in B$ are Murray--von Neumann equivalent projections, then for every
$\varepsilon>0$ and every dense $*$-subalgebra $B_0\subseteq B$ there exist
$p,q,v\in B_0$ such that
\[
\|p-e\|,\ \|q-f\|,\ \|v-w\|<\varepsilon
\]
for some partial isometry $w$ with $w^*w=e$ and $ww^*=f$, and the above
approximate relations hold.
\end{lemma}

\begin{proof}
This is the stability of the finite relations defining a corner copy of $M_2$.
Indeed, the exact relations
\[
e_{11}=e_{11}^*=e_{11}^2,\quad
e_{22}=e_{22}^*=e_{22}^2,\quad
e_{12}^*=e_{21},\quad
e_{12}e_{21}=e_{11},\quad
e_{21}e_{12}=e_{22}
\]
give a finite presentation of $M_2$, hence are weakly stable by semiprojectivity of
finite-dimensional $C^*$-algebras; cf.\ Lemma~\ref{lem:stable-fd} and~\cite{Loring}.
The converse follows by approximating an exact partial isometry witness in a dense
$*$-subalgebra.
\end{proof}
Fix $\varepsilon_0\in(0,1/4)$ so small that
$\varepsilon_0<\delta_{\mathrm{MvN}}/4$, and set
\[
\eta_0:=\min\{\eta(\varepsilon_0),\delta_{\mathrm{MvN}}/4\}.
\]
In the sequel, corrected projections are formed using this value of $\eta_0$ and the
corresponding function $\chi$ from Lemma~\ref{lem:K0-correct-proj}.

For $K\in\Ideal(A)$ and $n\in\N$, define the \emph{good approximate projection}
predicate
\[
\mathcal G(K,n)\iff \Phi^{(d(n))}(K,a_n^2-a_n)<\eta_0.
\]
If $\mathcal G(K,n)$ holds, let
\[
p_n^K:=\chi(a_n+M_{d(n)}(K))\in M_{d(n)}(A/K)
\]
be the corrected projection from Lemma~\ref{lem:K0-correct-proj}; if
$\mathcal G(K,n)$ fails, set $p_n^K:=0$.

We also define the \emph{zero-class} predicate
\[
\mathcal Z(K,n)\iff
\neg\mathcal G(K,n)\ \lor\ \Phi^{(d(n))}(K,a_n)<1/2.
\]
If $\mathcal G(K,n)$ holds, then $\mathcal Z(K,n)$ is equivalent to $p_n^K=0$.

\begin{lemma}
\label{lem:K0-dense-proj}
For every $K\in\Ideal(A)$ and every projection
$p\in M_d(A/K)$ there exists $n\in\N$ with $d(n)=d$ such that
$p_n^K$ is Murray--von Neumann equivalent to $p$.
Hence the set of classes $\{[p_n^K]:n\in\N\}$ is all of
$V(A/K)$.
\end{lemma}

\begin{proof}
Fix a projection $p\in M_d(A/K)$.
Choose $\varepsilon>0$ so small that $2\varepsilon+\varepsilon^2<\eta_0$ and
$\varepsilon<1/2$.
Since $\mathcal P^A$ is dense in the positive cone, pick $n$ with $d(n)=d$ and
\[
\|a_n+M_d(K)-p\|<\varepsilon.
\]
Then
\[
\|(a_n+M_d(K))^2-(a_n+M_d(K))\|
\le 2\varepsilon+\varepsilon^2<\eta_0,
\]
so $\mathcal G(K,n)$ holds.
By Lemma~\ref{lem:K0-correct-proj}, the corrected projection $p_n^K$ satisfies
\[
\|p_n^K-(a_n+M_d(K))\|<1/4.
\]
Therefore
\[
\|p_n^K-p\|<1/4+\varepsilon<1.
\]
Projections at distance strictly less than $1$ are unitarily equivalent, hence
Murray--von Neumann equivalent.
\end{proof}

For $m,n\in\N$, let
\[
r(m,n):=d(m)+d(n),
\]
and define
\[
b_{m,n}^{(1)}:=a_m\oplus 0_{d(n)},
\qquad
b_{m,n}^{(2)}:=0_{d(m)}\oplus a_n
\]
as elements of $M_{r(m,n)}(A)_+$.

For $m,n,k\in\N$, let
\[
s(m,n,k):=d(m)+d(n)+d(k),
\]
and define
\[
c_{m,n,k}^{(1)}:=a_m\oplus a_n\oplus 0_{d(k)},
\qquad
c_{m,n,k}^{(2)}:=0_{d(m)+d(n)}\oplus a_k
\]
as elements of $M_{s(m,n,k)}(A)_+$.

Now define relations
\[
\mathcal E\subseteq \Ideal(A)\times\N^2,
\qquad
\mathcal A\subseteq \Ideal(A)\times\N^3
\]
by the following clauses.

\smallskip
\noindent\emph{Equivalence relation.}
We declare $\mathcal E(K,m,n)$ to hold iff either
\[
\mathcal Z(K,m)\wedge \mathcal Z(K,n),
\]
or else there exists
$v\in M_{r(m,n)}(D^A)$ such that
\[
\Phi^{(r(m,n))}\bigl(K,(b_{m,n}^{(1)})^2-b_{m,n}^{(1)}\bigr)<\eta_0,
\]
\[
\Phi^{(r(m,n))}\bigl(K,(b_{m,n}^{(2)})^2-b_{m,n}^{(2)}\bigr)<\eta_0,
\]
and
\[
\Phi^{(r(m,n))}\bigl(K,v^*v-b_{m,n}^{(1)}\bigr)<\delta_{\mathrm{MvN}},
\qquad
\Phi^{(r(m,n))}\bigl(K,vv^*-b_{m,n}^{(2)}\bigr)<\delta_{\mathrm{MvN}}.
\]

\smallskip
\noindent\emph{Addition relation.}
We declare $\mathcal A(K,m,n,k)$ to hold iff one of the following alternatives holds:
\begin{enumerate}[label=\textup{(A\arabic*)}]
\item $\mathcal Z(K,m)$ and $\mathcal E(K,n,k)$;
\item $\mathcal Z(K,n)$ and $\mathcal E(K,m,k)$;
\item there exists $v\in M_{s(m,n,k)}(D^A)$ such that
\[
\Phi^{(s(m,n,k))}\bigl(K,(c_{m,n,k}^{(1)})^2-c_{m,n,k}^{(1)}\bigr)<\eta_0,
\]
\[
\Phi^{(s(m,n,k))}\bigl(K,(c_{m,n,k}^{(2)})^2-c_{m,n,k}^{(2)}\bigr)<\eta_0,
\]
and
\[
\Phi^{(s(m,n,k))}\bigl(K,v^*v-c_{m,n,k}^{(1)}\bigr)<\delta_{\mathrm{MvN}},
\qquad
\Phi^{(s(m,n,k))}\bigl(K,vv^*-c_{m,n,k}^{(2)}\bigr)<\delta_{\mathrm{MvN}}.
\]
\end{enumerate}

\begin{lemma}\label{lem:K0-relations}
The relations $\mathcal Z$, $\mathcal E$, and $\mathcal A$ are Borel in $K$.
Moreover, for every $K\in\Ideal(A)$:
\begin{enumerate}
\item $\mathcal Z(K,n)$ iff $p_n^K=0$;
\item $\mathcal E(K,m,n)$ iff $[p_m^K]=[p_n^K]$ in $V(A/K)$, equivalently iff
\[
p_m^K\oplus 0_{d(n)} \sim 0_{d(m)}\oplus p_n^K
\]
inside $M_{r(m,n)}(A/K)$;
\item $\mathcal A(K,m,n,k)$ iff
\[
[p_m^K]+[p_n^K]=[p_k^K]
\]
in the semigroup $V(A/K)$.
\end{enumerate}
\end{lemma}

\begin{proof}
Borelness is immediate from the definitions: all displayed conditions are strict
inequalities in Wijsman-continuous functions, with existential quantification only over
countable dense sets.

For (1), if $\neg\mathcal G(K,n)$ then by definition $p_n^K=0$.
If $\mathcal G(K,n)$ holds, then $p_n^K=0$ iff the spectrum of $a_n+M_{d(n)}(K)$ is
contained in $[0,1/4]$, which is equivalent to
$\|a_n+M_{d(n)}(K)\|<1/2$.

For (2), the forward implication follows from Lemma~\ref{lem:K0-stable-MvN} applied to
the approximate source/range data in the common amplification $M_{r(m,n)}(A/K)$; the
corrected projections of $b_{m,n}^{(1)}$ and $b_{m,n}^{(2)}$ are exactly
$p_m^K\oplus 0_{d(n)}$ and $0_{d(m)}\oplus p_n^K$.
Conversely, if these two amplified projections are Murray--von Neumann equivalent, approximate
an exact partial-isometry witness by an element of the dense $*$-subalgebra.

For (3), argue exactly as in (2), now comparing
\[
p_m^K\oplus p_n^K\oplus 0_{d(k)}
\quad\text{and}\quad
0_{d(m)+d(n)}\oplus p_k^K.
\]
\end{proof}

\begin{lemma}
\label{lem:Groth-presentation}
Let $M$ be a commutative monoid and let $\nu:\N\to M$ be surjective.
Assume given sets
\[
Z\subseteq \N,\qquad E\subseteq \N^2,\qquad A\subseteq \N^3
\]
such that
\[
\begin{aligned}
n\in Z &\iff \nu(n)=0,\\
(m,n)\in E &\iff \nu(m)=\nu(n),\\
(m,n,k)\in A &\iff \nu(m)+\nu(n)=\nu(k).
\end{aligned}
\]
Let $H\le \Z^{(\N)}$ be the subgroup generated by
\[
\{\mathbf e_n:n\in Z\}\cup
\{\mathbf e_m-\mathbf e_n:(m,n)\in E\}\cup
\{\mathbf e_m+\mathbf e_n-\mathbf e_k:(m,n,k)\in A\}.
\]
Then
\[
\Z^{(\N)}/H\cong G(M),
\]
where $G(M)$ denotes the Grothendieck group of $M$.
\end{lemma}

\begin{proof}
Let $\iota:M\to G(M)$ be the canonical map and define
\[
\pi:\Z^{(\N)}\to G(M),\qquad \pi(\mathbf e_n)=\iota(\nu(n)).
\]
By the defining properties of $Z,E,A$, every generator of $H$ lies in $\ker(\pi)$,
so $\pi$ factors through a surjective homomorphism
\[
\overline\pi:\Z^{(\N)}/H\to G(M).
\]

It remains to prove injectivity.
Take $x\in\ker(\pi)$ and write
\[
x=\sum_{i=1}^r \mathbf e_{m_i}-\sum_{j=1}^s \mathbf e_{n_j}.
\]
Since $\pi(x)=0$ in the Grothendieck group, there exists $c\in M$ such that
\[
\sum_{i=1}^r \nu(m_i)+c=\sum_{j=1}^s \nu(n_j)+c
\]
in $M$.  Choose $\ell\in\N$ with $\nu(\ell)=c$.

We claim that every finite sum of $\nu$-values can be reduced, modulo $H$, to a
single generator.  More precisely, for any indices $q_1,\dots,q_t$ there exists
$u\in\N$ such that
\[
\sum_{i=1}^t \mathbf e_{q_i}-\mathbf e_u\in H
\qquad\text{and}\qquad
\nu(u)=\sum_{i=1}^t \nu(q_i).
\]
This is proved by induction on $t$.
For $t=1$ there is nothing to show.
If $t\ge 2$, choose $u'$ with
\[
\nu(u')=\nu(q_1)+\nu(q_2),
\]
possible by surjectivity of $\nu$.
Then $(q_1,q_2,u')\in A$, so
\[
\mathbf e_{q_1}+\mathbf e_{q_2}-\mathbf e_{u'}\in H,
\]
and the induction hypothesis applies to $u',q_3,\dots,q_t$.

Applying this reduction to the two sums
\[
m_1,\dots,m_r,\ell
\qquad\text{and}\qquad
n_1,\dots,n_s,\ell
\]
yields indices $u,v\in\N$ with
\[
\sum_{i=1}^r \mathbf e_{m_i}+\mathbf e_\ell-\mathbf e_u\in H,
\qquad
\sum_{j=1}^s \mathbf e_{n_j}+\mathbf e_\ell-\mathbf e_v\in H,
\]
and
\[
\nu(u)=\sum_{i=1}^r \nu(m_i)+c=\sum_{j=1}^s \nu(n_j)+c=\nu(v).
\]
Hence $(u,v)\in E$, so $\mathbf e_u-\mathbf e_v\in H$.
Subtracting the two displayed relations and cancelling $\mathbf e_\ell$, we obtain
$x\in H$.
Thus $\ker(\overline\pi)=0$, and $\overline\pi$ is an isomorphism.
\end{proof}

\begin{theorem}\label{thm:K0}
Equip $\Ideal(A)$ with the Wijsman topology.  There is a Borel map
\[
\kappa_{K_0}:\Ideal(A)\to\Sub(\Z^{(\N)})
\]
such that
\[
\Z^{(\N)}/\kappa_{K_0}(K)\cong K_0(A/K)
\qquad (K\in\Ideal(A)).
\]
Moreover, for each fixed $z\in\Z^{(\N)}$, the section
\[
\{K\in\Ideal(A): z\in \kappa_{K_0}(K)\}
\]
is $F_\sigma$ in the Wijsman topology.
\end{theorem}

\begin{proof}
Let $(\mathbf e_n)_{n\in\N}$ denote the standard basis of $\Z^{(\N)}$.
For each $K\in\Ideal(A)$, let $H_K\le \Z^{(\N)}$ be the subgroup generated by the
following countable family:

\begin{enumerate}
\item $\mathbf e_n$ whenever $\mathcal Z(K,n)$ holds;
\item $\mathbf e_m-\mathbf e_n$ whenever $\mathcal E(K,m,n)$ holds;
\item $\mathbf e_m+\mathbf e_n-\mathbf e_k$ whenever $\mathcal A(K,m,n,k)$ holds.
\end{enumerate}

Fix $K$ and write $B=A/K$.
Let
\[
\nu_K:\N\to V(B),\qquad \nu_K(n):=[p_n^K].
\]
By Lemma~\ref{lem:K0-dense-proj}, the map $\nu_K$ is surjective.
By Lemma~\ref{lem:K0-relations}, the relations $\mathcal Z$, $\mathcal E$, and
$\mathcal A$ detect exactly the zero element, equality, and addition in the monoid
$V(B)$.
Therefore Lemma~\ref{lem:Groth-presentation} yields
\[
\Z^{(\N)}/H_K\cong G(V(B))=K_0(B)=K_0(A/K).
\]
Define $\kappa_{K_0}(K):=H_K$.

We next show that the map $K\mapsto H_K$ is Borel as a map into
$\Sub(\Z^{(\N)})$.
Fix $z\in\Z^{(\N)}$.  Then
\[
z\in H_K
\]
iff there exist finitely many generators $g_1,\dots,g_r$ of the above three types and
signs $\varepsilon_j\in\{\pm1\}$ such that
\[
z=\sum_{j=1}^{r}\varepsilon_j g_j
\]
and each of the corresponding activation relations $\mathcal Z$, $\mathcal E$, or
$\mathcal A$ holds at~$K$.
Since the potential generators form a fixed countable set and the activation relations are
Borel by Lemma~\ref{lem:K0-relations}, the set
\[
\{K:z\in H_K\}
\]
is a countable union of Borel sets, hence Borel.
Therefore $K\mapsto H_K$ is Borel as a map into $\Sub(\Z^{(\N)})$.

Finally, the coordinate-complexity claim is read off directly from the construction.
The predicate $\mathcal G(K,n)$ is open, so $\neg\mathcal G(K,n)$ is closed, and
$\mathcal Z(K,n)$ is $F_\sigma$.  The witness alternatives in the definitions of
$\mathcal E$ and $\mathcal A$ are open, while the alternatives involving
$\mathcal Z$ are $F_\sigma$; consequently $\mathcal E(K,m,n)$ and
$\mathcal A(K,m,n,k)$ are $F_\sigma$ predicates in~$K$.
For fixed $z\in\Z^{(\N)}$, membership $z\in H_K$ is equivalent to the existence of a
finite formal derivation of $z$ as a signed sum of generators of types~\textup{(1)}--\textup{(3)},
together with the corresponding activation relations.
Hence
\[
\{K\in\Ideal(A): z\in H_K\}
\]
is a countable union of finite intersections of $F_\sigma$ sets, therefore it is
$F_\sigma$.
\end{proof}

\begin{remark}
\label{rem:K0-general}
Inspection of the proof of Theorem~\ref{thm:K0} shows that it applies verbatim to
any separable \emph{unital} $C^*$-algebra $C$.
For a non-unital algebra $B$, however, one has
\[
K_0(B)=\ker\bigl(K_0(\widetilde B)\to K_0(\C)\cong \Z\bigr),
\]
so an internal treatment requires the relative/unitised picture rather than the
plain projection monoid $V(B)$.
We do not develop that extra bookkeeping here, because the present paper needs
only Borelness of $K_1$ and the higher groups.
\end{remark}

\subsection{The \texorpdfstring{$K_1$}{K1}-assignment}\label{subsec:K1}

Let $S=C_0((0,1))$ and set
\[
A_S:=A\otimes_{\max} S.
\]
Since $S$ is nuclear, this agrees canonically with $A\otimes_{\min}S$.
For $K\in\Ideal(A)$ define
\[
J_K:=\overline{K\odot S}\triangleleft A_S.
\]
Then
\[
A_S/J_K\cong (A/K)\otimes S = S(A/K).
\]

This is the first place where non-unitality matters: the suspension
$A_S/J_K\cong S(A/K)$ is non-unital, so the preceding internal $K_0$-presentation
cannot be applied verbatim.
We therefore combine the Borel suspension map $K\mapsto J_K$ with the already-known
Borel computation of $K$-theory in the standard $\Gamma/\Xi$ codings.

\begin{theorem}\label{thm:K1}
Equip $\Ideal(A)$ with the Wijsman topology.  There is a Borel map
\[
\kappa_{K_1}:\Ideal(A)\to\Sub(\Z^{(\N)})
\]
such that
\[
\Z^{(\N)}/\kappa_{K_1}(K)\cong K_1(A/K)
\qquad (K\in\Ideal(A)).
\]
\end{theorem}

\begin{proof}
Let $S=C_0((0,1))$ and set
\[
A_S:=A\otimes_{\min} S,
\]
which agrees canonically with $A\otimes_{\max} S$ because $S$ is nuclear.
For $K\in\Ideal(A)$ define
\[
J_K:=\overline{K\odot S}\triangleleft A_S.
\]
Then
\[
A_S/J_K\cong (A/K)\otimes S = S(A/K).
\]
By Proposition~\ref{prop:tensor-borel} (with $\bullet=\min$), the map
\[
K\longmapsto J_K
\]
from $\Ideal(A)$ to $\Ideal(A_S)$ is Borel.

Choose a dense sequence of contractions in $A_S$ which generates $A_S$, and apply
Lemma~\ref{lem:fixed-ideal-to-Xi} to pass Wijsman-continuously from
$\Ideal(A_S)$ to the seminorm coding $\Xi$.
In that coding, Farah--Toms--T\"ornquist proved that the $K_0$- and $K_1$-assignments
are Borel; see Remark~\ref{rem:K-FTT-compare}.
Applying their Borel $K_0$-assignment to the Borel code of the suspended quotient
$A_S/J_K$ yields a Borel code for
\[
K_0(A_S/J_K)=K_0\bigl(S(A/K)\bigr).
\]
By Bott periodicity,
\[
K_0\bigl(S(A/K)\bigr)\cong K_1(A/K),
\]
and composing these Borel maps gives the required $\kappa_{K_1}$.
\end{proof}

\begin{remark}
\label{rem:Kn-KOn}
For each $m\in\N$, let $S^{m}(B)$ denote the $m$-fold suspension of a $C^*$-algebra~$B$.
Iterating the Borel suspension assignment and composing with the standard-coding
Borel $K_0$-assignment yields a Borel map
\[
\kappa_{K_m}:\Ideal(A)\to\Sub(\Z^{(\N)})
\]
such that
\[
\Z^{(\N)}/\kappa_{K_m}(K)\cong K_0\bigl(S^{m}(A/K)\bigr).
\]
By complex Bott periodicity,
\[
K_m(A/K)\cong K_0\bigl(S^{m}(A/K)\bigr)\cong K_{m\bmod 2}(A/K).
\]
Thus all higher $K$-groups are Borel in the present quotient coding.
If one develops the relative/unitised internal $K_0$-construction alluded to in
Remark~\ref{rem:K0-general}, the same conclusion becomes entirely internal.
We do not optimise the precise Wijsman rank of $\kappa_{K_1}$.
\end{remark}

\begin{lemma}
\label{lem:Sub-map-complexity}
Let $X$ be a Polish space and let
\[
\Theta:X\to\Sub(\Z^{(\N)})
\]
be a map.
Assume that for every $z\in\Z^{(\N)}$ the section
\[
A_z:=\{x\in X:\ z\in \Theta(x)\}
\]
is $F_\sigma$.
Then:
\begin{enumerate}
\item for every finite subsets $E,F\subseteq \Z^{(\N)}$, the basic clopen cylinder
\[
U_{E,F}:=\{H\in\Sub(\Z^{(\N)}):\ E\subseteq H,\ F\cap H=\varnothing\}
\]
has preimage $\Theta^{-1}(U_{E,F})\in\Delta^0_3(X)$;
\item for every open $U\subseteq \Sub(\Z^{(\N)})$, the preimage $\Theta^{-1}(U)$ is
$\Sigma^0_3$.
\end{enumerate}
In particular, $\Theta$ is of Baire class~$2$.
\end{lemma}

\begin{proof}
For fixed finite $E,F$, the set
\[
\Theta^{-1}(U_{E,F})
=
\Bigl(\bigcap_{z\in E}A_z\Bigr)\cap
\Bigl(\bigcap_{w\in F}(X\setminus A_w)\Bigr)
\]
is the intersection of an $F_\sigma$ set with a $G_\delta$ set.
Equivalently,
\[
\Theta^{-1}(U_{E,F})
=
\Bigl(\bigcap_{z\in E}A_z\Bigr)\setminus
\Bigl(\bigcup_{w\in F}A_w\Bigr),
\]
hence it is the difference of two $F_\sigma$ sets.
In a Polish space, the difference of two $F_\sigma$ sets lies in $\Delta^0_3$:
if $A=\bigcup_n C_n$ and $B$ is $F_\sigma$, with each $C_n$ closed, then
\[
A\setminus B=\bigcup_n (C_n\setminus B),
\]
and each $C_n\setminus B$ is $G_\delta$; thus $A\setminus B\in\Sigma^0_3$.
Its complement
\[
(X\setminus A)\cup B
\]
is also $\Sigma^0_3$, so $A\setminus B\in\Pi^0_3$ as well.

Now every open subset of the zero-dimensional Polish space
$\Sub(\Z^{(\N)})$ is a countable union of such basic clopen cylinders.
Therefore the preimage of every open set is $\Sigma^0_3$.
\end{proof}

\begin{corollary}
\label{cor:K0-map-complexity}
The map $\kappa_{K_0}$ of Theorem~\ref{thm:K0} is of Baire class~$2$.
More explicitly, for each $z\in \Z^{(\N)}$,
\[
\kappa_{K_0}^{-1}(\{H:z\in H\})
\]
is $F_\sigma$, preimages of basic clopen cylinders in $\Sub(\Z^{(\N)})$ are
$\Delta^0_3$, and preimages of open sets are $\Sigma^0_3$.
\end{corollary}

\begin{proof}
Combine the last assertion of Theorem~\ref{thm:K0} with
Lemma~\ref{lem:Sub-map-complexity}.
\end{proof}

%%%%%%%%%%%%%%%%%%%%%%%%%%%%%%%%%%%%%%%%%%%%%%%%%%%%%%%%%%%%%%
\section{Tensor products and kernel maps}\label{sec:tensor-products}
%%%%%%%%%%%%%%%%%%%%%%%%%%%%%%%%%%%%%%%%%%%%%%%%%%%%%%%%%%%%%%

Let $A$ and $B$ be separable $C^*$-algebras.  For
$\bullet\in\{\max,\min\}$ write $A\otimes_\bullet B$ for the corresponding
tensor product, and equip each ideal space with the Wijsman topology.
For $I\in\Ideal(A)$ and $J\in\Ideal(B)$ define
\[
\mathcal J_\bullet(I,J)=\ker(q_I\otimes_\bullet q_J)\subseteq A\otimes_\bullet B.
\]
When $\bullet=\max$ one has
$\mathcal J_{\max}(I,J)=\overline{I\odot B+A\odot J}$.
For $\bullet=\min$, this description holds whenever one factor is exact;
see~\cite[Prop.~2.3.3]{BrownOzawa}.  The exactness assumption is used precisely here:
without it, the kernel of the quotient map for the minimal tensor product need not be
this elementary closure, and the approximation argument below would have no such
finite-tensor decomposition to start from.

\begin{proposition}\label{prop:tensor-usc}
Assume either $\bullet=\max$, or $\bullet=\min$ and one factor is exact.
Then for each $z\in A\otimes_\bullet B$ the function
\[
(I,J)\longmapsto\dist\bigl(z,\mathcal J_\bullet(I,J)\bigr)
\]
is upper semicontinuous on $\Ideal(A)\times\Ideal(B)$.
\end{proposition}

\begin{proof}
Write $d(I,J)=\dist(z,\overline{I\odot B+A\odot J})$.  Fix $\varepsilon>0$ and choose
$w\in I\odot B+A\odot J$ such that $\|z-w\|_\bullet<d(I,J)+\varepsilon$.
Write
$w=\sum_{k=1}^m(i_k\otimes b_k+a_k\otimes j_k)$
with $i_k\in I$, $j_k\in J$.
If $(I_n,J_n)\to(I,J)$ Wijsman, then $d(i_k,I_n)\to 0$ and $d(j_k,J_n)\to 0$.
Choose $i_k^{(n)}\in I_n$ and $j_k^{(n)}\in J_n$ with $\|i_k-i_k^{(n)}\|\to 0$
and $\|j_k-j_k^{(n)}\|\to 0$.  Set
\[
w_n=\sum_{k=1}^m(i_k^{(n)}\otimes b_k+a_k\otimes j_k^{(n)})\in I_n\odot B+A\odot J_n.
\]
By continuity of the tensor product bilinear map, $\|w_n-w\|_\bullet\to 0$.
Hence for large~$n$,
\[
\dist(z,\mathcal J_\bullet(I_n,J_n))\le\|z-w_n\|_\bullet
\le\|z-w\|_\bullet+\|w-w_n\|_\bullet<d(I,J)+2\varepsilon.
\]
Taking $\limsup$ and letting $\varepsilon\downarrow 0$ gives upper semicontinuity.
\end{proof}

While upper semicontinuity falls short of continuity, it already
pins down the Borel complexity of the tensor-ideal assignment.

\begin{proposition}
\label{prop:tensor-borel}
Under the hypotheses of Proposition~\ref{prop:tensor-usc}, the map
\[
\mathfrak T_\bullet:\Ideal(A)\times\Ideal(B)\to\Ideal(A\otimes_\bullet B),
\quad
(I,J)\mapsto\mathcal J_\bullet(I,J),
\]
is Borel for any admissible topology on the three ideal spaces
(since all admissible topologies share the same Borel $\sigma$-algebra;
see Remark~\ref{rem:rank-convention}).
For the \emph{Wijsman} topology one can say more: $\mathfrak T_\bullet$
is of Borel class~$2$ (\emph{i.e.}\ $\Sigma^0_2$-measurable).
More precisely, for each $z\in A\otimes_\bullet B$ and
$r\in\Q_{>0}$, the preimage
$\{(I,J):\dist(z,\mathcal J_\bullet(I,J))<r\}$
is $\Sigma^0_2$ (equivalently, $F_\sigma$) in the Wijsman topology.
This is the regularity that the complexity arguments in
Section~\ref{sec:D-stability} require.
\end{proposition}

\begin{proof}
The Wijsman topology on $\Ideal(A\otimes_\bullet B)$ is the initial topology
for the family $K\mapsto d(z,K)$, so its subbasis consists of both the sublevel sets
$\{K:\dist(z,K)<r\}$ and the superlevel sets $\{K:\dist(z,K)>r\}$.
By Proposition~\ref{prop:tensor-usc}, for each fixed~$z$ the map
$(I,J)\mapsto\dist\bigl(z,\mathcal J_\bullet(I,J)\bigr)$
is upper semicontinuous, hence the strict sublevel sets
$\{(I,J):\dist(z,\mathcal J_\bullet(I,J))<r\}$ are open (and therefore
$F_\sigma=\Sigma^0_2$ since $\Ideal(A)\times\Ideal(B)$ is metrisable).
For the superlevel sets, upper semicontinuity gives
$\{(I,J):\dist(z,\mathcal J_\bullet(I,J))\ge q\}$ closed for each~$q$, and
$\{\dist>r\}=\bigcup_{q\in\Q,\,q>r}\{\dist\ge q\}$
is therefore $F_\sigma$ ($\Sigma^0_2$).
Thus preimages of both families of Wijsman subbasic opens are $\Sigma^0_2$.
\end{proof}

\begin{theorem}
\label{thm:tensor-continuity}
Let $A$ and $B$ be separable $C^*$-algebras.
\begin{enumerate}
\item For the maximal tensor product, the map $\mathfrak T_{\max}$ is continuous
for the Wijsman topologies on all three ideal spaces.
\item If one factor, say~$A$, is exact, then $\mathfrak T_{\min}$ is also
Wijsman-continuous.
\end{enumerate}
\end{theorem}

\begin{proof}
Because the Wijsman topology is metrisable, it suffices to prove sequential continuity.

Let
\[
(I_m,J_m)\longrightarrow (I,J)
\]
in $\Ideal(A)\times \Ideal(B)$.
Set
\[
T:=\N\cup\{\infty\},
\]
where $\N$ is discrete and $\infty$ is the point at infinity.
Define
\[
I_t=
\begin{cases}
I_m,& t=m\in\N,\\
I,& t=\infty,
\end{cases}
\qquad
J_t=
\begin{cases}
J_m,& t=m\in\N,\\
J,& t=\infty,
\end{cases}
\]
and write
\[
A_t:=A/I_t,\qquad B_t:=B/J_t.
\]

For $a\in A$ define a section
\[
\widehat a(t):=a+I_t\in A_t,
\]
and similarly for $b\in B$.

\smallskip
\noindent\emph{The quotient families are continuous fields.}
For every fixed $a\in A$, the function
\[
t\longmapsto \|\widehat a(t)\|
=
\|a+I_t\|
=
\dist(a,I_t)
\]
is continuous on $T$: continuity at points of $\N$ is automatic because they are isolated,
and continuity at $\infty$ is exactly the Wijsman convergence $I_m\to I$.
By Dixmier's criterion for continuous fields of quotient algebras
\cite[Prop.~10.3.2]{Dixmier1977}, the family $\{A_t:t\in T\}$, together with the
sections generated by the $\widehat a$'s, forms a continuous field of $C^*$-algebras
over~$T$.
The same argument applies to the family $\{B_t:t\in T\}$.

\smallskip
\noindent\emph{Maximal tensor products.}
By the standard continuity theorem for fibrewise maximal tensor products of continuous
fields; see, for example, Kirchberg--Wassermann \cite{KirchbergWassermann1995}, the
fibrewise maximal tensor product is again a continuous field over $T$, whose fibre at $t$ is
\[
A_t\otimes_{\max} B_t.
\]
Therefore, for every algebraic tensor
\[
w=\sum_{r=1}^{N} a_r\odot b_r\in A\odot B,
\]
the section
\[
t\longmapsto
w_t:=
\sum_{r=1}^{N}(a_r+I_t)\otimes (b_r+J_t)
\in A_t\otimes_{\max} B_t
\]
has continuous norm:
\[
t\longmapsto \|w_t\|_{A_t\otimes_{\max} B_t}
\quad\text{is continuous.}
\]

Now let $z\in A\otimes_{\max} B$ be arbitrary.
Given $\varepsilon>0$, choose $w\in A\odot B$ with $\|z-w\|<\varepsilon$.
Then for every $t\in T$,
\[
\bigl|\|z_t\|-\|w_t\|\bigr|
\le \|z_t-w_t\|
\le \|z-w\|
<\varepsilon,
\]
so continuity of $t\mapsto \|w_t\|$ implies continuity of $t\mapsto \|z_t\|$.

Finally, for every $t\in T$ one has the canonical quotient identification
\[
(A\otimes_{\max} B)/\mathcal J_{\max}(I_t,J_t)
\cong
A_t\otimes_{\max} B_t.
\]
Hence
\[
\dist\bigl(z,\mathcal J_{\max}(I_t,J_t)\bigr)
=
\|z_t\|_{A_t\otimes_{\max} B_t}.
\]
Therefore
\[
\dist\bigl(z,\mathcal J_{\max}(I_m,J_m)\bigr)
\longrightarrow
\dist\bigl(z,\mathcal J_{\max}(I,J)\bigr)
\qquad (m\to\infty)
\]
for every $z\in A\otimes_{\max} B$, which is exactly Wijsman continuity of
$\mathfrak T_{\max}$.

\smallskip
\noindent\emph{Minimal tensor products under exactness.}
Assume that $A$ is exact.
Then each quotient $A_t=A/I_t$ is exact.
By Kirchberg--Wassermann's exactness theorem for fibrewise minimal tensor products
(see \cite[Thm.~4.5]{KirchbergWassermann1995}), the fibrewise minimal tensor product
of the above two continuous fields is a continuous field over $T$, whose fibre at $t$ is
\[
A_t\otimes_{\min} B_t.
\]
Thus, exactly as above, for every
\[
w=\sum_{r=1}^{N} a_r\odot b_r\in A\odot B
\]
the function
\[
t\longmapsto \|w_t\|_{A_t\otimes_{\min} B_t}
\]
is continuous, and by density of $A\odot B$ in $A\otimes_{\min} B$ the same holds for every
$z\in A\otimes_{\min} B$.

Because $A$ is exact, one has for every $t\in T$ the canonical quotient identification
\[
(A\otimes_{\min} B)/\mathcal J_{\min}(I_t,J_t)
\cong
A_t\otimes_{\min} B_t,
\]
so
\[
\dist\bigl(z,\mathcal J_{\min}(I_t,J_t)\bigr)
=
\|z_t\|_{A_t\otimes_{\min} B_t}.
\]
Hence
\[
\dist\bigl(z,\mathcal J_{\min}(I_m,J_m)\bigr)
\longrightarrow
\dist\bigl(z,\mathcal J_{\min}(I,J)\bigr)
\qquad (m\to\infty)
\]
for every $z\in A\otimes_{\min} B$, \emph{i.e.}\ $\mathfrak T_{\min}$ is Wijsman-continuous.
\end{proof}

\begin{remark}\label{rem:tensor-cont}
The Borelness result (Proposition~\ref{prop:tensor-borel}) is what our complexity
arguments require; Theorem~\ref{thm:tensor-continuity} gives a strengthening under
additional hypotheses.  We do \emph{not} assert Wijsman continuity of
$(I,J)\mapsto\mathcal J_\bullet(I,J)$ in full generality for the minimal tensor
product without exactness.  Lower semicontinuity of
$(I,J)\mapsto\dist(z,\mathcal J_{\min}(I,J))$ can fail, and the additional
structure needed to force continuity is subtle.

In particular, the $K_1$-assignment of Theorem~\ref{thm:K1} uses only the
\emph{Borel} tensor-ideal assignment (Proposition~\ref{prop:tensor-borel}), not
Wijsman continuity.
\end{remark}

\subsection{Projective and injective tensor products of Banach algebras}

Let $A$ and $B$ be separable Banach algebras.  For $\alpha\in\{\pi,\varepsilon\}$ write
$A\widehat\otimes_\alpha B$ for the projective ($\pi$) or injective ($\varepsilon$)
completed tensor product of the underlying Banach spaces.
The projective tensor product $A\widehat\otimes_\pi B$ is always a Banach algebra.
The injective tensor product $A\widehat\otimes_\varepsilon B$ is a Banach algebra
provided that the multiplication on $A\odot B$ extends continuously (which is the case,
for instance, when $A$ and $B$ are uniform algebras).

If $I\triangleleft A$ and $J\triangleleft B$ are closed
(two-sided) ideals, set
\[
\mathcal J_\alpha(I,J):=\overline{I\otimes B + A\otimes J}^{\ \|\cdot\|_\alpha}
\subseteq A\widehat\otimes_\alpha B.
\]
Then $\mathcal J_\pi(I,J)$ is always a closed ideal of $A\widehat\otimes_\pi B$, and
$\mathcal J_\varepsilon(I,J)$ is a closed ideal of $A\widehat\otimes_\varepsilon B$
whenever the latter is a Banach algebra.
The following topological regularity results hold strictly at the level of Banach
spaces and their closed subspaces, and hence apply to both norms without requiring
additional algebraic assumptions.

\begin{proposition}\label{prop:banach-tensor-usc}
Fix $\alpha\in\{\pi,\varepsilon\}$ and $z\in A\widehat\otimes_\alpha B$.
With the Wijsman topologies on $\Ideal(A)$ and $\Ideal(B)$, the function
\[
(I,J)\longmapsto \dist\bigl(z,\mathcal J_\alpha(I,J)\bigr)
\]
is upper semicontinuous, hence Borel.
\end{proposition}

\begin{proof}
Fix $\varepsilon>0$ and choose $w\in I\otimes B + A\otimes J$ with
$\|z-w\|_\alpha<\dist(z,\mathcal J_\alpha(I,J))+\varepsilon$.
Write $w=\sum_{k=1}^m(i_k\otimes b_k+a_k\otimes j_k)$.

If $(I_n,J_n)\to(I,J)$ Wijsman, pick $i_k^{(n)}\in I_n$ and $j_k^{(n)}\in J_n$ with
$i_k^{(n)}\to i_k$ and $j_k^{(n)}\to j_k$ in norm.
Set
\[
w_n=\sum_{k=1}^m(i_k^{(n)}\otimes b_k+a_k\otimes j_k^{(n)})
\in I_n\otimes B+A\otimes J_n.
\]
Since $(x,y)\mapsto x\otimes y$ is jointly continuous into both
$\widehat\otimes_\pi$ and $\widehat\otimes_\varepsilon$, we have
$\|w_n-w\|_\alpha\to 0$.
Hence, for large $n$,
\[
\dist(z,\mathcal J_\alpha(I_n,J_n))\le \|z-w_n\|_\alpha
<\dist(z,\mathcal J_\alpha(I,J))+2\varepsilon.
\]
Taking $\limsup$ and letting $\varepsilon\downarrow 0$ gives upper semicontinuity.
\end{proof}

\begin{proposition}\label{prop:banach-tensor-borel}
For $\alpha\in\{\pi,\varepsilon\}$ the map
\[
\mathfrak T_\alpha:\Ideal(A)\times\Ideal(B)\to \Ideal(A\widehat\otimes_\alpha B),
\qquad (I,J)\mapsto \mathcal J_\alpha(I,J),
\]
is Borel for any admissible topology on the three ideal spaces.
For the Wijsman topology it is of Baire class~$1$:
for each $z\in A\widehat\otimes_\alpha B$ and $r\in\Q_{>0}$, the preimage
$\{(I,J):\dist(z,\mathcal J_\alpha(I,J))<r\}$ is open.
\end{proposition}

\begin{proof}
The Wijsman subbasis consists of both $\{K:\dist(z,K)<r\}$ and
$\{K:\dist(z,K)>r\}$ (Lemma~\ref{lem:wijsman-subbasic}).
By Proposition~\ref{prop:banach-tensor-usc}, the function
$(I,J)\mapsto \dist\bigl(z,\mathcal J_\alpha(I,J)\bigr)$ is upper semicontinuous, so
$\{(I,J):\dist(z,\mathcal J_\alpha(I,J))<r\}$ is open.
For the superlevel direction, $\{\dist>r\}=\bigcup_{q\in\Q,\,q>r}\{\dist\ge q\}$
is $F_\sigma$ (since $\{\dist\ge q\}$ is closed for u.s.c.\ functions).
Thus preimages of both families of Wijsman subbasic opens are $F_\sigma$,
giving Borelness (indeed Baire class~$1$) of $\mathfrak T_\alpha$.
\end{proof}

%%%%%%%%%%%%%%%%%%%%%%%%%%%%%%%%%%%%%%%%%%%%%%%%%%%%%%%%%%%%%%
\section{\texorpdfstring{$D$}{D}-absorption and approximate intertwining}\label{sec:D-stability}
%%%%%%%%%%%%%%%%%%%%%%%%%%%%%%%%%%%%%%%%%%%%%%%%%%%%%%%%%%%%%%

Fix a separable unital $C^*$-algebra~$D$.  We say that a separable unital
$C^*$-algebra $B$ is \emph{$D$-absorbing} if
$B\cong B\otimes_{\min}D$.  When $D$ is strongly self-absorbing, this is the
usual notion of \emph{$D$-stability}; see \cite[II.8.5.1]{BlackadarOA}
and \cite{TomsWinter}.

\subsection{Elliott intertwining}\label{subsec:elliott}

A central theme in the classification programme for $C^*$-algebras is the passage from
approximate morphisms to genuine isomorphisms.  The fundamental device for this is the
\emph{Elliott intertwining argument}, which first appeared in
Elliott's classification of AF-algebras~\cite{Elliott1976} and has since become one of
the most widely used tools in the field.  The argument says, roughly, that if two
separable $C^*$-algebras $B$ and $C$ can be connected by a back-and-forth sequence of
maps that \emph{approximately} compose to the identity on ever-larger finite subsets,
then $B\cong C$ \emph{genuinely}.

This theorem matters for us for two reasons.  First, the
\emph{reduction from an isomorphism to finite-level witnesses} is the mechanism
that makes $D$-absorption ($B\cong B\otimes_{\min}D$) expressible as an analytic condition
(Theorem~\ref{thm:Dstab}).  Without the intertwining theorem, one
would need to quantify over an entire isomorphism---an object living in a
non-separable function space---and the resulting condition would a priori be
$\Sigma^1_1$ rather than Borel.
Second, it provides the $\Sigma^1_1$ upper bound for isomorphism of
separable $C^*$-algebras (Remark~\ref{rem:isom-analytic}): the approximate
intertwining witnesses form a countable Borel family, so the existence of a
complete intertwining is an analytic condition.

The proof of the intertwining theorem is entirely constructive (a standard
diagonal argument), and is by now well-established; see
R\o rdam--Larsen--Laustsen~\cite[Theorem~2.3.3]{RordamLarsenLaustsen}
and Elliott~\cite{Elliott1976}.
We do not reproduce it here, but state it for reference and proceed
directly to its descriptive set-theoretic consequences.

\begin{theorem}
\label{thm:elliott-intertwining}
Let $B$ and $C$ be separable unital $C^*$-algebras.  Suppose there exist sequences of
unital $*$-homomorphisms $\phi_n:B\to C$ and $\psi_n:C\to B$ such that for every
$b\in B$ and $c\in C$,
\[
\|\psi_n\phi_n(b)-b\|\to 0,
\qquad
\|\phi_{n+1}\psi_n(c)-c\|\to 0.
\]
Then there exists a unital $*$-isomorphism $\Phi:B\to C$.  \qed
\end{theorem}

\subsection{Finite-level intertwining and Borel complexity}

\begin{remark}\label{rem:isom-analytic}
In any standard coding of separable $C^*$-algebras with a named dense sequence,
the isomorphism relation is analytic.  One way to see this is to code a candidate
Elliott intertwining by a pair of sequences of finite partial maps between the
named dense sets.  The coherence conditions asserting that these finite maps are
approximately $*$-homomorphic and approximately inverse on larger and larger
initial segments form a Borel relation on the Polish space of such sequences.
By Theorem~\ref{thm:elliott-intertwining}, existence of such a coded intertwining
is equivalent to genuine isomorphism.  Thus $\cong$ is analytic.
Moreover, the resulting orbit-equivalence complexity is maximal in the standard
sense: Sabok proved that isomorphism of separable $C^*$-algebras, already of simple
separable AI algebras, is complete among orbit equivalence relations; see
\cite{SabokCompleteness}.  This is distinct from, and stronger in a different
direction than, merely observing analyticity as a subset of a product coding space.
\end{remark}

\begin{theorem}\label{thm:Dstab}
Let $D$ be a separable unital exact $C^*$-algebra.
Then the class of $K\in\Ideal(A)$ such that $A/K$ is $D$-absorbing is analytic.
\end{theorem}

\begin{proof}
Write $B_K=A/K$ and
\[
C_K:=(A/K)\otimes_{\min} D.
\]
Since $D$ is exact, Proposition~\ref{prop:tensor-borel} applies to the minimal tensor
product, so the assignment
\[
K\longmapsto \mathcal J_{\min}(K,\{0\})\in \Ideal(A\otimes_{\min} D)
\]
is Borel. Therefore the quotient assignment $K\mapsto C_K$ is Borel in the standard
coding of separable $C^*$-algebras.

By Remark~\ref{rem:isom-analytic}, the isomorphism relation for separable $C^*$-algebras
is analytic in the quotient coding. Pulling this analytic relation back along the Borel
map
$K\mapsto (B_K,C_K)$
shows that
$\{K\in\Ideal(A): B_K\cong C_K\}$
is analytic.
\end{proof}

\begin{remark}\label{rem:Dstab-SSA}
For strongly self-absorbing $D$, one expects a sharper $\Pi^0_3$ description of
$D$-stability via approximately central approximately multiplicative u.c.p.\ maps
on finite sets.
However, obtaining that bound in the present coding requires a separate
finite-dimensional operator-system lemma for coding u.c.p.\ maps from a fixed
separable domain into quotients.
Since that lemma is not developed here, we content ourselves with the analytic
upper bound of Theorem~\ref{thm:Dstab}.
\end{remark}

%%%%%%%%%%%%%%%%%%%%%%%%%%%%%%%%%%%%%%%%%%%%%%%%%%%%%%%%%%%%%%
\section{TROs and ternary ideals}\label{sec:TRO}
%%%%%%%%%%%%%%%%%%%%%%%%%%%%%%%%%%%%%%%%%%%%%%%%%%%%%%%%%%%%%%

A ternary ring of operators (TRO) is a closed subspace
$T\subseteq\mathcal B(H,K)$ closed under $[x,y,z]=xy^*z$.
A closed subspace $J\subseteq T$ is a ternary ideal if
\[
    [T,T,J]+[T,J,T]+[J,T,T]\subseteq J.
\]

\begin{lemma}\label{lem:freeTRO}
There exists a separable TRO $T_{\mathrm{free}}$ generated by a sequence
$(u_n)$ of contractions such that for every separable TRO $T$ and every
sequence $(t_n)\subseteq \operatorname{Ball}(T)$ there exists a unique
contractive ternary homomorphism
$\pi:T_{\mathrm{free}}\to T$
with $\pi(u_n)=t_n$ for all $n$.
In particular every separable TRO is a quotient of $T_{\mathrm{free}}$ by a
closed ternary ideal.
\end{lemma}

\begin{proof}
Let $\mathcal T_{\mathrm{alg}}$ be the algebraic free ternary system on symbols
$u_n$.  Define
\begin{multline*}
\|x\|_{\max}
:=
\sup\bigl\{\|\rho(x)\|:\ \rho:\mathcal T_{\mathrm{alg}}\to S
\text{ is a ternary homomorphism}\\
\text{into a TRO }S,\ 
\|\rho(u_n)\|\le 1\ \forall n\bigr\}.
\end{multline*}
An induction on ternary word length shows that $\|x\|_{\max}<\infty$ for every
$x\in \mathcal T_{\mathrm{alg}}$: if $x=[x_1,x_2,x_3]$, then
$\|\rho(x)\|\le\|\rho(x_1)\|\,\|\rho(x_2)\|\,\|\rho(x_3)\|$,
and the right-hand side is uniformly bounded over all admissible $\rho$ by the
inductive hypothesis.

Let $N=\{x:\|x\|_{\max}=0\}$, set $T_0=\mathcal T_{\mathrm{alg}}/N$, and let
$T_{\mathrm{free}}$ be the completion of $T_0$.  By construction the images of
the $u_n$ are contractions.

Given $(t_n)\subseteq \operatorname{Ball}(T)$, the algebraic universal property
gives a unique ternary homomorphism
$\rho:\mathcal T_{\mathrm{alg}}\to T$ with $\rho(u_n)=t_n$.
By definition of $\|\cdot\|_{\max}$, $\rho$ is contractive, so it factors
through $T_0$ and extends uniquely to a contractive ternary homomorphism
$\pi:T_{\mathrm{free}}\to T$.

Finally, if $T$ is separable, choose a dense sequence $(t_n)$ in
$\operatorname{Ball}(T)$.  The resulting map $\pi$ has dense range, and the
range of a ternary homomorphism is closed because such a homomorphism extends
to a $*$-homomorphism between linking $C^*$-algebras.  Hence $\pi$ is
surjective, so $T\cong T_{\mathrm{free}}/\ker\pi$.
\end{proof}

\begin{proposition}\label{prop:TRO-polish}
The space of closed ternary ideals in $T_{\mathrm{free}}$ is closed in
$\SB(T_{\mathrm{free}})$ for any admissible topology, hence Polish.
The quotient map $J\mapsto T_{\mathrm{free}}/J$ is a Borel surjection onto separable
TROs (up to isomorphism).
\end{proposition}

\begin{proof}
Apply Proposition~\ref{prop:ideal-closed} to the ternary operation
$[x,y,z]=xy^*z$.
\end{proof}

Write $V=T_{\mathrm{free}}/J$ and let $D^T\subseteq T_{\mathrm{free}}$ be a fixed
countable dense subset; set $\Phi_T(J,x)=\|x+J\|_{V}$.

\begin{proposition}\label{prop:TRO-complexity}
In the Wijsman topology on the ternary ideal space of $T_{\mathrm{free}}$:
\begin{enumerate}
\item The class of $J$ such that $[x,y,z]=[z,y,x]$ in $V=T_{\mathrm{free}}/J$ for all
$x,y,z\in V$ is closed.
\item The class of $J$ such that $V$ is rectangular AF, that is, the closure of an
increasing union of finite-dimensional sub-TROs, is $\Pi^0_2$.
\end{enumerate}
\end{proposition}

\begin{proof}
For ternary commutativity it suffices to test $x,y,z\in D^T$:
\[
        \Phi_T(J,[x,y,z]-[z,y,x])=0.
\]
Each condition is closed, and the intersection is countable.

For rectangular AF, fix a finite tuple $\bar x\in(D^T)^{<\N}$ and $n\in\N$.  A witness
consists of a finite direct sum of rectangular matrix spaces and approximate
rectangular matrix units from $D^T$, together with rational linear combinations of
these units approximating the entries of $\bar x$ within $1/n$.  The rectangular
matrix-unit relations are stable by applying the finite-dimensional perturbation
lemma in the linking $C^*$-algebra.  Hence, for fixed finite data, the witness
conditions are finitely many strict inequalities in the continuous functions
$J\mapsto\Phi_T(J,t)$, and so define an open set.  The search over all finite
rectangular types, dense tuples, and rational coefficients is countable.  Thus, for
fixed $(\bar x,n)$, the witness set is open; intersecting over all $\bar x$ and $n$
gives a $G_\delta$ set.
\end{proof}

We do not assign a finite Wijsman rank to TRO simplicity.  The dense-set argument that
works for purely algebraic simplicity has the same obstruction as in the Banach-algebra
case discussed after Proposition~\ref{prop:simple}: non-full elements need not be
approximated by non-full elements from a prescribed dense set.

%%%%%%%%%%%%%%%%%%%%%%%%%%%%%%%%%%%%%%%%%%%%%%%%%%%%%%%%%%%%%%
\section{Countable structures: a complexity catalogue}
\label{sec:countable-complexity}
%%%%%%%%%%%%%%%%%%%%%%%%%%%%%%%%%%%%%%%%%%%%%%%%%%%%%%%%%%%%%%

All spaces in this section are compact zero-dimensional Polish spaces (closed subspaces
of Cantor cubes), and atomic predicates are clopen.  For a fixed word $w$ in the free
object, the predicate ``$[w]=1$ in the quotient'' is clopen.  Consequently:
one universal quantifier over a countable domain yields a closed
($\Pi^0_1$) condition;
one existential quantifier yields an open ($\Sigma^0_1$) condition;
alternations proceed in the expected way.

\subsection{Countable groups}\label{subsec:groups}

We code $G=F_\infty/N$ with $N\in\NSub(F_\infty)\subseteq 2^{F_\infty}$.

\paragraph{Order and structure.}
\begin{itemize}
\item \emph{Finite:} $\Sigma^0_2$ ($\exists$ finite transversal
$T\subseteq F_\infty$ s.t.\ $\forall g\in F_\infty\;\bigvee_{t\in T} gt^{-1}\in N$;
each membership test is clopen, so the inner condition is closed, and the outer
existential over finite $T$ gives $\Sigma^0_2$).
\item \emph{Finitely generated:} $\Sigma^0_3$ ($\exists$ finite $S\subseteq F_\infty$
s.t.\ $\forall g\in F_\infty$, $g\in\langle S\rangle N$).
\item \emph{Cyclic:} $\Sigma^0_3$ ($\exists g\;\forall h\;\exists n:\ h\in\langle g\rangle N$).
\item \emph{Abelian:} closed ($[F_\infty,F_\infty]\subseteq N$).
\item \emph{Simple:} $\Pi^0_2$.
$G_N$ is simple iff for every $g$ and every generator $x_i$,
either $g\in N$ or there exist
$w_1,\dots,w_k\in F_\infty$ and signs $\varepsilon_j\in\{\pm1\}$ such that
$x_i(\prod_{j=1}^k w_j g^{\varepsilon_j} w_j^{-1})^{-1}\in N$.
Since $g\in N$ is clopen and the existentially quantified membership test is
a countable union of clopen conditions (hence open), each instance
of the outer $\forall g\;\forall x_i$ is open, and the intersection is $G_\delta=\Pi^0_2$.
\end{itemize}

\paragraph{Commutativity and centrality.}
\begin{itemize}
\item \emph{Nilpotent of class $\le c$:} closed (all basic commutators of weight
$c\!+\!1$ lie in~$N$).  ``Nilpotent'' $=\bigcup_c$ is $\Sigma^0_2$.
\item \emph{Solvable of derived length $\le d$:} closed.  ``Solvable'' is $\Sigma^0_2$.
\item \emph{Perfect ($G=[G,G]$):} $G_\delta$ ($\Pi^0_2$).
For each generator~$x_i$, the condition
$\exists w\in[F_\infty,F_\infty]:\ x_iw^{-1}\in N$ is a countable union of clopen sets,
hence open.
Intersecting over all generators gives $G_\delta$.
\item \emph{Centreless:} $\Pi^0_2$
($\forall g\;\exists h:\ g\in N\lor [g,h]\notin N$; the disjunction is
clopen $\cup$ clopen = clopen, so $\forall\exists(\text{clopen})$ gives $\Pi^0_2$).
\item \emph{Non-trivial centre:} $\Sigma^0_2$
($\exists g:\ g\notin N\land\forall h:\ [g,h]\in N$; the conjunction of a
clopen set with a countable intersection of clopen sets is closed,
so $\exists(\text{closed})$ gives $\Sigma^0_2$).
\end{itemize}

\paragraph{Torsion.}
\begin{itemize}
\item \emph{Torsion-free:} closed
($\forall g\;\forall n\ge 2:\ g^n\in N\implies g\in N$).
\item \emph{Periodic:} $\Pi^0_2$ ($\forall g\;\exists n\ge 2:\ g^n\in N$).
\item \emph{Locally finite:} $\Pi^0_3$ ($\forall$ finite $S$, $\langle S\rangle N/N$
is finite; each instance is $\Sigma^0_2$ since finiteness is $\Sigma^0_2$,
and $\forall(\Sigma^0_2)=\Pi^0_3$).
\end{itemize}

\paragraph{Group extensions and virtual properties.}
The remaining properties in this direction
(finite-by-cyclic, finite-by-abelian, polycyclic, cyclic-by-finite, virtually
cyclic, virtually abelian, virtually nilpotent, residually finite, Hopfian, and
related variants) require additional coding of finite quotients, finite-index
subgroups, or endomorphisms.  Their precise Borel complexity is not established
in the present paper, so we do not assign ranks to them here.

We record only the easy derived-series case:
\begin{itemize}
\item \emph{Metabelian:} closed, since $G_N$ is metabelian iff
$F_\infty^{(2)}\subseteq N$.
\end{itemize}

\subsubsection{Universality and embeddability}

A fundamental asymmetry between the countable-group and $C^*$-algebraic settings
is the failure of universality for embeddings.
There is no \emph{universal} countable group with respect to embeddings: no single
countable group contains an isomorphic copy of every countable group.
Indeed, there are $2^{\aleph_0}$ pairwise non-isomorphic finitely generated groups
(already among $2$-generated quotients of the free group $F_2$), whereas any
fixed countable group has only countably many finitely generated subgroups
(since each is generated by a finite subset).
Thus no countable group can contain copies of all finitely generated groups,
and \emph{a fortiori} none can contain copies of all countable groups;
see, e.g.,~\cite[Chapter~IV]{LyndonSchupp} for related standard facts.
In the quotient coding this means that the parameter space $\NSub(F_\infty)$ cannot
be replaced by the ideal space of a single countable group, in contrast with the
$C^*$-algebraic setting where $C^*_{\max}(F_\infty)$ is universal for separable
$C^*$-algebras.
Nevertheless, $F_\infty$ is universal for \emph{presentations}: every countable group
is a quotient of~$F_\infty$, and the normal subgroup lattice $\NSub(F_\infty)$
serves as the standard Borel space for the classification problem.

Gromov~\cite{Gromov1999} initiated a systematic study of the \emph{space of
finitely generated groups} (quotients of $F_k$ for fixed~$k$, equipped with the
Chabauty topology), viewing group-theoretic properties as subsets of this space
and asking for their topological and measure-theoretic complexity.
The present catalogue may be seen as an extension of this viewpoint to all
countable groups (using $F_\infty$ in place of~$F_k$) within the framework
of the Borel hierarchy.

\subsubsection{Sofic groups}

Sofic groups were introduced by Gromov~\cite{Gromov1999} as a common generalisation of amenable
and residually finite groups.  They form a robust class of groups that are
well-approximated by finite symmetric groups.  The motivation for studying sofic groups
stems from the fact that they are known to satisfy several major conjectures that remain
open for arbitrary groups, such as Gottschalk's surjunctivity conjecture and
Kaplansky's direct finiteness conjecture for group rings.  Gromov~\cite{Gromov1999}
asked whether \emph{every} group is sofic; this remains one of the most prominent open
problems in group theory.

\begin{definition}\label{def:sofic}
A countable group $G$ is \emph{sofic} if for every finite subset $F\subseteq G$ and
every $\varepsilon>0$, there exist a natural number $d$ and a map
$\sigma:F\to\mathrm{Sym}(d)$ (not necessarily a homomorphism) such that:
\begin{enumerate}
\item $\sigma(1)=\mathrm{id}$ (if $1\in F$);
\item for all $u,v,w\in F$ with $uv=w$ in $G$, the normalised Hamming distance
satisfies $d_H(\sigma(u)\sigma(v),\sigma(w))<\varepsilon$;
\item for all distinct $u,v\in F$,
$d_H(\sigma(u),\sigma(v))>1-\varepsilon$.
\end{enumerate}
\end{definition}

\begin{proposition}\label{prop:sofic}
The class of $N\in\NSub(F_\infty)$ such that $F_\infty/N$ is sofic is $G_\delta$
(\emph{i.e.}~$\Pi^0_2$).
\end{proposition}

\begin{proof}
Let $G_N=F_\infty/N$.  Fix a finite set $S\subseteq F_\infty$ containing $1$ and closed
under inverses.  Fix $k\in\N$ and write $\varepsilon=1/k$.  For $d\in\N$ and a map
$\sigma:S\to\mathrm{Sym}(d)$, we say that $\sigma$ is an \emph{$(S,\varepsilon)$-sofic
approximation for $G_N$} if the following hold:
\begin{enumerate}[label=\textup{(\alph*)}]
\item $\sigma(1)=\mathrm{id}$;
\item (\emph{multiplicativity on the partial multiplication table})
for all $u,v,w\in S$ with $uvw^{-1}\in N$ (\emph{i.e.}\ $uv=w$ in $G_N$), we have
$d_H(\sigma(u)\sigma(v),\sigma(w))<\varepsilon$;
\item (\emph{almost injectivity})
for all $u,v\in S$ with $uv^{-1}\notin N$ (\emph{i.e.}\ $u\ne v$ in $G_N$), we have
$d_H(\sigma(u),\sigma(v))>1-\varepsilon$.
\end{enumerate}
(Here $d_H$ denotes the normalised Hamming metric on $\mathrm{Sym}(d)$.)

For fixed $(S,k,d,\sigma)$, the set of $N$ for which $\sigma$ is an
$(S,1/k)$-sofic approximation depends only on membership of the finitely many words
\[
\{u:\ u\in S\}\ \cup\ \{uvw^{-1}:\ u,v,w\in S\}\ \cup\ \{uv^{-1}:\ u,v\in S\}
\]
in $N$, together with finitely many Hamming-metric inequalities involving the fixed
permutations $\sigma(u)$.  Since each predicate ``$w\in N$'' is clopen in the Cantor
topology on $2^{F_\infty}$, this set of $N$ is clopen.

Let $U_{S,k}$ be the union of these clopen sets over all $d\in\N$ and all
$\sigma:S\to\mathrm{Sym}(d)$.  Then $U_{S,k}$ is open.

Finally, $G_N$ is sofic iff for every finite $S$ and every $k$ there exists such a
$(S,1/k)$-approximation, \emph{i.e.}\ if and only if
$N\in \bigcap_{S,k} U_{S,k}$,
a countable intersection of open sets.  Hence the sofic codes form a $G_\delta$ set.
\end{proof}

\subsection{Countable abelian groups}\label{subsec:abelian}

We code $A=\Z^{(\N)}/H$ with $H\in\Sub(\Z^{(\N)})$.

\paragraph{Basic structure.}
\begin{itemize}
\item \emph{Finite:} $\Sigma^0_2$
($\exists$ finite transversal $T$ s.t.\ $\forall x\;\bigvee_{t\in T} x-t\in H$).
\item \emph{Finitely generated:} $\Sigma^0_3$
($\exists$ finite $S$ s.t.\ $\forall x\;\exists\vec n:\ x-\sum n_ss\in H$).
\item \emph{Cyclic:} $\Sigma^0_3$
($\exists g\;\forall x\;\exists n:\ x-ng\in H$).
\item \emph{Torsion-free:} closed
($\forall x\;\forall n\ge 2:\ nx\in H\implies x\in H$;
each instance is clopen, so the intersection is closed).
\item \emph{Torsion:} $\Pi^0_2$
($\forall x\;\exists n\ge 1:\ nx\in H$).
\item \emph{Divisible:} $\Pi^0_2$
($\forall a\;\forall n\ge 1\;\exists b:\ nb-a\in H$).
\item \emph{First Ulm subgroup is zero ($G^1=0$):} $\Pi^0_3$.
An abelian group $A_H=\Z^{(\N)}/H$ has $G^1=0$ (\emph{i.e.}, has no non-zero element divisible by every positive integer) iff
\[
\forall x\in \Z^{(\N)}\ 
\Bigl(
x\in H\ \lor\ \exists n\ge 1\ \forall y\in \Z^{(\N)}:\ ny-x\notin H
\Bigr).
\]
For fixed $x$ and $n$, the inner universal condition is closed; hence the
displayed formula has complexity $\Pi^0_3$.
(For torsion-free groups, being reduced is equivalent to $G^1=0$.
We do not need, and therefore do not discuss, the descriptive complexity of
reducedness for general countable abelian groups here.)
\end{itemize}

\paragraph{Torsion decomposition.}
\begin{itemize}
\item \emph{$p$-primary:} $\Pi^0_2$ ($\forall x\;\exists k:\ p^kx\in H$).
\item \emph{Direct sum of cyclic groups:} analytic ($\Sigma^1_1$) in general, since it
quantifies over an entire sequence of potential cyclic generators.
\end{itemize}

\paragraph{As $\Z$-modules.}
\begin{itemize}
\item \emph{Free abelian of finite rank $r$:} $\Sigma^0_3$
($\exists$ independent set of size $r$ generating modulo~$H$).
\item \emph{Injective $\Leftrightarrow$ divisible:} $\Pi^0_2$.
\end{itemize}

\paragraph{Slenderness.}
A torsion-free abelian group $A$ is \emph{slender} if every homomorphism
$\Z^\N\to A$ depends on finitely many coordinates.

\begin{proposition}\label{prop:slender-Pi03}
For countable abelian groups,
\[
\mathrm{Slender}=\mathrm{TorsionFree}\cap \{G : G^1=0\}.
\]
Consequently, in the coding by subgroups of $\Z^{(\N)}$, the class of slender groups is
$\Pi^0_3$.
\end{proposition}

\begin{proof}
By Nunke's theorem~\cite{Nunke1962}, an abelian group $G$ is slender iff it contains no subgroup
isomorphic to any of $\Q$, $\Z^\omega$, $\Z(p)$, or $J_p$ (for any prime $p$).
If $G$ is countable, then $\Z^\omega$ and $J_p$ cannot embed into $G$ because they are
uncountable. The obstruction $\Z(p)$ is exactly torsion, so for countable groups the
criterion reduces to: $G$ is slender iff $G$ is torsion-free and contains no copy of
$\Q$.

For torsion-free abelian groups, ``contains no copy of $\Q$'' is equivalent to $G^1=0$.
Indeed, every copy of $\Q$ is divisible, so a group containing it has non-zero elements
divisible by every integer.
Conversely, if a torsion-free group $G$ has a non-zero element $x$ divisible by every
integer, the unique roots of $x$ generate a subgroup isomorphic to $\Q$.

Thus, for countable abelian groups,
\[
\mathrm{Slender}=\mathrm{TorsionFree}\cap \{G : G^1=0\}.
\]
By the complexity bounds recorded above, torsion-free is closed ($\Pi^0_1$) and $G^1=0$
is $\Pi^0_3$. Hence slenderness is $\Pi^0_3$.
\end{proof}

\begin{question}\label{qu:slender-optimal}
Is the $\Pi^0_3$ upper bound for slender countable abelian groups optimal?
In particular, is slenderness $\Pi^0_3$-complete in this coding?
\end{question}

\subsection{Countable rings}\label{subsec:rings}

We code $R=R_\infty/I$ where $R_\infty=\Z\langle X_n:n\in\N\rangle$ and
$I\in\Id(R_\infty)$.

\paragraph{Basic algebraic structure.}
\begin{itemize}
\item \emph{Commutative:} closed ($\forall a,b:\ ab-ba\in I$).
\item \emph{Finite:} $\Sigma^0_2$ (finite transversal).
\item \emph{Integral domain (commutative case):} closed
($\forall a,b:\ ab\in I\implies a\in I\lor b\in I$).
\item \emph{Division ring:} $\Pi^0_2$
($\forall a\;\exists b:\ a\in I\lor(ab-1\in I\text{ and }ba-1\in I)$;
the first disjunct is clopen and the second is a countable union of clopen conditions,
so each instance is open and the intersection is $G_\delta$).
\item \emph{Field:} $\Pi^0_2$ (commutative $\cap$ division ring; intersection of closed and $\Pi^0_2$).
\item \emph{Reduced:} closed
($\forall a\;\forall n\ge 1:\ a^n\in I\implies a\in I$).
\item \emph{Nil ring:} $\Pi^0_2$ ($\forall a\;\exists n\ge 1:\ a^n\in I$).
\item \emph{Boolean ring:} closed ($\forall a:\ a^2-a\in I$).
\end{itemize}

\paragraph{Ideals and factor structure.}
\begin{itemize}
\item \emph{Simple:} $\Pi^0_2$
($\forall a\;\exists m,b_j,c_j:\ a\in I\lor\sum_{j=1}^mb_jac_j-1\in I$;
$a\in I$ is clopen, the existential clause is a countable union of clopen conditions,
hence open; $\forall(\text{open})=G_\delta$).
\item \emph{Prime (two-sided):} $\Pi^0_2$
($\forall a,b\;\exists r:\ a\in I\lor b\in I\lor arb\notin I$;
each disjunct is clopen, so the disjunction is open and $\forall(\text{open})=G_\delta$).
\item \emph{Semiprime:} $\Pi^0_2$
($\forall a:\ (\forall r:\ ara\in I)\implies a\in I$).
\item \emph{Local:} $\Pi^0_2$
($\forall a:\ [\exists u:\ au-1,ua-1\in I]\lor[\exists v:\ (1\!-\!a)v-1,v(1\!-\!a)-1\in I]$;
each $\exists$ clause is open, so the disjunction is open and $\forall(\text{open})=G_\delta$).
\item \emph{Von Neumann regular:} $\Pi^0_2$
($\forall a\;\exists x:\ a-axa\in I$; $\forall\exists(\text{clopen})=G_\delta$).
\item \emph{Unit-regular:} $\Pi^0_2$
($\forall a\;\exists u,v:\ uv-1,vu-1,a-aua\in I$;
$\forall\exists(\text{clopen})=G_\delta$).
\end{itemize}

\paragraph{Dedekind finiteness.}

\begin{proposition}\label{prop:ring-df}
Let $R=R_\infty/I$ be coded by $I\in\Id(R_\infty)$.  The class of codes $I$ such that
$R$ is \emph{Dedekind finite} (directly finite), \emph{i.e.}\
$\forall a,b\in R\ (ab=1\ \Rightarrow\ ba=1)$,
is a closed subset of $\Id(R_\infty)$ (hence $\Pi^0_1$).
Its complement (Dedekind infiniteness) is open.
\end{proposition}

\begin{proof}
Fix $u,v\in R_\infty$.  In the quotient $R_\infty/I$ we have $uv=1$ iff $uv-1\in I$,
and $vu=1$ iff $vu-1\in I$.  For fixed $(u,v)$ the set
\[
E_{u,v}:=\{I:\ (uv-1\in I)\Rightarrow(vu-1\in I)\}
\]
is the complement of the clopen cylinder
$\{I:\ uv-1\in I\ \wedge\ vu-1\notin I\}$, hence is clopen.
Dedekind finiteness is $\bigcap_{u,v\in R_\infty}E_{u,v}$, a countable intersection
of clopen sets, therefore closed.
The complement is the corresponding countable union of clopen sets, hence open.
\end{proof}

\begin{remark}\label{rem:ring-df-contrast}
The contrast with the Banach-algebra case (Proposition~\ref{prop:df-banach}) is
instructive: in countable rings the Dedekind finiteness condition involves no norm
bounds and the atomic predicates are clopen, yielding a closed ($\Pi^0_1$) set.
In Banach algebras the atomic predicates involve continuous (not clopen) norms, but
the open-fragment/Neumann-series argument above still yields a closed set.
\end{remark}

\subsection{Lattices and Boolean algebras}\label{subsec:lattices}

We code lattices as quotients $L_\infty/\theta$ with $\theta$ a lattice congruence
on the free lattice $L_\infty$ in the signature $(\wedge,\vee)$.  If considering
bounded lattices, we use $(\wedge,\vee,0,1)$.

\paragraph{Lattice properties.}

\begin{itemize}
\item \emph{Finite:} $\Sigma^0_2$ (finite transversal).
\item \emph{Distributive:} closed (satisfies the identity
$x\wedge(y\vee z)=(x\wedge y)\vee(x\wedge z)$).
\item \emph{Modular:} closed (satisfies the Horn condition
$x\le b\implies x\vee(a\wedge b)=(x\vee a)\wedge b$).
\item \emph{Complemented (bounded case):} $\Pi^0_2$.
For a fixed $x$, the condition $\exists y:\ (x\wedge y=0)\text{ and }(x\vee y=1)$
is open (a single existential over the countable domain).  Intersecting over all~$x$
gives~$\Pi^0_2$.
\end{itemize}

\paragraph{Boolean algebras.}
If we use the Boolean-algebra signature $(\wedge,\vee,{}',0,1)$, Boolean algebras form
a variety and hence a closed subspace of the corresponding congruence space on
$B_\infty$.  If we use only the lattice signature $(\wedge,\vee,0,1)$, the property of
being a Boolean algebra is $\Pi^0_2$ (distributive, which is closed, and complemented,
which is $\Pi^0_2$).  Working within the closed space of Boolean algebras $\Id(B_\infty)$:

\begin{itemize}
\item \emph{Atomless:} $\Pi^0_2$
($\forall x>0\;\exists y:\ 0<y<x$).
\item \emph{Atomic:} $\Pi^0_3$
($\forall x>0\;\exists y\le x:\ y\text{ is an atom}$, where ``$y$ is an atom'' is the
closed condition $y>0\wedge\forall z\le y\ (z=0\vee z=y)$; the inner $\forall\exists$
scheme yields~$\Pi^0_3$).
\item \emph{Complete:} $\Sigma^0_2$.  A countable Boolean algebra is complete if and only
if it is finite, so this reduces to the property of being finite.
\end{itemize}

\subsection{Optimality of Borel bounds}
\label{subsec:optimality}

The bounds established above are not merely upper estimates: several fundamental
properties are \emph{complete} for their respective Borel classes, demonstrating
that the definability method assigns optimal ranks.

\begin{theorem}\label{thm:finite-sigma2-complete}
In the parameter space $\NSub(F_\infty)$, the property of being
a finite group is $\Sigma^0_2$-complete.  The same holds for rings, abelian groups,
and lattices in their respective parameter spaces.
\end{theorem}

\begin{proof}
Membership in $\Sigma^0_2$ was noted in Section~\ref{subsec:groups}.  For completeness,
we reduce from the canonical $\Sigma^0_2$-complete set
$\mathrm{Fin} = \{ \alpha \in 2^\N : \exists N\; \forall n \ge N\; (\alpha(n) = 0) \}$.

Define $V:=\bigoplus_{k\in\N}\Z/2\Z$ with canonical basis $(e_k)_{k\in\N}$.
For $\alpha\in 2^\N$ set
\[
k_\alpha(n):=\sum_{i<n}\alpha(i)\in\N,
\]
and define a homomorphism $\pi_\alpha:F_\infty\to V$ by $\pi_\alpha(x_n)=e_{k_\alpha(n)}$.
Let 
\[
    N_\alpha:=\ker(\pi_\alpha)\in\NSub(F_\infty).
\]
For the abelian-group coding, define the surjective homomorphism
\[
\widetilde\pi_\alpha:\Z^{(\N)}\to V,
\qquad
\widetilde\pi_\alpha(e_n)=e_{k_\alpha(n)},
\]
and let
\[
H_\alpha:=\ker(\widetilde\pi_\alpha)\in \Sub(\Z^{(\N)}).
\]
Then
\[
\Z^{(\N)}/H_\alpha\cong \bigoplus_{j\in\operatorname{im}(k_\alpha)} \Z/2\Z,
\]
which is finite iff $\alpha\in\mathrm{Fin}$ and countably infinite otherwise.
Exactly as in the free-group coding, membership of a fixed element of $\Z^{(\N)}$ in
$H_\alpha$ depends only on finitely many values of $\alpha$, so the map
$\alpha\mapsto H_\alpha$ is continuous.

Then $G_\alpha:=F_\infty/N_\alpha$ is abelian of exponent dividing $2$, and in $G_\alpha$
we have $x_n=x_{n+1}$ whenever $\alpha(n)=0$ (since then $k_\alpha(n)=k_\alpha(n+1)$).
Thus $G_\alpha$ has one $\Z/2\Z$-generator for each $\sim_\alpha$-equivalence class
of $\N$ generated by the edges $n\sim_\alpha n{+}1$ when $\alpha(n)=0$.

If $\alpha\in\mathrm{Fin}$, then $\alpha(n)=0$ for all $n\ge N$, so $k_\alpha(n)$ is bounded
and $G_\alpha$ is generated by finitely many elements of order $2$, hence is finite.
If $\alpha\notin\mathrm{Fin}$, then $k_\alpha(n)\to\infty$, so $G_\alpha\cong\bigoplus_\N\Z/2\Z$
and is infinite.

Finally, the map $\alpha\mapsto N_\alpha$ is continuous.
Indeed, for a fixed reduced word $w\in F_\infty$ let $M$ be the largest index of a generator
appearing in $w$.  Then $\pi_\alpha(w)$ depends only on the values $\alpha(0),\dots,\alpha(M{-}1)$,
hence the set $\{\alpha: w\in N_\alpha\}=\{\alpha:\pi_\alpha(w)=0\}$ is clopen in $2^\N$.
Therefore $\alpha\mapsto N_\alpha$ is continuous as a map into $2^{F_\infty}$.

For rings, let
\[
V_\alpha:=\bigoplus_{j\in\operatorname{im}(k_\alpha)}\mathbb F_2 e_j,
\qquad
R_\alpha:=\mathbb F_2\oplus V_\alpha,
\]
with multiplication
\[
(a,v)(b,w):=(ab,aw+bv),
\]
so that $V_\alpha^2=0$.
Define the unique unital ring homomorphism
\[
\rho_\alpha:R_\infty\to R_\alpha
\]
by
\[
\rho_\alpha(1)=(1,0),
\qquad
\rho_\alpha(X_n)=(0,e_{k_\alpha(n)}) \quad (n\in\N).
\]
This map is surjective, since $(1,0)$ is the unit of $R_\alpha$ and the elements
$(0,e_{k_\alpha(n)})$ generate the square-zero ideal $V_\alpha$.
Then $R_\alpha$ is finite iff $\operatorname{im}(k_\alpha)$ is finite, \emph{i.e.}\ if and only if
$\alpha\in\mathrm{Fin}$; otherwise $R_\alpha$ is countably infinite.

For lattices, let $I_\alpha:=\operatorname{im}(k_\alpha)$ and let $L_\alpha$ be
the lattice consisting of all finite subsets of $I_\alpha$ together with the top
element $I_\alpha$, ordered by inclusion.  This lattice is finite iff
$\alpha\in\mathrm{Fin}$ and countably infinite otherwise.  It is generated by the
top element together with the singletons $\{k_\alpha(n)\}$, so there is a
surjective lattice homomorphism
$\lambda_\alpha:L_\infty\to L_\alpha$
sending one chosen free generator to $I_\alpha$ and the remaining generators to
the singletons $\{k_\alpha(n)\}$.  Let $\theta_\alpha=\ker(\lambda_\alpha)$.

In each case, membership of a fixed word / polynomial / lattice term in the
corresponding kernel depends only on finitely many values
$k_\alpha(0),\dots,k_\alpha(M)$, hence only on finitely many coordinates of
$\alpha$.  Therefore the maps
$\alpha\mapsto N_\alpha$, $\alpha\mapsto H_\alpha$,
$\alpha\mapsto \ker(\rho_\alpha)$, $\alpha\mapsto \theta_\alpha$
are continuous, completing the proof.
\end{proof}

\begin{theorem}\label{thm:torsion-divisible-pi2-complete}
In the parameter space $\Sub(\Z^{(\N)})$, the properties
of being a torsion group and of being a divisible group are
both $\Pi^0_2$-complete.
\end{theorem}

\begin{proof}
Membership in $\Pi^0_2$ is recorded in Section~\ref{subsec:abelian}.  We reduce from
the canonical $\Pi^0_2$-complete set
\[
        P_\infty=\{\alpha\in2^\N: \exists^\infty n\, \alpha(n)=1\}.
\]

\emph{Torsion.}
For $\alpha\in2^\N$, let $H_\alpha\le\Z^{(\N)}$ be generated by
\[
        \{2^{m-n}e_n:n\le m\text{ and }\alpha(m)=1\}.
\]
Then $\Z^{(\N)}/H_\alpha$ is torsion iff $\alpha\in P_\infty$.  Indeed, if
$\alpha$ has infinitely many ones, then for every $n$ there is $m\ge n$ with
$\alpha(m)=1$, so a power of $2$ annihilates the image of $e_n$; every element has
finite support, hence is torsion.  If $\alpha$ has only finitely many ones and $n$ is
larger than the last one, then no non-zero multiple of $e_n$ belongs to $H_\alpha$, so
the image of $e_n$ has infinite order.

The map $\alpha\mapsto H_\alpha$ is continuous.  For a fixed
$z=\sum c_ne_n\in\Z^{(\N)}$ with finite support, membership $z\in H_\alpha$ holds
iff, for every $n$ with $c_n\ne0$, there is some
\[
        m\in\{n,n+1,\dots,n+\nu_2(|c_n|)\}
\]
with $\alpha(m)=1$; here $\nu_2(q)$ is the exponent of $2$ in $q$.  Thus the cylinder
$\{\alpha:z\in H_\alpha\}$ depends on only finitely many coordinates of
$\alpha$.

\emph{Divisible.}
For $\alpha\in2^\N$, define
\[
        d_n(\alpha)=
        \begin{cases}
        1, & \alpha(n)=0,\\
        (n+2)!, & \alpha(n)=1,
        \end{cases}
\]
and let $J_\alpha\le\Z^{(\N)}$ be generated by
\[
        e_n-d_n(\alpha)e_{n+1}\qquad(n\in\N).
\]
Equivalently, $\Z^{(\N)}/J_\alpha$ is the direct limit of copies of $\Z$ with
bonding maps multiplication by $d_n(\alpha)$.

If $\alpha\in P_\infty$, then the indices with $\alpha(n)=1$ are unbounded.  Given
an element represented by $ae_i$ and an integer $q\ge1$, choose $t>i$ such that
$\prod_{k=i}^{t-1}d_k(\alpha)$ is divisible by $q$; this is possible because one of
the factorial factors can be chosen large enough.  Since
\[
        e_i=\Bigl(\prod_{k=i}^{t-1}d_k(\alpha)\Bigr)e_t
\]
in the quotient, $ae_i$ is divisible by $q$.  Thus the whole quotient is divisible.
If $\alpha\notin P_\infty$, then all but finitely many $d_n(\alpha)$ are equal to
$1$, and the quotient embeds as the cyclic subgroup $M^{-1}\Z\subseteq\Q$
for some integer $M$; this group is not divisible.

The map $\alpha\mapsto J_\alpha$ is continuous.  Put
$D_0=1$ and $D_i=\prod_{k<i}d_k(\alpha)$.  The homomorphism
\[
        e_i\longmapsto 1/D_i
\]
from $\Z^{(\N)}$ into $\Q$ has kernel exactly $J_\alpha$.  Hence, for a fixed
$z=\sum_{i\le r}c_ie_i$, the condition $z\in J_\alpha$ is
\[
        \sum_{i\le r} c_i/D_i=0,
\]
which depends only on $\alpha(0),\dots,\alpha(r-1)$.  This proves continuity and
therefore $\Pi^0_2$-hardness for both properties.
\end{proof}

%%%%%%%%%%%%%%%%%%%%%%%%%%%%%%%%%%%%%%%%%%%%%%%%%%%%%%%%%%%%%%
\section{A \texorpdfstring{$\Pi^1_1$}{Pi11}-complete property in a quotient coding}\label{sec:Pi11-example}
%%%%%%%%%%%%%%%%%%%%%%%%%%%%%%%%%%%%%%%%%%%%%%%%%%%%%%%%%%%%%%

We give an explicit non-Borel example within the quotient paradigm, showing that some
natural Banach-space properties of quotients cannot be located within the Borel
hierarchy.  This calibrates the reach of the definability method developed in earlier
sections.

\subsection{Closed subsets of Cantor space and commutative quotients}

Let $K_0=2^\N$ be the Cantor space and let $\mathscr F(K_0)$ be the hyperspace of
non-empty closed subsets of $K_0$ with its Effros Borel structure (equivalently, the
standard Borel structure induced by the Vietoris topology).

For $F\in\mathscr F(K_0)$ define the closed ideal
\[
I_F=\{f\in C(K_0): f|_F=0\}\triangleleft C(K_0).
\]
Then the restriction map $C(K_0)\to C(F)$ is a surjective $*$-homomorphism with kernel
$I_F$, hence $C(K_0)/I_F\cong C(F)$ as $C^*$-algebras.

\begin{lemma}\label{lem:F-to-I-borel}
The map $F\mapsto I_F$ from $\mathscr F(K_0)$ onto $\Ideal(C(K_0))$ is a
Borel isomorphism.
\end{lemma}

\begin{proof}
For $f\in C(K_0)$,
\[
\dist(f,I_F)=\|f|_F\|_\infty=\sup_{x\in F}|f(x)|.
\]
Hence $F\mapsto I_F$ is Borel.  The map is bijective by the ideal theory of
commutative $C^*$-algebras.

To see that the inverse is Borel, fix a countable clopen basis
$\mathcal U$ of $K_0$.  For each $U\in\mathcal U$, the indicator $1_U$ is a
continuous function on $K_0$, and
\[
F\cap U\neq\varnothing
\iff
\sup_{x\in F}1_U(x)=1
\iff
\dist(1_U,I_F)=1.
\]
Since the hit sets $\{F:F\cap U\neq\varnothing\}$ generate the Vietoris
Borel structure on $\mathscr F(K_0)$, the inverse map is Borel.
\end{proof}

\subsection{Countability of closed sets is \texorpdfstring{$\Pi^1_1$}{Pi11}-complete}

Let $\CTbl\subseteq\mathscr F(K_0)$ be the set of countable closed subsets
of~$K_0$.

\begin{theorem}\label{thm:countable-closed-Pi11}
$\CTbl$ is $\Pi^1_1$-complete.
\end{theorem}

\begin{proof}
It is classical that $\CTbl$ is coanalytic: a closed set $F\subseteq K_0$ is
uncountable iff it contains a perfect subset.

$\Pi^1_1$-hardness goes back to Hurewicz; see
\cite[Theorem~27.5 and Section~33.B]{Kechris}.  In particular, for every
uncountable Polish space $X$, the classes of countable compact subsets
$K_{\mathrm{No}}(X)$ and countable closed subsets $F_{\mathrm{No}}(X)$ are
$\Pi^1_1$-complete.  Since $K_0=2^\N$ is compact, every closed subset of $K_0$
is compact, and the Wijsman/Effros/Vietoris Borel structures coincide on
$\mathscr F(K_0)$.  Hence $\CTbl$ is $\Pi^1_1$-complete.
\end{proof}

\subsection{Separable dual of \texorpdfstring{$C(F)$}{C(F)}}

\begin{lemma}\label{lem:CF-dual-separable}
For $F\in\mathscr F(K_0)$, the dual Banach space $C(F)^*$ is separable iff $F$ is
countable.
\end{lemma}

\begin{proof}
If $F$ is countable, then every regular Borel measure on $F$ is atomic and
$C(F)^*\cong \ell_1(F)$, which is separable.

If $F$ is uncountable, then the family of point masses $(\delta_x)_{x\in F}\subseteq C(F)^*$
is uncountable and satisfies $\|\delta_x-\delta_y\|=2$ for $x\neq y$.
Hence $C(F)^*$ contains an uncountable $2$-separated subset and is non-separable.
\end{proof}

\begin{theorem}\label{thm:dual-separable-Pi11}
In the quotient coding $F\mapsto C(K_0)/I_F\cong C(F)$, the property
\[
\{F\in\mathscr F(K_0):\ C(F)^*\ \text{is separable}\}
\]
is $\Pi^1_1$-complete.  Consequently, the corresponding class of ideals
\[
\{I\in\Ideal(C(K_0)):\ C(K_0)/I\ \text{has separable dual}\}
\]
is $\Pi^1_1$-complete.
\end{theorem}

\begin{proof}
By Lemma~\ref{lem:CF-dual-separable}, the set in question is exactly $\CTbl$,
which is $\Pi^1_1$-complete by Theorem~\ref{thm:countable-closed-Pi11}.
Since $F\mapsto I_F$ is a Borel isomorphism (Lemma~\ref{lem:F-to-I-borel}),
$\Pi^1_1$-completeness transfers to the ideal coding.
\end{proof}

\begin{remark}\label{rem:Pi11-calibration}
Theorem~\ref{thm:dual-separable-Pi11} shows that certain natural properties of
quotient algebras are \emph{provably not Borel}, even for commutative $C^*$-algebras.
This stands in contrast to the properties in our catalogue
(Sections~\ref{sec:banach-complexity}--\ref{sec:Cstar-properties}), all of which are
Borel at finite levels.  The distinction reflects a fundamental boundary: properties
whose definitions involve quantification over uncountable sets (here, over all
regular Borel measures, or equivalently over all closed subsets of~$F$) tend to escape
the Borel hierarchy, whereas those expressible by countable quantification over
a dense set of witnesses remain Borel.
\end{remark}

\subsection{Superatomic Boolean algebras}\label{subsec:superatomic}

The parameterisation of countable Boolean algebras provides another application
of the $\CTbl$ space.  Let $B_\infty$ be the free Boolean algebra on countably many
generators.  The Stone space of $B_\infty$ is canonically homeomorphic to the Cantor
space $K_0=2^\N$.

By Stone duality, there is a canonical bijection between ideals
$I\in\Id(B_\infty)$ and closed subsets $F\in\mathscr F(K_0)$: the quotient Boolean
algebra $B_\infty/I$ is isomorphic to the algebra of clopen subsets of~$F$.

A Boolean algebra is called \emph{superatomic} if every homomorphic image contains
an atom.  A classical result in topology and logic establishes that a countable
Boolean algebra is superatomic if and only if its Stone space is countable.

\begin{theorem}\label{thm:superatomic-Pi11}
In the parameter space of countable Boolean algebras coded by ideals in the free
Boolean algebra $\Id(B_\infty)$, the property
\[
\{I\in\Id(B_\infty):\ B_\infty/I\ \text{is superatomic}\}
\]
is $\Pi^1_1$-complete.
\end{theorem}

\begin{proof}
Under the Stone duality correspondence $I\leftrightarrow F$, the map
$F\mapsto I_F$ from $\mathscr F(K_0)$ to $\Id(B_\infty)$ is a Borel isomorphism.
By the characterisation above, $B_\infty/I_F$ is superatomic if and only if
$F\in\CTbl$.

Since $\CTbl$ is $\Pi^1_1$-complete
(Theorem~\ref{thm:countable-closed-Pi11}), the preimage of $\CTbl$ under this
Borel isomorphism is exactly the set of ideals generating superatomic quotients.
Therefore superatomicity is $\Pi^1_1$-complete.
\end{proof}

\begin{remark}\label{rem:slenderness-resolved}
Proposition~\ref{prop:slender-Pi03} shows that, for countable abelian groups,
slenderness is already Borel: indeed
\[
\mathrm{Slender}=\mathrm{TorsionFree}\cap\{G:G^1=0\}
\in \Pi^0_3.
\]
Thus the only remaining issue is the optimality of this upper bound.
\end{remark}

%%%%%%%%%%%%%%%%%%%%%%%%%%%%%%%%%%%%%%%%%%%%%%%%%%%%%%%%%%%%%%
\section{Extensions and open questions}
%%%%%%%%%%%%%%%%%%%%%%%%%%%%%%%%%%%%%%%%%%%%%%%%%%%%%%%%%%%%%%

\begin{enumerate}
\item \emph{Completeness.}
Theorem~\ref{thm:Cstar-Pi2-complete} proves optimality of the $G_\delta$ upper bounds
for AF-ness, real-rank bounds and topological stable-rank bounds.  The remaining
low-rank $C^*$-classes in this paper include MF-ness and approximate divisibility
($G_\delta$), and quasidiagonality, property~(SP), and nuclear dimension~$\le n$
($\Pi^0_3$).  We have only $\Pi^0_2$ lower bounds for property~(SP) and nuclear
dimension.  Are any of the remaining upper bounds sharp?
The $\Pi^1_1$-completeness results in Section~\ref{sec:Pi11-example}---separable
dual (Theorem~\ref{thm:dual-separable-Pi11}) and superatomicity
(Theorem~\ref{thm:superatomic-Pi11})---show that non-Borel phenomena span both
continuous and discrete quotient codings.  It would be interesting to identify further
such examples among less elementary properties.

\item \emph{Other invariants.}
The Borel coding of $K_0$ (Theorem~\ref{thm:K0}) and $K_1$ (Theorem~\ref{thm:K1})
extends to all higher $K$-groups via Bott periodicity (Remark~\ref{rem:Kn-KOn}).
It remains to code traces and broader Elliott-type invariants as countable Borel data.

\item \emph{Tensor-product regularity.}
Refine tensor-product kernel regularity: when does the tensor-ideal
assignment become continuous for natural choices of topology beyond
the cases treated in Theorem~\ref{thm:tensor-continuity}?

\item \emph{Slender abelian groups.}
The class of slender groups is $\Pi^0_3$ in the present
coding.  Determine whether this upper bound is optimal; in particular, is
slenderness $\Pi^0_3$-complete?

\item \emph{Dedekind finiteness for Banach algebras.}
Proposition~\ref{prop:df-banach} shows that Dedekind finiteness is closed in the
Wijsman topology.  Determine whether this closed set is non-open, or whether natural
subclasses exhibit sharper completeness phenomena.  The unbounded-witness phenomenon
of Daws--Horv\'ath~\cite{DawsHorvathCJM} remains relevant for uniform versions, but
not for the upper bound proved here.

\item \emph{Simplicity for Banach algebras and TROs.}
The present paper records Borelness of simplicity for separable $C^*$-algebras via
standard codings, but deliberately does not assert a finite Wijsman-rank bound for
Banach-algebra or TRO simplicity.  Determine the optimal complexity in these quotient
codings, both in the unital and non-unital settings.

\item \emph{Isomorphism complexity in the non-$C^*$ settings.}
For $P=C^*_{\max}(F_\infty)$, the relation
$K\sim L \iff P/K\cong P/L$
is precisely isomorphism of separable $C^*$-algebras, whose complexity is
already known by \cite{FarahTomsTornquist,SabokCompleteness}.
It would be interesting to determine the analogous complexity for the
Banach-algebra, TRO, and other quotient codings developed here.

\item \emph{Other categories.}
Apply the framework to Banach lattices with additional structure
(Banach function spaces, KB-spaces), operator systems, Jordan algebras, and other
operator-algebraic structures.
\end{enumerate}

\section*{Acknowledgements}
The author is grateful to Ilijas Farah for his interest in this work and for
helpful comments on an earlier draft, especially concerning substructure versions
of the Wijsman-closedness argument, congruence-based parametrisations, and the
comparison with existing codings of separable $C^*$-algebras. Thanks are extended to Christian Rosendal for numerous spot on remark made during the author's presentation about certain aspects of the present work. The author also thanks Jennifer Pi and Stuart White for communicating Austin Shiner's related forthcoming work. The author gratefully acknowledges support received from
NCN Sonata-Bis~13 (2023/50/E/ST1/00067).

%%%%%%%%%%%%%%%%%%%%%%%%%%%%%%%%%%%%%%%%%%%%%%%%%%%%%%%%%%%%%%
\bibliographystyle{amsplain}
\begin{thebibliography}{10}

\bibitem{Arveson1969}
W.~B.~Arveson,
\emph{Subalgebras of $C^*$-algebras},
Acta Math.\ \textbf{123} (1969), 141--224.

\bibitem{Beer}
G.~Beer,
\emph{Topologies on closed and closed convex sets},
Mathematics and its Applications, vol.~268,
Kluwer Academic Publishers, Dordrecht, 1993.

\bibitem{BalcerzakKania2023}
M.~Balcerzak and T.~Kania,
\emph{The existence of UFO implies projectively universal morphisms},
Proc.\ Amer.\ Math.\ Soc.\ \textbf{151} (2023), no.~9, 3737--3742.

\bibitem{Blackadar}
B.~Blackadar,
\emph{$K$-theory for operator algebras},
2nd ed., Mathematical Sciences Research Institute Publications, vol.~5,
Cambridge University Press, Cambridge, 1998.

\bibitem{BlackadarOA}
B.~Blackadar,
\emph{Operator algebras: Theory of $C^*$-algebras and von Neumann algebras},
Encyclopaedia of Mathematical Sciences, vol.~122, Operator Algebras and
Non-Commutative Geometry~III,
Springer-Verlag, Berlin, 2006.

\bibitem{BeckerKechris}
H.~Becker and A.~S.~Kechris,
\emph{The descriptive set theory of Polish group actions},
London Mathematical Society Lecture Note Series, vol.~232,
Cambridge University Press, Cambridge, 1996.

\bibitem{BlanchardTensorCX}
E.~Blanchard,
\emph{Tensor products of $C(X)$-algebras over $C(X)$},
Ast\'erisque \textbf{232} (1995), 81--92.


\bibitem{BrownOzawa}
N.~P.~Brown and N.~Ozawa,
\emph{$C^*$-algebras and finite-dimensional approximations},
Graduate Studies in Mathematics, vol.~88,
American Mathematical Society, Providence, RI, 2008.

\bibitem{BrownPedersen1991}
L.~G.~Brown and G.~K.~Pedersen,
\emph{$C^*$-algebras of real rank zero},
J. Funct. Anal. \textbf{99} (1991), 131--149,
\url{https://doi.org/10.1016/0022-1236(91)90056-B}.

\bibitem{CuthDolezalDouchaKurka2022}
M.~C{\'u}th, M.~Dole\v{z}al, M.~Doucha, and O.~Kurka,
\emph{Polish spaces of Banach spaces},
Forum Math.\ Sigma \textbf{10} (2022), Paper No.~e26, 28 pp.

\bibitem{CuthDolezalDouchaKurka2024}
M.~C{\'u}th, M.~Dole\v{z}al, M.~Doucha, and O.~Kurka,
\emph{Polish spaces of Banach spaces: Complexity of isometry and isomorphism classes},
J.\ Inst.\ Math.\ Jussieu \textbf{23} (2024), no.~4, 1919--1957.

\bibitem{DawsHorvath}
M.~Daws and B.~Horv{\'a}th,
\emph{A purely infinite Cuntz-like Banach $*$-algebra with no purely infinite ultrapowers},
J.\ Funct.\ Anal.\ \textbf{283} (2022), no.~1, 109488.

\bibitem{DawsHorvathCJM}
M.~Daws and B.~Horv{\'a}th,
\emph{Ring-theoretic (in)finiteness in reduced products of Banach algebras},
Canadian J.\ Math.\ \textbf{73} (2021), no.~5, 1423--1458.

\bibitem{Dixmier1977}
J.~Dixmier,
\emph{$C^*$-algebras},
North-Holland Mathematical Library, vol.~15,
North-Holland, Amsterdam--New York--Oxford, 1977.

\bibitem{DragaKania}
S.~Draga and T.~Kania,
\emph{When is multiplication in a Banach algebra open?},
Linear Algebra Appl.\ \textbf{538} (2018), 149--165.

\bibitem{Elliott1976}
G.~A.~Elliott,
\emph{On the classification of inductive limits of sequences of semisimple
finite-dimensional algebras},
J.~Algebra \textbf{38} (1976), 29--44.

\bibitem{EffrosRuanBook}
E.~G.~Effros and Z.-J.~Ruan,
\emph{Operator spaces},
London Mathematical Society Monographs, New Series, vol.~23,
Oxford University Press, Oxford, 2000.

\bibitem{ElliottFarahPaulsenRosendalTomsTornquist}
G.~A.~Elliott, I.~Farah, V.~I.~Paulsen, C.~Rosendal, A.~S.~Toms, and
A.~T{\"o}rnquist,
\emph{The isomorphism relation for separable $C^*$-algebras},
Math.~Res.~Lett.~\textbf{20} (2013), no.~6, 1071--1080.

\bibitem{FarahTomsTornquist}
I.~Farah, A.~S.~Toms, and A.~T{\"o}rnquist,
\emph{The descriptive set theory of $C^*$-algebra invariants},
Int.\ Math.\ Res.\ Not.\ IMRN \textbf{2013}, no.~22, 5196--5226.

\bibitem{FarahTomsTornquistCrelle}
I.~Farah, A.~S.~Toms, and A.~T{\"o}rnquist,
\emph{Turbulence, orbit equivalence, and the classification of nuclear $C^*$-algebras},
J.\ Reine Angew.\ Math.\ \textbf{688} (2014), 101--146.

\bibitem{FarahHartLupiniRobertTikuisisVignatiWinter}
I.~Farah, B.~Hart, M.~Lupini, L.~Robert, A.~Tikuisis, A.~Vignati, and
W.~Winter,
\emph{Model theory of $C^*$-algebras},
Mem.~Amer.~Math.~Soc.~\textbf{271} (2021), no.~1324.

\bibitem{FarahMagidor2018}
I.~Farah and M.~Magidor,
\emph{Omitting types in logic of metric structures},
J. Math. Logic \textbf{18} (2018), no.~2, 1850006.

\bibitem{FriedmanStanley}
H.~Friedman and L.~Stanley,
\emph{A Borel reducibility theory for classes of countable structures},
J.\ Symbolic Logic \textbf{54} (1989), no.~3, 894--914.

\bibitem{FreeBLl1}
A.~Avil{\'e}s, J.~Rodr{\'i}guez, and P.~Tradacete,
\emph{The free Banach lattice generated by a Banach space},
J.\ Funct.\ Anal.\ \textbf{274} (2018), no.~10, 2955--2977.

\bibitem{FuchsII}
L.~Fuchs,
\emph{Infinite abelian groups. Vol.~II},
Pure and Applied Mathematics, vol.~36-II,
Academic Press, New York--London, 1973.

\bibitem{GaoDST}
S.~Gao,
\emph{Invariant descriptive set theory},
Pure and Applied Mathematics (Boca Raton), vol.~293,
CRC Press, Boca Raton, FL, 2009.

\bibitem{Gromov1999}
M.~Gromov,
Endomorphisms of symbolic algebraic varieties,
J.\ Eur.\ Math.\ Soc.\ \textbf{1} (1999), no.~2, 109--197.

\bibitem{GodefroySaintRaymond2018}
G.~Godefroy and J.~Saint-Raymond,
\emph{Admissible topologies on the space of closed subspaces of a separable Banach space},
J.\ Funct.\ Anal.\ \textbf{275} (2018), no.~9, 2379--2407.


\bibitem{Hjorth}
G.~Hjorth,
\emph{Classification and orbit equivalence relations},
Mathematical Surveys and Monographs, vol.~75,
American Mathematical Society, Providence, RI, 2000.

\bibitem{HermanVaserstein1984}
R.~H.~Herman and L.~N.~Vaserstein,
\emph{The stable range of $C^*$-algebras},
Invent.\ Math.\ \textbf{77} (1984), no.~3, 553--555.

\bibitem{KaniaNoUniversalBanachAlg}
T.~Kania,
\emph{There is no universal separable Banach algebra},
preprint, arXiv:2511.04314, 2025.

\bibitem{KaniaMaslany}
T.~Kania and N.~Ma{\'s}lany,
\emph{Differential embeddings into algebras of topological stable rank~$1$},
Proc.\ Roy.\ Soc.\ Edinburgh Sect.~A, published online 20 November 2024,
1--25, doi:10.1017/prm.2024.108.

\bibitem{Kirchberg1993}
E.~Kirchberg,
\emph{On non-semisplit extensions, tensor products and exactness of group $C^*$-algebras},
Invent.\ Math.\ \textbf{112} (1993), no.~3, 449--489.

\bibitem{KirchbergWassermann1995}
E.~Kirchberg and S.~Wassermann,
\emph{Operations on continuous bundles of $C^*$-algebras},
Math.\ Ann.\ \textbf{303} (1995), no.~4, 677--698.

\bibitem{Kechris}
A.~S.~Kechris,
\emph{Classical descriptive set theory},
Graduate Texts in Mathematics, vol.~156,
Springer-Verlag, Berlin, 1995.

\bibitem{KechrisAsianCstar}
A.~S.~Kechris,
\emph{The descriptive classification of some classes of $C^*$-algebras},
in \emph{Proceedings of the Sixth Asian Logic Conference (Beijing, 1996)},
World Scientific, River Edge, NJ, 1998, 121--149.

\bibitem{KPTT}
A.~Kavruk, V.~I.~Paulsen, I.~G.~Todorov, and M.~Tomforde,
\emph{Quotients, exactness, and nuclearity in the operator system category},
Adv.\ Math.\ \textbf{235} (2013), 321--360.

\bibitem{Loring}
T.~A.~Loring,
\emph{Lifting solutions to perturbing problems in $C^*$-algebras},
Fields Institute Monographs, vol.~8,
American Mathematical Society, Providence, RI, 1997.

\bibitem{LyndonSchupp}
R.~C.~Lyndon and P.~E.~Schupp,
\emph{Combinatorial group theory},
Ergebnisse der Mathematik und ihrer Grenzgebiete, vol.~89,
Springer-Verlag, Berlin--Heidelberg, 1977.

\bibitem{NiwinskiBL}
M.~Niwi{\'n}ski,
\emph{Polish spaces of separable Banach lattices},
in preparation.

\bibitem{Novikov1955}
P.~S.~Novikov,
\emph{Ob algoritmicheskoj nerazreshimosti problemy tozhdestva slov v teorii grupp}
[On the algorithmic unsolvability of the word problem in group theory],
Trudy Matematicheskogo Instituta imeni V.~A.~Steklova \textbf{44} (1955), 3--143
(Russian).

\bibitem{Nunke1962}
R.~J.~Nunke,
\emph{Slender groups},
Acta Sci. Math. (Szeged) \textbf{23} (1962), 67--73.

\bibitem{PisierOS}
G.~Pisier,
\emph{Introduction to operator space theory},
London Mathematical Society Lecture Note Series, vol.~294,
Cambridge University Press, Cambridge, 2003.

\bibitem{PiSzachniewiczTartarotti2026}
J.~Pi, M.~Szachniewicz, and M.~Tartarotti,
\emph{A game-theoretic unital classification theorem for $C^*$-algebras},
preprint, arXiv:2601.01735, 2026.

\bibitem{RordamLarsenLaustsen}
M.~R{\o}rdam, F.~Larsen, and N.~J.~Laustsen,
\emph{An introduction to $K$-theory for $C^*$-algebras},
London Mathematical Society Student Texts, vol.~49,
Cambridge University Press, Cambridge, 2000.

\bibitem{Rieffel1983}
M.~A.~Rieffel,
\emph{Dimension and stable rank in the $K$-theory of $C^*$-algebras},
Proc.\ London Math.\ Soc.\ (3) \textbf{46} (1983), no.~2, 301--333.

\bibitem{SabokCompleteness}
M.~Sabok,
\emph{Completeness of the isomorphism problem for separable $C^*$-algebras},
Invent.\ Math.\ \textbf{204} (2016), no.~3, 833--868.

\bibitem{ShinerForthcoming}
A.~Shiner,
work in preparation on quotient encodings of separable $C^*$-algebras and Borel
functoriality in $KK$-theory.

\bibitem{TomsWinter}
A.~S.~Toms and W.~Winter,
\emph{Strongly self-absorbing $C^*$-algebras},
Trans.\ Amer.\ Math.\ Soc.\ \textbf{359} (2007), 3999--4029.

\bibitem{WeggeOlsen}
N.~E.~Wegge-Olsen,
\emph{$K$-theory and $C^*$-algebras: a friendly approach},
Oxford Science Publications,
Oxford University Press, New York, 1993.

\bibitem{WinterZacharias}
W.~Winter and J.~Zacharias,
\emph{The nuclear dimension of $C^*$-algebras},
Adv.\ Math.\ \textbf{224} (2010), no.~2, 461--498.

\bibitem{Zielinski2016}
J.~Zielinski,
\emph{The complexity of the homeomorphism relation between compact metric spaces},
Adv. Math. \textbf{291} (2016), 635--645.


\end{thebibliography}

\end{document}
